% =====================================================================
%  Harald Høffding, DEN NYERE FILOSOFIS HISTORIE. En Fremstilling af
%  Filosofiens Historie fra Renaissansens Slutning til vore Dage.
%  Anden gennemsete Udgave.  Første Bind: København, Det Nordiske Forlag,
%  Ernst Bojesen, 1903.  Andet Bind: København, Gyldendalske Boghandel,
%  Nordisk Forlag, 1904.  Aarhus Stiftsbogtrykkeri.
%  A SELECTION: the INDLEDNING (vol. I pp. XI-XV), the TIENDE BOG
%  "Filosofien i Tyskland 1850-1880" (vol. II pp. 493-565), and its notes
%  113-130 from the "Noter til andet Bind" (vol. II pp. 596-600) -- the parts
%  translated in translation.tex.
%  DIPLOMATIC TRANSCRIPTION -- GENERATED; DO NOT EDIT BY HAND.
%  Made by ebook-method/nyere-filosofis-historie/build.py --emit from this
%  template, the decisions (decisions.tsv) and the patches (patches.tsv) there.
%
%  WITNESSES.  (1) The SAGA e-book (Den_nyere_filosofis_historie.epub, ISBN
%  9788726363043), an OCR of the 2nd edition: base text, paragraphs, italics,
%  small caps, notes.  (2) The Google scans of the two volumes at the Internet
%  Archive (dennyerefilosof01hfgoog, dennyerefilosof00hfgoog), read by tesseract
%  (ebook-method/collate/ocr_layer.py): every letter, the page turns, the
%  letterspacing.  Every disagreement no rule settles was read at the image
%  (decisions.tsv).  Print governs; misprints are kept and marked PRINTED AS IS.
%  The scans lose the tails of commas and semicolons; where the scan shows only a
%  point or a colon for the e-book's comma or semicolon, the e-book's mark stands.
%
%  CONVENTIONS.  "% --- p. N ---" + \opage{N} at the first word of printed page N
%  (a word broken over the turn is split with %); the Introduction's pages are
%  roman.  The note anchors as printed, \textsuperscript{113}); the notes follow
%  the Tenth Book, numbered as printed (113-130, with an added 117 b), each
%  beginning "S. N." with the page it belongs to.  Small capitals \textsc;
%  italics \textit; letterspacing \emph.  Smaller type (the print sets some
%  passages of the body in it) is not marked.  Running heads, folios and
%  signature marks are not transcribed.
% =====================================================================
\documentclass[11pt]{article}
\usepackage[utf8]{inputenc}
\usepackage[T1]{fontenc}
\usepackage{textalpha}
\usepackage{libertinus}
\usepackage{libertinust1math}
\usepackage[danish]{babel}
\usepackage[protrusion=true,expansion=true]{microtype}
\usepackage{setspace}
\usepackage[margin=1.3in]{geometry}
\usepackage{fancyhdr}
\usepackage{hyperref}
\hypersetup{pdftitle={Den nyere Filosofis Historie (udvalg)}, pdfauthor={Harald Høffding}, hidelinks}
\pagestyle{fancy}
\fancyhf{}
\fancyhead[LE,RO]{\thepage}
\fancyhead[C]{\textit{Harald Høffding: Den nyere Filosofis Historie (1903--04), udvalg}}
\setstretch{1.3}
\usepackage{marginnote}
\usepackage{graphicx}
\DeclareUnicodeCharacter{0254}{\reflectbox{c}}   % ɔ (id est)
\newcommand{\opage}[1]{\marginnote{\footnotesize\textit{[#1]}}}

\begin{document}

% --- front matter: title of vol. I (PDF 14) ---
\begin{center}
{\LARGE DEN NYERE}\\[0.5em]
{\LARGE FILOSOFIS HISTORIE}\\[1.5em]
EN FREMSTILLING AF\\
FILOSOFIENS HISTORIE FRA RENAISSANSENS SLUTNING\\
TIL VORE DAGE\\[1.5em]
AF\\[1em]
{\large HARALD HØFFDING}\\[2em]
ANDEN GENNEMSETE UDGAVE\\[1.5em]
FØRSTE BIND\\
FRA RENAISSANSENS SLUTNING TIL ROUSSEAU\\[6em]
KØBENHAVN\\
DET NORDISKE FORLAG\\
ERNST BOJESEN\\
1903
\end{center}
% (vol. II, PDF 15: ... ANDET BIND / FRA LESSING OG KANT TIL VORE DAGE / KØBENHAVN /
%  GYLDENDALSKE BOGHANDEL / NORDISK FORLAG / 1904)
\clearpage

% ======== FØRSTE BIND (vol. I, 1903), INDLEDNING, pp. XI-XV ========
%
% --- p. XI ---
\opage{XI}

\begin{center}\large INDLEDNING.\end{center}

Det vil være naturligt, at der, naar en ny Fremstilling af en vigtig Periode i
Filosofiens Historie udkommer, spørges, hvilken Opfattelse af Filosofien
Forfatteren lægger til Grund, og hvilken Betydning og Værdi han tillægger
Filosofiens Historie. Mit Svar herpaa er, at de Studier, hvis foreløbige
Afslutning betegnes ved denne Bogs Udgivelse, netop havde til Hensigt at
orientere med Hensyn til, hvad Filosofi egentligt er. Som vi lære et Menneske
at kende gennem hans Biografi, saaledes maa vi kunne lære en Videnskab at kende
gennem dens Historie. Og det vil være saameget des naturligere at prøve denne
Vej, som Erfaring jo stadigt viser, at der paa Filosofiens Omraade gør sig
modstridende Synsmaader gældende, saa at der ikke kan henvises til nogen
Fremstilling af Filosofien, der skulde kunne siges at give et udtømmende Begreb
om den. Det vil da være naturligt her --- ligesom paa Religionens Omraade ---
at anvende den sammenlignende Metode. Filosofiens Historie behandler de Forsøg,
som enkelte Tænkere have gjort paa at drøfte Erkendelsens og Livets sidste
Problemer. Det vigtigste Resultat vil nu være at faa Rede paa, hvilke Problemer
Filosofien beskæftiger sig med, hvorledes disse Problemer stilles til de
forskellige Tider, og hvoraf de forskellige Problemstillinger og Forsøg paa
Problemløsning i det Enkelte bestemmes. Hvis et saadant Resultat
tilnærmelsesvis kan naas, vil Studiet af Filosofiens Historie være af ikke
ringe Betydning for den fortsatte Filosoferen.

For mit Vedkommende er jeg ved Studiet af den nyere Filosofis Historie bleven
bestyrket i den Antagelse, jeg ad anden %
% --- p. XII ---
\opage{XII}
Vej var kommen til, at den filosofiske Forsken drejer sig om fire
Hovedproblemer. Disse fire Problemer skulle nu kort karakteriseres. En mere
indgaaende Drøftelse og Sammenligning af dem har jeg givet i Afhandlingen
\textit{Filosofiske Problemer} (Universitetsprogram 1902).

1. \emph{Erkendelsesproblemet} (det \emph{logiske} Problem). Hvor forskellige
de enkelte Videnskaber end ere i Henseende til Genstand og Metode, saa arbejde
de dog alle med den menneskelige Tænkning. Hver Gang de danne et Begreb, udtale
en Dom eller drage en Slutning, forudsætte de Tænkningens almindelige Former og
Principer. Der viser sig da her Muligheden af en særegen Disciplin, der
undersøger de Former, i hvilke Tænkningen bevæger sig, og de Principer, den
efter sin Natur maa lægge til Grund, hvilke Genstande den saa end angaar. Denne
Disciplin, \textit{den formelle Logik} behandler dog kun en Del af
Erkendelsesproblemet. Hine Former og Principer føre ikke udover Tænkningen
selv, men gøre det kun muligt for denne at være i Overensstemmelse med sig
selv, at være konsekvent. Hver Gang de anvendes paa givne Fænomener, som vor
Tænkning ikke selv har konstrueret, men maa tage, som de komme, opstaar det
Spørgsmaal, med hvilken Ret denne Anvendelse sker, --- med hvilken Ret vi
antage, at ikke blot vor Tænkning, men selve Tilværelsen, der udtrykker sig i
de givne Fænomener, er konsekvent, stemmer overens med sig selv. Herved viser
sig Muligheden af en Disciplin, der undersøger Betingelserne for en Erkendelse
af Tilværelsen og Grænserne for en saadan Erkendelse. Denne Disciplin er
\textit{Erkendelsesteorien.}

2. \emph{Tilværelsesproblemet} (det \emph{kosmologiske} Problem) opstaar, naar
der spørges om, hvorledes den Tilværelse, af hvilken vi ere Led, maa antages at
være beskaffen, naar vi drage de sidste Konsekvenser af Alt, hvad vi vide,
eller ved de sandsynligste Hypotesers Hjælp kunne formode. Vi kalde dette
Problem det kosmologiske (af Kosmos, Verden, og Logos, Lære), fordi det fører
til Drøftelse af de Muligheder, der fremstille sig for Tanken, naar den søger
at sammenarbejde det i Erfaringen Givne til en almindelig Verdensanskuelse,
eller naar den ad Spekulationens dristige Vej forsøger at konstruere en saadan.
De for%
% --- p. XIII ---
\opage{XIII}
skellige filosofiske Systemer ere Forsøg i denne Retning, Forsøg, hvis Værdi
beror paa, hvor betydningsfulde og omfattende Erfaringer, de lægge til Grund,
og hvor stor Konsekvens og Kombinationsevne, der lægger sig for Dagen i deres
Konstruktioner.

3. \emph{Vurderingsproblemet} (det \emph{etisk}-\emph{religiøse} Problem)
opstaar ved, at vi ikke blot forholde os anskuende og forstaaende til
Tilværelsen, men ogsaa faa vor Følelse sat i Bevægelse ved den, saa at vi
udtale Domme, der tillægge eller frakende den Værdi. Af særlig Betydning ere de
Domme, vi udtale om menneskelige Handlinger, vore egne og Andres. Al saadan
Bedømmelse hviler --- ligesom al Erkendelse og Forstaaelse --- paa visse
Forudsætninger, som det gælder at paavise og bestemme. Denne Opgave hører under
det etiske Problem. Naar Vurderingen ikke blot angaar menneskelige Handlinger
og Indretninger, men Tilværelsen, Livet i det Hele, opstaar det religiøse
Problem, der især fører til at drøfte Forholdet mellem de etiske Idealer og den
virkelige Tilværelse, saa at det bliver en Kombination af det kosmologiske og
det etiske Problem.

4. \emph{Bevidsthedsproblemet} (det \emph{psykologiske} Problem). For alle de
tre antydede Problemers Vedkommende viser det sig, at deres Behandling
forudsætter et erfaringsmæssigt Kendskab til det menneskelige Bevidsthedsliv.
Psykologien fremstiller den menneskelige Erkendelses \textit{faktiske}
Udvikling, som maa kendes, før man kan drøfte Betingelserne for Erkendelsens
\textit{Gyldighed.} Og da Forholdet mellem det Aandelige og det Materielle er
et Hovedpunkt indenfor Tilværelsesproblemet, forudsættes Psykologien ogsaa i
Kosmologien. Hvad endeligt Vurderingsproblemet angaar, drøfter Psykologien dels
Beskaffenheden af selve de Følelser, der føre til at anstille Vurderinger, dels
de Muligheder, det virkelige Bevidsthedsliv frembyder for en videre Udvikling i
den Retning, Vurderingen fordrer. Som Følge af dette nøje Forhold til de
filosofiske Problemer, er Psykologien selv at regne for en Del af Filosofien,
og man vil derfor ogsaa omvendt ganske naturligt i Psykologien komme til at
berøre filosofiske Problemer. Og selv om man vil erklære alle de tre
foregaaende Problemer for uløselige eller for opstaaede ved Misforstaaelse,
vilde dog det psykologiske Problem, Spørgsmaalet %
% --- p. XIV ---
\opage{XIV}
om Bevidsthedslivets Natur og Love, staa tilbage som Filosofiens sidste Borg.

--- Spørges der om de Faktorer, der efter Sagens Natur ville faa Indflydelse
paa disse Problemers Behandling og Besvarelse, maa først nævnes den
Filosoferendes Personlighed. De nævnte Problemer have det fælles, at de ligge
ved Grænsen af vor Erkendelse, hvor de exakte Metoder svigte; det kan da ikke
være Andet, end at Forskerens Personlighed maa komme til at bestemme hans
Tankegang, uden at han behøver at være sig det bevidst. Den »personlige
Ligning« maa være af større Betydning i Filosofien end paa andre videnskabelige
Omraader. Den historiske, sammenlignende Metode bliver derfor her af særlig
Vigtighed; ved dens Hjælp kan det personlige Element lettest paavises. Dette
personlige Element er ikke altid Noget, der skal skydes til Side; dets
Tilstedeværelse er ofte en Betingelse for, at et Problem kan blive til. Der er
Tanker, som kun kunne blive til paa en bestemt psykologisk Jordbund. --- Det
vil for det Andet komme an paa, hvilke Iagttagelser der lægges til Grund. Her
faar især Naturvidenskabens Udvikling stor Betydning for den nyere Filosofi.
Som det vil vise sig, bestemmes den nyere Filosofis afgørende Problemer ved
selve det Faktum, at den moderne Naturvidenskab blev til. Men dertil komme ---
især for Vurderingsproblemets Vedkommende --- ogsaa historiske Forhold,
aandelige Bevægelser paa andre Omraader. --- Endeligt bestemmes Problemstilling
og Problemløsning af den Konsekvens, hvormed de grebne Udgangspunkter
fastholdes og gennemføres.

Baade Karakteristiken og Kritiken af ethvert filosofisk Forsøg vil stedse føre
tilbage til disse tre Punkter. Blandt dem lægger jeg i den følgende
Fremstilling mest Vægt paa de to første. En Inkonsekvens er hos en stor Tænker
ofte netop en naturlig Følge af, at \textit{flere} Tankerækker have aabnet sig
for hans Geni, uden at han har kunnet forfølge dem saa langt, at han har
opdaget deres indbyrdes Modstrid. Og det kan være af den største Betydning, at
disse forskellige Tankerækker drages frem. Allerbedst er det naturligvis, naar
Dybsindighed og Konsekvens forenes.

Et Forsøg paa Problemopstilling eller Problemløsning kan %
% --- p. XV ---
\opage{XV}
have en dobbelt Interesse. For det Første kan det betragtes som Symptom, som
historisk Udtryk for aandelige Strømninger. Fra denne Side er Filosofiens
Historie en Del af den almindelige Kulturhistorie. For det Andet kan det
undersøges med Hensyn til den virkelige, definitive Klarhed, der er naat ved
det. Denne Betragtning er mere specielt filosofisk og teknisk. De to Sider, fra
hvilke de filosofiske Fænomener kunne betragtes, ville naturligvis for de
forskellige Fænomeners Vedkommende staa i højst forskelligt Forhold til
hinanden. Undertiden vil den kulturhistoriske Interesse endogsaa kunne staa i
omvendt Forhold til den rent filosofiske.

Naar jeg i Korthed skal angive, hvorved min Fremstilling af den nyere Filosofis
Historie er forskellig fra sin Forgænger i vor Literatur, Brøchners 1873---74
udkomne Værk, saa vil jeg først nævne den særlige Vægt, jeg lægger paa den
personlige Faktor og paa Forholdet til Erfaringsvidenskaben, samt paa den
kulturhistoriske Betydning af de filosofiske Fænomener. Dernæst vil jeg henvise
til min Bestræbelse for at give Problemstillingen Forrangen for
Problemløsningen. Løsningerne kunne dø, og Problemerne leve dog; ellers havde
Filosofien ikke havt saa langt et Liv, som den har havt. Endeligt følger det af
sig selv, at et fornyet Kildestudium saavel som Benyttelse af de senere Aars
frodige Literatur have stillet mange Fænomener i et nyt Lys. Men trods alt
Dette har det dog været mig en kær Tanke ved dette Værk at kunne fortsætte
Brøchners Arbejde paa Filosofiens Histories Omraade, og jeg vil kun haabe, at
min Bog ikke maa være uværdig til at helliges Mindet om denne ædle Forsker,
hvis Beundring for Tankens Heroer paa én Gang var saa dyb og saa forstaaende.

\begin{center}\rule{2cm}{0.4pt}\end{center}

% ======== ANDET BIND (vol. II, 1904), TIENDE BOG, pp. 493-565 ========
\clearpage
%
% --- p. 493 ---
\opage{493}

\begin{center}\large TIENDE BOG.\end{center}

\begin{center}FILOSOFIEN I TYSKLAND 1850--1880.\end{center}

Af de to store Retninger i det 19. Aarhundredes Filosofi var det Positivismen,
der bedst havde vedligeholdt Forbindelsen dels med det 18. Aarhundredes
Tankegang, dels med Erfaringsvidenskaben. Romantikens Filosofi var derimod et
bestemt og bevidst Forsøg paa Reaktion mod begge disse Retninger; man vilde jo
endog gøre det helt om, der var blevet grundlagt ved Naturvidenskabens Opstaaen
i det 17. Aarhundrede! I Tyskland, Arnestedet for Romantiken og dens Filosofi,
var ved Aarhundredets Midte, denne Retning herskende. Den kritiske Understrøms
Tilhængere vare de eneste, der fastholdt Kontinuiteten med de andre
Videnskaber, naar man ser bort fra Schopenhauer og Feuerbach, der paa dette
Tidspunkt endnu stode som ensomme Tænkere, for hvem Anerkendelsens Tid ikke var
kommen.

Ligesaa lidt som Positivismen var at forklare som Reaktion mod Romantiken,
ligesaa lidt er den filosofiske Bevægelse, der finder Sted i Tyskland efter
Aarhundredets Midte, at forklare som en Fortsættelse af den fransk-engelske
Positivisme. Den har sine Forudsætninger i Tyskland selv. Den faar især sine
Problemer ved det fornyede Opsving af Naturvidenskaben henimod Aarhundredets
Midte. Ikke blot blive nu naturvidenskabelige Studier og Resultater Genstand
for mere almindelig Interesse end i Aarhundredets første Del, hvor Poesi,
Religion og spekulativ Filosofi havde lagt Beslag paa Sindene. Men Natur%
% --- p. 494 ---
\opage{494}
videnskaben vendte nu selv med klarere Bevidsthed tilbage til de store
Principer, dens Grundlæggere i sin Tid havde opstillet. Fordringen om en
Forstaaelse af den ydre Natur som en Række af paaviselige Aarsager og
Virkninger, med andre Ord, om en mekanisk Naturforklaring, blev nu hævdet igen,
ligesom i det 17. Aarhundrede. Og samtidigt hermed opdagedes og paavistes en ny
stor, omfattende Naturlov: den Lov, at der i den fysiske Natur ingen Energi
opstaar eller forgaar, men at det, der sker ved enhver tilsyneladende Opstaaen
eller Forgaaen, er en Omsætning af Energien til en anden Form, og det saaledes,
at de forskellige Former for Energi staa i et bestemt kvantitativt Forhold til
hverandre. Denne Lov --- der tilligemed Darwins Lære om Arternes Oprindelse
gennem Kampen for Tilværelsen er det betydningsfuldeste Resultat af de
naturvidenskabelige Undersøgelser i det 19. Aarhundrede --- maatte nødvendigvis
sætte den filosofiske Tænkning i Bevægelse, ligesom i sin Tid den kopernikanske
Teori og Galileis Grundlæggelse af Mekaniken havde ført Filosofien ind paa nye
Baner. Det maatte særligt for den tyske Filosofi blive et Hovedspørgsmaal:
hvorvidt kunne de Tanker, Romantikens Filosofi udviklede, fastholdes, naar man
gaar ind paa den nye naturvidenskabelige Betragtningsmaade? Dette Spørgsmaal
besvaredes paa tre forskellige Maader. Den moderne Materialisme forkastede hine
Tanker som rene Illusioner og proklamerede den Lære, at kun det Materielle
existerer, som en nødvendig Konsekvens af Naturvidenskaben. Lotze og Fechner
forsøgte derimod at paavise den spekulative Religionsfilosofis Grundtanke som
sidste og afsluttende Forudsætning for det Billede af Tilværelsen, der kan
opbygges ved den naturvidenskabelige Metodes Hjælp. Endeligt nærmede Albert
Lange og Eugen Dühring sig den kritiske Filosofi og Positivismen ved at
fremdrage Erkendelsesproblemets Betydning og ved at hævde den praktiske
Idealismes Uafhængighed af Erfaringsvidenskaben, samtidigt med, at Erfaringens
Ret til at bestemme Verdensbilledets faktiske Indhold indskærpedes. %
% --- p. 495 ---
\opage{495}

\begin{center}1. ROBERT MAYER OG PRINCIPET OM ENERGIENS BESTAAEN.\end{center}

Ikke blot i Filosofien, ogsaa i andre Videnskaber bevæger Erkendelsen sig
fremad paa rytmisk Maade, gennem Aktioner og Reaktioner. Ved Slutningen af det
18. Aarhundrede havde Naturvidenskaben en stor og frugtbar Periode, idet nogle
af de vigtigste Sandheder paa Kemiens og Fysiologiens Omraader opdagedes.
LAVOISIER indførte den kvantitative Metode i Kemien, og ved dens Hjælp
godtgjordes Sandheden af den gamle Tanke, at intet Stof opstaar eller forgaar,
men at den samme Stofmasse bestaar trods alle Forandringer, under forskellige
Former. Kemien blev grundlagt som exakt Videnskab. PRIESTLEY, INGENHOUSS,
SENEBIER og SAUSSURE fandt Hovedlovene for Planternes og Dyrenes Stofskifte og
grundlagde derved den store Lære om Stoffets Kredsløb i Naturen, hvorved den
organiske og den uorganiske Verden fremtræde i deres nøje indbyrdes
Forbindelse. Man opdagede, at det er under Sollysets Indflydelse, at det
organiske Stof dannes i de grønne Planteceller ved Optagelse og Opløsning af
Luftens Kulsyre. Det i Plantecellerne oplagrede Kulstof tjener saa Dyrene til
Ernæring; ved de dyriske Funktioner sker der en Forbrænding, og den derved
dannede Kulsyre gaar ved Udaandingen ud i Luften, hvorfra saa det samme
Kredsløb kan begynde paany. En stor kosmisk Sammenhæng var da her bleven
paavist\textsuperscript{113}). Det af Kopernikus, Bruno, Kepler og Galilei
udkastede Verdensbillede, som Newton havde beriget ved Paavisningen af den
almindelige Tiltrækning som sammenholdende og ordnende Baand, fik herved en ny
og betydningsfuld Berigelse.

Men Tiden var for optagen af Revolution og Verdenskrig, af Romantik og
Spekulation, snart ogsaa af Ortodoxi og Mystik til, at disse store Ideer kunde
faa den Indflydelse paa den almindelige Opfattelse, som de fortjente. Selv
indenfor Naturvidenskaben hæmmedes deres Sejr af andre Tendenser. Det var
indenfor Videnskaben om det organiske Liv, at der især blev gjort Modstand mod
de nye Synspunkter, fordi man endnu beherskedes af Reaktionen mod den
kartesianske Lære om Organ%
% --- p. 496 ---
\opage{496}
ismen som en Maskine. Man udledte de for Organismen ejendommelige Ytringer af
en særegen Livskraft, der skulde være helt forskellig fra de andre
Naturkræfter. Denne Forklaringsmaade, den saakaldte Vitalisme, maatte føre til
at afvise Forsøgene paa at drage de organiske Fænomener ind i Stoffets
universelle Kredsløb\textsuperscript{114}). Desuden var den herskende Interesse
hos Botanikere og Zoologer den, at beskrive Formerne og bestemme deres
systematiske Plads, ikke at finde den Udviklingsproces, gennem hvilken Formerne
blive til, og de Aarsager, der bestemme denne Udviklingsproces. Den romantiske
Naturfilosofi begunstigede denne æstetiske og formale Opfattelse --- og var jo
selv fremgaaet af den. Naar »Naturfilosofien« ofte faar hele Skylden for, at
Synspunkterne fra det 18de Aarhundredes Slutning ikke tidligere sejrede i
Botaniken og i Zoologien, maa det ikke glemmes, at den nok saa meget var
Virkning af Naturvidenskabens Tilstand som Aarsag til denne.

Omtrent fra 1840 kom der en Forandring i denne Henseende. Man begyndte at
opfatte Livet som Andet end et Spil af Former eller en Aabenbaring af Ideer.
\textsc{Dumas} og \textsc{Liebig} indskærpede de kemiske Processers Betydning
for Plante- og Dyrelivet. Og man begyndte ogsaa fra medicinsk Side at fordre en
streng mekanisk Forklaring af de Processer, der gaa for sig i Organismen. I
nogle af sine første Skrifter virkede \textsc{Hermann Lotze} i denne Retning,
især i \textit{Allgemeine Pathologie und Therapie als mechanische
Naturwissenschaften} (1842) og i Afhandlingen \textit{Leben, Lebenskraft}
(1843, i R. Wagners Handwörterbuch der Physiologie). Den Opgave blev stillet at
gøre Fysiologien til en rent mekanisk Videnskab. Mest epokegørende blev dog den
Grundsætning, Lægen og Fysikeren \textsc{Robert Mayer} (f. 1814, d. 1878)
opstillede i sit Skrift \textit{Die organische Bewegung in ihrem Zusammenhange
mit dem Stoffwechsel} (1845): at der under Livsprocessen kun foregaar en
Forvandling af Kraft saavel som af Stof, aldrig en Skabelse af dem. Mayer viser
her tilbage til den store Lov, han faa Aar i Forvejen havde opdaget, idet han
havde stillet Loven om Kraftens Bestaaen ved Siden af den fra Lavoisier
stammende Lov om Stoffets Bestaaen. %
% --- p. 497 ---
\opage{497}

Robert Mayers Forskerliv er baaren af en eneste Tanke, som han tidligt kom ind
paa; men det var hans Livs Tragedie, at han ikke formaaede at gennemføre den i
det Enkelte ad experimental Vej. Som Barn havde han forsøgt at konstruere et
perpetum %
% PRINTED AS IS: „perpetum“ (p. 497)
mobile, og det gjorde et stærkt Indtryk paa ham, at det ikke lykkedes. Fra den
Tid af grublede han stadigt over Forholdet mellem Aarsag og Virkning, først og
fremmest indenfor Fysiologien, der laa ham som Mediciner nærmest, men derefter
ogsaa indenfor Kemien og Fysiken. Paa en Rejse til Ostindien som Skibslæge fik
han den nye Ide om Kraftens Uforgængelighed i Naturen, under Indflydelse dels
af Undersøgelser af den organiske Varmes Oprindelse, dels af den Kendsgerning,
at der ved Bølgebevægelsen i Havet opstaar Varme. Hans Udgangspunkt er altsaa
bestemt Iagttagelse. Men denne Iagttagelse førte ham til at spørge, om der ikke
i Fysiken skulde gælde en lignende Lov som i Kemien, saa at ikke blot Stoffet,
men ogsaa Kraften var uforgængelig, og enhver Forandring en blot Omsætning. Og
han kom da tilbage til sine tidligere Grublerier om Aarsag og Virkning. Det var
for ham en selvindlysende Sætning, at der ikke kunde være Mere i Virkningen end
i Aarsagen: causa æqvat effectum! Det følger af »den logiske Grunds Lov«! Mayer
skelner altsaa ikke, ligesaa lidt som de dogmatiske Filosofer, mellem Grund og
Aarsag; derved kommer han, ligesom William Hamilton og Herbert Spencer, til at
betragte det som en selvfølgelig Ting, at Aarsagen ikke kan blive til Intet,
naar Virkningen opstaar, men at Virkningen maa yde Ækvivalens for det, der
syntes at forsvinde med Aarsagen. Hidtil havde man slaaet sig til Ro med, at en
Bevægelse ophørte, naar den mødte en tilstrækkeligt stærk Modstand, og med, at
der ved Gnidningsmodstanden opstod Varme. Men --- spurgte Mayer --- bliver
Bevægelsen da til Intet, og opstaar der saa Varme af Intet? Hvis det er saa,
saa overskæres »Videnskabens røde Traad«, ligesaa vel, som hvis vi i Kemien
vilde antage, at Ilt og Brint blev til Intet, naar de forenedes, og at saa
Vandet blev til af Intet. Ligesom vi antage, at Ilt og Brint forvandles til
Vand, saaledes maa vi ogsaa antage, at Bevægelsen ikke bliver til Intet, men
forvandles til Varme\textsuperscript{115}). Og ved Udtrykket \textit{for%
% --- p. 498 ---
\opage{498}
vandles} vilde Mayer kun betegne, at der bestaar et konstant kvantitativt
Forhold mellem Aarsagen, som forsvinder, og Virkningen, som kommer i dens Sted.
Hvis den fremførte Tankegang er rigtig, maa Erfaringen godtgøre et saadant
Ækvivalensforhold mellem de forskellige Kraftformer, der afløse hverandre. Paa
Grundlag af foreliggende Forsøg søgte Mayer allerede i sin første Afhandling
(\textit{Bemerkungen über die Kräfte der unbelebten Natur.} Trykt i Liebigs
»Annalen der Chemie und Pharmacie«. 1842) at bestemme det kvantitative Forhold
mellem Varme og Bevægelse. Det var denne Ide om et uforanderligt
Størrelsesforhold, som var Hovedsagen for Mayer, og som i den følgende Tid
viste sig saa frugtbar for Naturvidenskaben. Den var en betydningsfuld
Udvidelse af den allerede af Huyghens og Leibniz hævdede abstrakte Lov om
Kraftens Bestaaen til at gælde for Forholdet mellem forskellige Naturkræfter
indbyrdes. Mayer udledte af sin Tankegang, at der egentligt kun existerer en
eneste Kraft, der fremtræder under forskellige Former, som indbyrdes staa i
bestemt kvantitativt Forhold til hverandre.

Skønt Mayers Forsøg paa at give en filosofisk Udledning af Loven om Energiens
Bestaaen ikke slaar til, saa har hans Tankegang dog erkendelsesteoretisk
Interesse. Ti det var i sig selv fuldt berettiget at opkaste det Spørgsmaal, om
der ikke mellem den Begivenhed, vi kalde Aarsagen, og den Begivenhed, vi kalde
Virkningen, skulde finde et lignende Forhold Sted som mellem logisk Grund og
logisk Følge. Og til at opkaste dette Spørgsmaal fører Mayers Tankegang
naturligt. Saa maa Erfaringen afgøre, hvad Svaret skal blive. --- Det er
interessant, at den danske Fysiker \textsc{Colding}, der Aaret efter, at Mayers
første Afhandling var bleven trykt, ad sin egen Vej opstillede Sætningen om
Energiens Bestaaen og bekræftede den ved Forsøg, gik ud fra, at den egentligt
var en Fornuftsætning, og at ligeledes \textsc{Helmholtz} i sin Afhandling
\textit{Ueber die Erhaltung der Kraft} (1847) gik ud fra et
erkendelsesteoretisk Postulat. Englænderen \textsc{Joule}, der ligesom Colding
og Helmholtz ad sin egen Vej var kommen til samme Resultat som Mayer, --- et
mærkeligt Exempel paa, hvorledes Flere samtidigt kunne være paa Spor efter en
og samme Opdagelse, --- gaar %
% --- p. 499 ---
\opage{499}
overvejende experimentalt frem, men antyder dog det a priori Usandsynlige i, at
Kraft skulde tilintetgøres uden tilsvarende Virkning\textsuperscript{116}).

Den nye Sætning havde Vanskelighed ved at vinde Anerkendelse, men trængte
efterhaanden igennem i Løbet af 18 Aar (1842---1860), især da den viste sig
frugtbar ved at lede til nye Undersøgelser og Opdagelser. Dens store Betydning
bestaar netop i de bestemte Spørgsmaal, den fører Naturforskeren til at stille
hver Gang en ejendommelig Kraftytring fremtræder eller ophører. Ligesom det er
den almindelige Aarsagssætnings Betydning, at den ved enhver Forandring, der
sker, fører os til at se os om efter en forudgaaende Forandring, hvis Virkning
den nye Forandring kunde være, saaledes er det Betydningen af Sætningen om
Kraftens (eller som man nu efter engelske Forskeres Forslag hellere siger:
Energiens) Bestaaen, at man strax kommer til at stille den Opgave at faa at
vide, i hvilket Forhold de hinanden afløsende Kraftytringer staa til hinanden.

Filosofisk set maatte Spørgsmaalet især blive det, hvorledes de aandelige
Fænomeners Forhold til den nye Lov var. Det er i denne Henseende
karakteristisk, at næsten alle Opdagerne af den gik ud fra afgjort
spiritualistiske og teleologiske Anskuelser. Mayer har gentagne Gange erklæret
sig selv for bestemt Modstander af Materialismen og udtalt som sin
Overbevisning, at de videnskabelige Sandheder forholde sig til den kristelige
Religion som Bække og Floder til Verdenshavet. Paa Naturforskermødet i
Innsbruck 1869 afgav han saaledes en Erklæring i denne Retning til Forargelse
for Carl Vogt og hans Venner. Coldings Standpunkt ses af følgende Udtalelse:
»Min første Tanke om, at Naturkræfterne maatte være uforgængelige, har jeg
hentet fra den Anskuelse, at Kræfterne i Naturen maatte være beslægtede med det
Aandelige i Naturen, med den evige Fornuft saavel som med den menneskelige
Aand. Det var altsaa den religiøse Opfattelse af Livet, som ledte mig paa
Tanken om Naturkræfternes Uforgængelighed«. Colding og Joule tænke sig, at Gud
ved Skabelsen har lagt en vis Sum af Kraft ind i Naturen, og at denne Sum
hverken kan forøges eller forminskes i sin Helhed, men kun fordeles paa
forskellig Maade, ligesom %
% --- p. 500 ---
\opage{500}
Descartes i sin Tid antog om Bevægelsen. Og de antog begge, at med Energiens
Bestaaen var det Værdifuldes Bestaaen i Verden sikret, idet de ligesaa lidt som
Leibniz skelnede mellem selve Energien og dens Anvendelse til Gunst for Liv og
Udvikling\textsuperscript{117}).

Men en nærmere Undersøgelse og Begrundelse af disse Forudsætninger anstille de
nævnte Naturforskere ikke; de gik, ligesom det 17. Aarhundredes Filosofer ud
fra givne religiøse Forestillinger. En nærmere Undersøgelse af den paaviste
Lovs og idethele af den hidtil vundne Naturerkendelses Forhold til det
aandelige Livs Fænomener maatte derfor naturligt fremkaldes. Og det er da intet
Under, at Vejene deltes, og at forskellige Retninger traadte op overfor
hverandre. Givet var jo kun det, at naar en Art af fysisk Energi ophører,
indtræder en anden Art af fysisk Energi i et vist bestemt Kvantum i dens Sted.
Spørgsmaalet var, hvilken Betydning denne Erkendelse konsekvent maatte faa for
Verdensopfattelsen.

\begin{center}2. MATERIALISME.\end{center}

Den stærkeste Reaktion mod Romantikens Filosofi --- afset fra den krasse
Ortodoxi --- betegnes ved den materialistiske Literatur, der blomstrede i
Tyskland ved Aarhundredets Midte. Var hist Ideen Alt, saa blev nu Materien Alt.
Begejstring for Naturvidenskaben og de store Synspunkter, det var lykkedes den
at gennemføre ved Læren om Stoffets og Kraftens Bestaaen, kunde naturligt føre
til at betragte selve denne Lære som en tilstrækkelig Filosofi og forklare alle
Sider af Tilværelsen ved dens Hjælp. Den simpleste Maade at løse Spørgsmaalet
om den Betydning, Loven om Stoffets og Kraftens Bestaaen har for vor
Verdensanskuelse, vilde jo unægteligt være den at sige, at denne Lov indeholder
hele vor Verdensanskuelse. Den Filosofi, den moderne Materialisme opstiller,
vil ikke være Andet end den simple Konsekvens af Naturvidenskaben. Derved
syntes tillige at være vundet en fast Basis for alle vore Forestillinger og for
vor Livsførelse. I Stedet for et mystisk eller spiritual%
% --- p. 501 ---
\opage{501}
istisk Grundlag faa vi nu et aldeles haandgribeligt Grundlag at bygge paa i
Teori og Praxis. Og det er en Lære, der er let at popularisere. Den er barnlig,
anskuelig, let tilgængelig, og dens Fremstilling giver god Anledning til at
fremdrage en Mængde almeninteressante naturvidenskabelige Fakta. Flere af den
tyske Materialismes Ordførere vare dygtige og selvstændige Naturforskere,
saaledes \textsc{Carl Vogt} og \textsc{Jacob Moleschott}. Andre, som
\textsc{Louis Büchner}, virkede især ved klar, smagfuld og begejstret
Fremstilling. Ved sin populariserende Tendens har denne Literatur sin største
Fortjeneste. En Fylde af Kundskab er ved den spredt ud i store Kredse. Büchners
\textit{Kraft und Stoff} er en af de mest læste populærvidenskabelige Bøger i
vort Aarhundrede; fra 1855 til 1899 udkom 19 tyske Udgaver, og den er desuden
oversat paa flere andre Sprog. Selv om der er noget Dogmatisk ved
Materialismen, saa har den dog gjort god Gavn ved at sætte sin Dogmatisme mod
den kirkelige Dogmatisme og derved vække Eftertanken over Problemer, som ved
den romantiske Filosofis Opløsning vare trængte til Side. Den hele
materialistiske Bevægelse i Tyskland var derhos baaren af en idealistisk
Interesse for Humanitet og Fremskridt, og det var fuldt berettiget, naar
Büchner protesterede imod, at man sammenblandede Materialisme som Metode og
Teori med Materialisme i Betydning af en praktisk Livsretning. Materialismen
kan meget godt anerkende de højeste og ædleste Ideers og Følelsers Værdi ---
selv om den mener, at de ligesom alle aandelige Fænomener kun ere Produkter af
eller Former for materielle Processer.

Naar Materialismen paastaar at være den simple Konsekvens af Naturvidenskabens
Resultater, har det sin Interesse at se --- ikke blot, at Opdagerne af Loven om
Energiens Bestaaen gennemgaaende gik ud fra spiritualistiske Forudsætninger ---
men at de vigtigste Sammenstød i den materialistiske Strid fandt Sted mellem
Parter, der alle stode paa Naturvidenskabens Grund. Dette viser ikke
nødvendigvis, at Materialismen har Uret; men det viser, hvor vanskeligt det er
at drage de sande Konsekvenser, og hvor mange forskellige Motiver der faa
Indflydelse paa de Enkeltes Verdensanskuelse. Moleschotts berømte Bog \textit{Der}%
% --- p. 502 ---
\opage{502}
\textit{Kreislauf des Lebens} (1852) er rettet mod \textsc{Liebigs}
teologiserende Udtalelser i hans »Chemische Briefe«, i hvilke han særligt havde
angrebet en tidligere Sætning af Moleschott: »Ohne Phosphor kein Gedanke!«
Samtidigt udbrød en Strid mellem Fysiologen \textsc{Rudolph Wagner} i Göttingen
og Zoologen \textsc{Carl} \textsc{Vogt} i Genève. Denne Strid naaede sit
Højdepunkt, da Wagner paa et Naturforskermøde i Göttingen 1854 udtalte sig til
Gunst for Antagelsen af en æterisk Sjælesubstans, der bevæger Hjernens Fibre
som Musikeren Pianofortet og ved Deling [!] forplantes fra Forældre til Børn,
en Antagelse, han støttede paa Bibelens Lære, men indrømmede kun at kunne
fastholde ved en skarp Skelnen mellem Tro og Viden. Men naar man gaar ud over
den exakte Videns Omraade, har man efter Wagner kun en ubevislig Tro at holde
sig til, --- og han føjer til: »In Sachen des Glaubens liebe ich den
schlichten, einfachen Köhlerglauben am meisten«. Herimod optraadte Vogt med et
djervt Stridsskrift: \textit{Köhlerglaube und Wissenschaft} (1855). I et senere
Skrift \textit{Vorlesungen über den Menschen} (1863) giver Vogt en mere
videnskabelig Fremstilling af sit Standpunkt. Han lægger især Vægt paa, at
Hjernen er Bevidsthedens Organ, og at Bevidstheden forholder sig til Hjernen,
som enhver Funktion forholder sig til det tilsvarende Organ. Skønt han
indrømmer, at det ikke kan forklares, \textit{hvorledes} Bevidstheden opstaar i
Hjernens Celler, fastholder han, at den uopløseligt er knyttet til disse, og
han opgiver ikke den drillende Sætning, han allerede tidligere havde opstillet,
at Tanken staar i samme Forhold til Hjernen som Galden til Leveren eller Urinen
til Nyrerne, --- uagtet en Funktion dog staar i et andet Forhold til Organet
end et Produkt til det Sted, hvor det frembringes.

Moleschott støtter sig, som Titelen paa hans Skrift viser, især til Læren om
Stoffets Bestaaen. Den er for ham den store Tanke, som allerede det 18.
Aarhundredes Encyklopædister drog frem, som den følgende Naturvidenskab har
bekræftet, og som Fremtidens Tro maa bygges paa. Han føler Ærefrygt for det
store Kredsløb i Naturen. Bjergmanden henter i sit Ansigts Sved fosforsur Kalk
op fra Jorden --- og maaske gaar derved Stoffet til det bedste Hoved og de
højeste Tanker gennem %
% --- p. 503 ---
\opage{503}
hans Hænder: Bonden gøder sin Jord med den fosforsure Kalk; denne bliver
Bestanddel af Hveden, der nærer Menneskers Krop og Hjerne! Med Stoffet kredser
Livet gennem alle Verdens Dele, med Livet Tanken --- og af Tanken udspringer
igen Villien til at gøre Livet bedre og lykkeligere. Kunne vi blot tilføre
Organismen og Hjernen de bedste Stoffer, vil Tanken og Villien ogsaa faa den
højeste Udvikling. Naturforskeren er vor Tids Prometheus, og Kemien er den
højeste Videnskab. Det sociale Spørgsmaal finder sin Løsning, naar man blot
[sic] kan finde den rette Fordeling af de Stoffer, til hvilke Tankens og
Villiens Liv er knyttet. --- Der er her kun taget Hensyn til første Udgave af
Moleschotts Bog; i senere Udgaver er den bleven ændret og udvidet. Moleschott,
som er født i Holland 1822, virkede, da Bogen udkom, som Docent i Heidelberg,
men da man greb ind i hans Lærefrihed her, drog han til Zürich; senere virkede
han som Professor i Fysiologi i Turin og i Rom, hvor han døde 1893. Efter hans
Død er hans Selvbiografi udkommen. (\textit{Für meine Freunde.}
Lebenserinnerungen von Jacob Moleschott. Giessen 1895). Den giver et smukt
Billede af den ideelt sindede og efter alsidig Dannelse stræbende Naturforsker.
Til ret Opfattelse af hans Standpunkt tjener følgende Ytring, som fremkommer,
efterat Tankegangen i »Kreislauf der Lebens« er udviklet: »Ensidigt
»materialistisk« er dette kun for dem, der kunne forestille sig et Stof uden
Kraft eller en Kraft uden stoflig Bærer. Jeg var mig klart bevidst, at man
kunde vende hele Forestillingsmaaden om: da alt Stof er en Bærer af Kraft, er
kraftbegavet eller gennemtrængt af Aand (geistesdurchdrungen), kan man ligesaa
vel kalde den en idealistisk Anskuelse«. (S. 221). Moleschotts Standpunkt er
snarere at kalde Monisme end Materialisme: »Det drejer sig om en sand, udelelig
Toenighed, og Striden staar ikke mellem Materialisme og Spiritualisme, men
mellem Toenighed (Zweieinigkeit) og Tvedobbelthed (zwiespältige Anschauung)«.
(S. 222).

\textsc{Louis Büchner} (1824---1899) erklærer vel, at han ikke vil forklare,
hvorledes Forholdet er mellem Aand og Materie, Kraft og Stof; han vil kun
hævde, at der er et nødvendigt og uopløseligt Forhold imellem dem. Men allerede
dette er Mere, %
% --- p. 504 ---
\opage{504}
end der videnskabeligt kan bevises; og Büchner nærer desuden ingen Tvivl om, at
Aanden kun er en Egenskab ved Materien, Kraften kun en Egenskab ved Stoffet.
Impulsen til sin berømte Bog »Kraft und Stoff« havde han faaet ved Læsningen af
Moleschotts »Kreislauf des Lebens«, og det er ogsaa Stoffets Evighed, der er
det sidste Grundlag hos ham. Han mener, at Kraftens Bestaaen egentligt er en
selvindlysende Konsekvens af Stoffets Bestaaen og derfor allerede kunde være
bleven uddraget af Lavoisiers Kemi. Det var dog først i 5. Udgave af sin Bog,
at han indskød et særeget Kapitel om »Kraftens Udødelighed«, i hvilket Kraftens
Kredsløb stilles som nødvendigt Korrelat ved Siden af Stoffets Kredsløb,
saaledes at »begge tilsammen fra Evighed og i Evighed danne den Sum af
Fænomener, som vi kalde Verden«. Dog fastholder han stadigt den Antagelse, at
»i Materien bo alle Natur- og Aandskræfter; Materien er al Værens Urgrund«. Og
hvad det Aandelige angaar, saa betragter hau %
% PRINTED AS IS: „hau“ (p. 504)
det (uagtet han har fraskrevet sig al Indsigt i dets Forhold til Materien) som
et blot Produkt. »Ligesom Dampmaskinen frembringer Bevægelse, frembringer den
indviklede organiske Sammensætning af kraftbegavede Stoffer i det dyriske
Legeme en Totalsum af visse Virkninger, der, naar de ere forbundne til en
Enhed, af os kaldes Aand, Sjæl, Tanke«. (Kraft und Stoff. 7. Aufl. p. 130). I
en senere Udtalelse gaar Büchner (uden at han synes at lægge Mærke dertil) et
Skridt videre. Saaledes hedder det i en Udtalelse, der er rettet polemisk mod
Forfatteren af dette Værk: »Modsætningen mellem fysisk og psykisk Energi er kun
holdbar, saalænge man fra først af opfatter Begreberne om Legeme og Aand, eller
mere almindeligt: Stof og Kraft, paa dualistisk Maade, medens de for den
materialistisk-monistiske Opfattelse falde sammen til Et og maaske kun
fremstille \textit{forskellige Sider eller Ytringsformer for den ene og samme
Urgrund for alle Ting.} Netop Loven om Kraftens Bestaaen er det, der i
Sjæleproblemet nødvendigvis fører til materialistiske Konsekvenser, saaledes
som jeg, støttet til fysiologisk-psykologiske Kendsgerninger og Betragtninger,
tror at have godtgjort til Evidens i flere af mine Skrifter og nyligt igen i
det femte af mine Breve om » Das künftige Leben %
% --- p. 505 ---
\opage{505}
und die moderne Wissenschaft« (Leipzig 1889). Jeg tror her at have vist, at
psykisk Virksomhed \textit{ikke er eller kan være Andet end den mellem den graa
Hjernebarks Celler stedfindende Udstraaling af en ved ydre Indtryk indledet
Bevægelse«.} Af en Artikel i Ugebladet \textit{Menschthum.} Gotha 1889. Nr. 46).

Det er klart, at Büchner i denne Udtalelse sammenblander to forskellige
Teorier, den ene, der opfatter Aand og Materie som Ytringsformer for en og
samme Substans, og den anden, for hvem Materien er »Urgrunden« og psykiske
Virksomheder derfor konsekvent \textit{ikke Andet end} Bevægelse (eller
»Udstraaling af Bevægelse«). I et og samme Aandedræt bekender han sig til
begge. Den utrættelige Forkæmper for de naturvidenskabelige Resultaters
Indflydelse paa vor Verdensanskuelse har ikke set denne Inkonsekvens, som
maaske i hans Øjne ikke havde saa Meget at sige, da hans egentlige Formaal ikke
var at indskærpe en vis bestemt filosofisk Teori. I sin store Iver har han
tillige trot, at han »til Evidens har godtgjort« Mere end der kan godtgøres.

Med skarpere Kritik søgte \textsc{Heinrich Czolbe} (1819---1873), Læge ligesom
Büchner, at bestemme Konsekvenserne af Naturvidenskabens Lære. I sine første
Arbejder er han konsekvent Materialist. I \textit{Neue Darstellung des
Sensualismus} (1855) og \textit{Die Entstehung des Selbstbewusstseins} (1856),
saavel som i Afhandlingen \textit{Die Elemente der Psychologie vom Standpunkte
des Materialismus} (Zeitschrift für Philosophie. XXVI) udvikler han, at
Materialismen strider mod Naturvidenskaben, saalænge denne antager Læren om
Sansenervernes specifike Energi og altsaa antager en Forskel mellem
Sansefornemmelserne og de ydre Processer, der svare til dem. I Modsætning til
denne Lære opstiller Czolbe da den Antagelse, at det er væsentligt en og samme
Bevægelse, der forplanter sig fra Omverdenen gennem Sanseorganer og Nerver til
Hjernen. Paa intet Punkt foregaar der nogen kvalitativ Forandring, men
Forskellen mellem de forskellige Sanser er at forklare ved den forskellige
Intensitet, hvormed Bevægelsen gaar for sig i de forskellige Organer.
Bevidsthedens Enhed er at forklare ved, at Bevægelserne i Hjernen faa en i sig
selv tilbageløbende Ret%
% --- p. 506 ---
\opage{506}
ning; ved at gøre saadanne Kredsbevægelser mulige er Hjernen Bevidsthedens
Organ. Saaledes vil Czolbe kort og godt opfatte baade Fornemmelse og
Selvbevidsthed som Bevægelser i Rummet. Men han bliver --- klar og konsekvent
Tænker som han er --- opmærksom paa, at naar han uden videre gør Fornemmelse og
Selvbevidsthed til Et med Bevægelse, saa maa man ogsaa kunne vende Sætningen om
og sige, at hvor der er Bevægelse tilstede af en vis Intensitet og en vis
Retning, der er der ogsaa Bevidsthed, saa at Konsekvensen bliver Naturens
almindelige Besjæling. Saaledes førte den konsekvente Gennemførelse af
Materialismen ud over den\textsuperscript{117\,b}).

Senere (i Skriftet \textit{Die Grenzen und der Ursprung der menschlichen
Erkenntniss,} 1865, og i den overordentligt interessante Afhandling \textit{Die
Mathematik als Ideal für alle andere Erkenntniss und das Verhältniss der
empirischen Wissenschaften zur Philosophie.} Zeitschrift für exakte
Philosophie, 1866) indsaa Czolbe, at det var umuligt at forklare Tilværelsen
udfra ét Princip, hvad enten man med Büchner finder dette Princip i Materien,
eller med de spekulative Filosofer finder det i Aanden, eller med Teologerne i
Gud. En Forklaring naa vi kun ved at gaa ud fra flere Elementer, idet vi kalde
det for et Element, som vi ikke formaa at analysere. Saadanne Elementer, der
ikke igen kunne reduceres til hverandre, ere de materielle Atomer, de organiske
Kræfter og de sjælelige Elementer (hvilke sidste i en Sum udgøre
Verdenssjælen). Mellem disse tre Arter af Elementer finder der en harmonisk
Samvirken Sted, ved hvilken en hensigtsmæssig Natursammenhæng virkeliggøres.
Ikke i sin Oprindelse, men i sin Retning og Tendens godtgør Tilværelsen sig som
en Enhed. --- Saaledes viste det sig, at Tilværelsesproblemet var mere
indviklet, end Materialismen antog. I sine senere Undersøgelser, som kun
tildels ere udgivne, nærmede Czolbe sig til Spinozas Grundtanker, af hvilke han
vilde give en empirisk Udførelse. --- Om de forskellige Udviklingstrin af den
energiske Tænkers Filosoferen maa her henvises til \textsc{Vaihinger's}
Afhandling \textit{Die drei Phasen des Czolbeschen Naturalismus}
(Philosophische Monatshefte XII).

Hvad der fra først til sidst ligger til Grund hos Czolbe er %
% --- p. 507 ---
\opage{507}
Fordringen om Klarhed og Anskuelighed i alle Grundbegreber. Alt Mystisk og
Oversanseligt vilde han have bort og søgte saa vidt som muligt at tænke alle
Ideer med geometrisk Klarhed. Det var mere med Hensyn til Anskueligheden end
med Hensyn til den stringente Bevisførelse, at han proklamerede Matematiken som
Videnskabens Ideal. Han var i sin Ungdom begejstret for Hölderlins Digtning og
havde sat sig den Opgave at hævde den græske Klarhed og Livsglæde mod al Mystik
og Dualisme. Anskuelighed paa Tænkningens Omraade var for ham nøje beslægtet
med Glæden over den naturlige Verden. Ligesom det kan kræve Arbejde at
gennemføre Anskuelighedens Princip i Tænkningen, saaledes kan det koste
Resignation at gennemføre Glæden ved det virkelige Liv i Praxis; men baade paa
det teoretiske og det praktiske Omraade vil Czolbe bekæmpe hvad han kalder »den
Blödsinn der Transcendenz«. Han indrømmede allerede i sin materialistiske
Periode, at det inderste Motiv for hans Filosofi var den subjektive Trang til
Anskuelighed og til Fastholden af det Virkelige; Materialismen var ham et
Postulat --- der senere afløstes af andre Postulater. ---

\textsc{Ernst Haeckel} (f. 1834, fra 1865 Professor i Zoologi i Jena) regnes
ofte med til Materialismen, men har selv bestemt sit Standpunkt som en Monisme,
der fører udover Modsætningen mellem Spiritualisme og Materialisme og lægger
den store panteistiske Grundtanke om Naturens Enhed til Grund. For Monismen,
siger han, existerer der hverken Aand eller Materie i sædvanlig Betydning, men
kun Et, som paa én Gang er begge Dele. Det var i sin \textit{Generelle
Morphologie} (1862---1866), at han først udtalte sine almindelige filosofiske
Anskuelser. Det Sjælelige er for ham (ligesom for Czolbe) et oprindeligt
Element i Tilværelsen, skønt det existerer i overordentligt forskellige Grader,
ligefra Atomsjælen og Cellesjælen til de højeste Organismers Sjæle. Betænkelig
bliver denne Besjælingsteori hos Haeckel derved, at han ofte lader Sjælens
Indvirken tjene som Forklaring af organiske Bevægelser i Stedet for at søge en
rent naturvidenskabelig Forklaring. Paa den anden Side erklærer han undertiden
sjælelige Fænomener for blotte Virkninger af de mest sammensatte og mest
ustadige Kulstofforbindelser, Ægge%
% --- p. 508 ---
\opage{508}
hvidemolekulerne i Nervevævet. Han mener at være ude over Spiritualisme og
Materialisme, og dog finder man hos ham baade Sætninger, der lyde
spiritualistisk, og Sætninger, der lyde materialistisk. Han er ikke kommen til
sine Principer gennem skarp Gennemtænken af Lovene om Stoffets og Kraftens
Bestaaen. Det er Darwins Lære, der har sat hans Tænkning i Gang, og han har
søgt et højeste Princip, der kunde forenes med den store Enhed i Naturen, den
nye Lære pegede hen paa. Monismen fører efter Haeckel (Gen. Morphol. II. p.
445---451) til den mest ophøjede Gudsforestilling. I den store, Alt omspændende
Aarsagslov aabenbarer Guddommen sig som omfattende den hele Natur og virkende i
ethvert Naturfænomen. Den sædvanlige Gudsopfattelse er derimod Tro paa to
Guder, Amfiteisme, ikke Monoteisme, fordi den sætter Naturaarsagerne ved Siden
af Guddommen,

Haeckel var en af de første tyske Naturforskere, der sluttede sig til Darwin,
og han kæmpede med stor Fyrighed for dennes Teori, som han tillagde en
Sikkerhed og gav en Udvidelse, der ikke altid kunde billiges af Teoriens
kritiske og besindige Grundlægger. Han var utrættelig i at konstruere de
nulevende Arters Stamtræer, og han nærede ingen Betænkelighed ved at antage en
stadig Opstaaen af organisk Materie. Han har ikke noget klart Blik for
Hypotesers Begrænsning og for Nødvendigheden af Verifikation; kun derved
forklares, at han uden videre kan stille Darwins Hypotese sammen med Newtons.
Han bebrejdede da ogsaa Darwin, at denne lagde for stor Vægt paa Indvendingerne
mod Teorien. Den store Forsker rystede paa Hovedet ad sin unge Tilhængers Iver:
Your boldness sometimes makes me tremble, skrev han til ham (19. Novbr. 1868).
--- I Haeckels seneste »filosofiske« Skrift \textit{(Die Welträthsel.}
Gemeinverständliche Studien über monistische Philosophie. Bonn 1899) fremtræde
de omtalte Egenskaber i endnu grellere Form end i hans tidligere Udtalelser. ---

Læren om Stoffets og Kraftens Bestaaen og Læren om Arternes Udvikling gennem
naturlig Udvælgelse vare --- som allerede bemærket --- de naturvidenskabelige
Resultater, der i Perioden efter Aarhundredets Midte især maatte sætte Efter%
% --- p. 509 ---
\opage{509}
tanken i Bevægelse. Efterat vi nu have betragtet de Forsøg, der af
Naturforskere bleve gjorte paa at grundlægge Filosofi paa disse Resultater,
skulle vi gaa over til at undersøge den Maade, paa hvilken man fra
Fagfilosofers Side stillede sig til dem.

\begin{center}3. IDEALISTISKE KONSTRUKTIONER PAA REALISTISK BASIS.\end{center}

\begin{center}a. RUDOLPH HERMANN LOTZE.\end{center}

Den idealistiske Filosofi har i Lotze sin vigtigste Repræsentant i den sidste
Halvdel af Aarhundredet. I sin Personlighed og i sin Udviklingsgang staar han
som den, der paa en interessant Maade har optaget i sig de ideale Motiver, som
Romantikens Filosofi byggede paa, saavel som den strenge Gennemførelse af den
mekaniske Naturopfattelse, der henimod Aarhundredets Midte arbejdede sig frem i
Videnskaben. Lotze er en Mester i Begrebsudvikling, i at drøfte en Tanke
saaledes, at alle dens Nuancer komme frem, i atter og atter at tage et Problem
op og behandle det fra forskellige Synspunkter. Hans Ideal var i sidste Instans
det samme som det, der foresvævede Romantikens Filosofi: at udlede al Udvikling
og al Sammenhæng i Verden af en evig Ide, der skulde indeholde den sidste Grund
for Alt, hvad der sker, saavel som for Værdien af det, der sker. Men han saa
klart, at en saadan Udledning overstiger, hvad den menneskelige Tanke kan yde,
og at naar Romantikens Tænkere alligevel mente at have givet en saadan
Udledning, var det, fordi den poetisk-religiøse Tendens hos dem uvilkaarligt
gik i Et med den filosofiske. Lotze udskiller derfor det poetisk-religiøse
Element fra det spekulative. Hans egen fine poetiske Sans lod ham forstaa, at
deres romantiske Sammenblanding ikke var heldig. Og denne poetiske Sans er Et
med en Sans for individuelle Nuancer og Forhold, som den spekulative Filosofi
altfor ofte forflygtigede ved sine Abstraktioner. Gennem denne Sans hænger
Lotzes spekulative Interesse sammen med hans real%
% --- p. 510 ---
\opage{510}
istiske Tendens. Det gælder for ham om først og fremmest at opfatte Fænomenerne
i deres konkrete Natur og i deres bestemte, lovmæssige Sammenhæng; det er da
den filosofiske Tænknings Opgave at udfinde de Forudsætninger, paa hvilke denne
reale Sammenhæng hviler. Det digteriske, det naturvidenskabelige og det
filosofiske Element hos Lotze hænge saaledes nøje sammen, og han har været
udrustet som Faa til at behandle den Opgave, han havde sat sig, og som Tidens
aandelige Vilkaar lagde nær: at forsøge en Rekonstruktion af en idealistisk
Filosofi paa realistisk Basis. Det var hans Overbevisning, at Romantikens
Filosofi havde overset de reale Betingelser og den mekaniske Natursammenhæng,
uden hvilke selv de betydningsfuldeste Ideer ere hjælpeløse Idoler, ligesom det
ogsaa var hans Overbevisning, at Materialismen gjorde \textit{det} til Et og
Alt, som kun var en --- ganske vist nødvendig --- Form og Ramme for
Tilværelsens værdifulde Indhold. Hans Filosofi er væsentligt en Analyse af
selve Naturmekanismens Begreb, der skal vise, at dette Begreb med Nødvendighed
fører til Antagelsen af et ideelt Tilværelsesprincip, og at det i hvert
Tilfælde ikke udelukker Antagelsen af, at et saadant Princip er den evige Kilde
til alt Godt og Værdifuldt i Verden.

I Lotzes Udviklingsgang fremtræde de forskellige Motiver og Interesser, der
besjælede ham, klart. Han er født i Bautzen i Sachsen den 21. Maj 1817, i den
samme Egn, fra hvilken Lessing og Fichte stamme. Ved Universitetet i Leipzig
studerede han Filosofi, Medicin og Fysik. Han indviedes her i de to
Tankekredse, det senere skulde blive hans Opgave at forene, efter at have
gennemarbejdet dem hver for sig. Det var Æstetikeren og Religionsfilosofen Chr.
Hermann Weisse, som var hans Lærer i Filosofi, og Lotze har senere udtalt, at
han ikke blot skylder Weisse mange enkelte Paavirkninger, men at han særligt
blev indført af ham i en Kreds af Tanker, han senere ikke har følt Trang til at
opgive. Weisse var den mest fremragende Repræsentant for den filosofiske
Teisme. Gennem ham staar Lotzes religionsfilosofiske Opfattelse i historisk
Forbindelse med Schellings --- og derigennem med den gamle Jacob Böhmes. Lotze
fremhæver i et Tilbageblik paa sin Udvikling \textit{(Streitschriften.} %
% --- p. 511 ---
\opage{511}
1857. p. 5 ff.), at han tilegnede sig den spekulative Filosofis Ideer ikke som
sluttet Lærebygning, men som en ejendommelig Form for aandelig Dannelse. Den
Forudsætning, at Tingenes sidste Grund kun kunde tænkes som aandelig, er
tidligt bleven grundlagt hos ham og aldrig sluppen. Medicin og Fysik studerede
han under E. H. Weber, Volkmann og G. Th. Fechner. Her lærte han den
naturvidenskabelige Metode og Opfattelse at kende paa første Haand. I et og
samme Aar tog han Doktorgraden i Filosofi og Medicin og optraadte nu som Docent
i begge Fag, ogsaa efterat der 1842 var overdraget ham et Professorat i
Filosofi. Efterat have virket nogle Aar i Leipzig blev han Professor i
Gøttingen, som Herbarts Efterfølger, og her falde hans vigtigste Arbejder.
Aaret før sin Død modtog han en Kaldelse til Berlin, men bukkede her snart
under for en Sygdom, han længe havde lidt af (1881). Hans Liv var stille, viet
Studium, Tænkning og akademisk Undervisning. Med sjelden Forening af Alsidighed
og Grundighed forstod han at orientere sig paa højst forskellige Omraader,
Noget, hvorom ikke blot hans medicinske og filosofiske Værker, men ogsaa en hel
Række mindre Afhandlinger og Anmeldelser (udgivne efter hans Død i fire Bind
under Titel \textit{Kleine Schriften)} vidne. Fra sine strengere videnskabelige
Arbejder udhvilede han sig i Beskæftigelse med Kunst og Literatur. Saaledes
oversatte han lejlighedsvis »Antigone« paa latinske Vers.

Som medicinsk Forfatter satte Lotze sig, som allerede tidligere berørt, til
Opgave at hævde Fysiologiens Karakter som mekanisk Naturvidenskab. De organiske
Fænomeners Ejendommelighed vil han ikke have forklaret ved Paaberaabelse af en
mystisk Livskraft, men ved Paavisning af den bestemte, lovmæssige Maade, paa
hvilken de almindelige Naturkræfter virke indenfor Organismen. Som overalt i
Naturen, saaledes maa vi ogsaa her finde Forklaringen i reale Elementers
Vexelvirkning. Ikke ved at være unddraget den mekaniske Natursammenhæng, men
ved den særegne Maade, paa hvilken sammenhængende Rækker af Virkninger opstaa,
adskiller det organiske Liv sig fra den uorganiske Verden. I den spekulative
Naturfilosofi var Ideen om Organismen gjort til Type for hele Verdensopfattelsen; %
% --- p. 512 ---
\opage{512}
at hin Ide blev korrigeret, maatte da blive af almindelig filosofisk Betydning,
og man forstaar, at ikke blot Lotzes medicinske, men ogsaa hans filosofiske
Interesse tilfredsstilledes ved hans Arbejder i denne Retning.
\textit{(Allgemeine Pathologie und Therapie als mechanische
Naturwissenschaften.} 1842. --- \textit{Allgemeine Physiologie des körperlichen
Lebens.} 1851). Disse Arbejder gave endog nogle Materialister Anledning til at
se en Meningsfælle i Lotze, uagtet denne bestemt udtalte, at Mekanismen for ham
kun var en Del af hans Naturopfattelse, ikke den hele. Det er karakteristisk
for Lotze, at han saaledes deler Problemerne. Sine almindelige filosofiske
Anskuelser havde han udviklet i sin \textit{Metaphysik} (1841). Men han følte
sig ingenlunde færdig, da han begyndte sin Lærervirksomhed, og særligt hvad de
afsluttende Ideer angaar, henviste han stadigt til fremtidige Arbejder. I
\textit{Medicinische Psychologie oder Physiologie der Seele} (1852) udtalte han
sig udførligt om Forholdet mellem det Aandelige og det Materielle og gav
tillige psykologiske Undersøgelser af blivende Interesse. En længe næret Plan
udførte han i den følgende Tid (1856---1864) ved Udgivelsen af sin
\textit{Mikrokosmus} (3 Bd.). Dette Værk, der er tænkt som et Sidestykke til
Humboldts »Kosmos« og Herders »Ideen«, giver en Psykologi i nøje Sammenhæng med
Fysiologi og Kulturhistorie og ender med en Udvikling af Lotzes kosmologiske og
religionsfilosofiske Ideer. For alle sine Interesser, alle de forskellige Veje,
hans Tanke og Følelse havde vandret, havde han her fundet Udtryk. Værket er
anlagt som populær Fremstilling og har vundet en betydelig Udbredelse. Efter
Udgivelsen af en aandfuld \textit{Geschichte der Aesthetik in Deutschland}
(1868) skred Lotze til at give en afsluttende systematisk Fremstilling af sin
Filosofi. Men kun to Dele af den lykkedes det ham at faa færdige \textit{(Drei
Bücher der Logik.} 1874. --- \textit{Drei Bücher der Metaphysik.} 1879). Det
tredie Bind, der skulde omfatte Æstetik, Etik og Religionsfilosofi, blev ikke
udarbejdet. Hans Filosofi er altsaa ufuldendt; det blev ham ikke forundt at
sætte Kronen paa Værket og vise, hvorledes de mange Tanketraade, han havde
spundet, kunde væves sammen til Enhed. --- Korte Fremstillinger af Lotzes Lære
paa de forskellige filosofiske Om%
% --- p. 513 ---
\opage{513}
raader har man i de efter hans Død udgivne Diktater \textit{(Grundzüge)} fra
hans Forelæsninger. Blandt disse maa især fremhæves \textit{Grundzüge der Logik
und Encyclopädie der Philosophie,} som (i sin anden Del) giver en kort og klar
Oversigt over hele hans System. Her give vi en Karakteristik af de tre
vigtigste Punkter i hans Filosofi.

\begin{center}\textit{α. Den mekaniske Naturopfattelse.}\end{center}

Lotzes Tænkning har to Udgangspunkter: en levende Følelse af det aandelige Livs
Værdi, af at det Højeste for os er knyttet til den aandelige Udvikling og dens
Idealer --- men tillige en fast Overbevisning om, at et System af mekaniske
Aarsager og Love er nødvendigt til Virkeliggørelse selv af de højeste Idealer.
Han kan, som han et Sted siger, tro paa virkende, men ikke paa hexende Ideer.
De to Udgangspunkter staa dog hos ham ikke helt isolerede. Han søger ved at gaa
ud fra det ene at vise Muligheden af at anerkende det andet. Det mest
karakteristiske og det mest videnskabelige Træk i hans Filosofi er hans Forsøg
paa gennem Analyse af selve Mekanismens Begreb at godtgøre, at det til sidste
Forudsætning har Begrebet om et Princip, der --- saasnart vi opfatte det paa en
for os forstaaelig Maade --- viser sig som Bærer af og Kilde til de højeste
Ideer. Lotze vil tage Tyren ved Hornene: ved at drage Realismens Konsekvenser
vil han grundlægge Idealismen. Kun Undersøgelser, der føres i Realismens Aand,
kunne efter hans Opfattelse bringe os nærmere henimod det Maal, Idealismen
sætter sig: Erkendelsen af Verden som Udtryk for en værdifuld Ide. Romantikens
Filosofi havde trot at kunne udlede Virkelighedens Former af den højeste Ide.
Dette har vist sig at være en uløselig Opgave. Derimod kunne vi forsøge ad
Slutningens Vej at gaa tilbage fra det Givne til dets Forudsætninger. Ikke en
\textit{Deduktion,}men \textit{Reduktion} er mulig. Tænkningen maa stedse
anvende sine Former paa et givet Stof. Baade den almindelige Dannelse og de
enkelte Videnskaber operere her med en Mængde Begreber, hvis Oprindelse,
Betydning og Gyldighed de ikke nærmere undersøge. Saadanne Begreber ere Aarsag
og Virk%
% --- p. 514 ---
\opage{514}
ning, Stof og Kraft, Formaal og Middel, Frihed og Nødvendighed, Materie og
Aand. Det er nu Filosofiens Opgave at bringe Enhed og Sammenhæng i
Forestillingskredsen ved at gøre disse i det praktiske Liv og i de enkelte
Videnskaber forudsatte Begreber til Genstand for en særegen Undersøgelse og
bestemme Grænserne for deres Gyldighed.\textit{ (Grundzüge der Logik und
Encyclopädie der Philosophie.} § 88).

Af disse Begreber er nu det vigtigste det, som vi forudsætte ved enhver
Undersøgelse af Virkeligheden, --- hver Gang vi foretage et Experiment, hver
Gang vi søge en Forklaring: Begrebet om en gennemgaaende, altomfattende
Aarsagssammenhæng. Dette Begreb er ikke grundet paa Erfaring, men forudsættes
ved enhver Erfaring. Men da den Forstaaelse, som hidtil er vunden af Naturen,
beror paa dets Realitet, saa kan det selv betegnes som Udtryk for en
Kendsgerning, og denne Kendsgerning, at det enkelte Element i vor Erfaring ved
en lovmæssig Sammenhæng er knyttet til andre Elementer, --- den mekaniske
Sammenhængs Kendsgerning, --- er det da, Filosofien skal gennemtænke og drage
alle Konsekvenserne af. Deducere den kan den ikke; men den kan maaske opdage,
hvad der ligger i den. En \textit{Mangfoldighed af reale Elementer i indbyrdes
Vexelvirkning} er det Grundlag, paa hvilket den mekaniske Naturopfattelse
bygger, og Lotze var, som vi have set, saa overbevist om dette Grundlags
Uundgaaelighed, at han selv arbejdede med Iver for at skaffe det Anerkendelse
indenfor Fysiologien, hvor hidtil Vitalismen havde raadet med sin Appel til en
eneste formende og styrende »Livskraft«.

Men fordi den mekaniske Sammenhæng er et uundgaaeligt Træk i vor
Verdensopfattelse, behøver den ikke at være det eneste, Alt bestemmende Træk.
Tvertimod: selv om dens Gyldighed er ubegrænset, kan den godt have en
underordnet Betydning i Verden som Helhed. (Die unbeschränkte Gültigkeit des
Mechanismus, aber zugleich seine durchaus untergeordnete Bedeutung im Ganzen
der Welt! \textit{Drei Bücher der Metaphysik.} p. 462). Og en nærmere Analyse
af Mekanismens Begreb fører os til at indse, at saaledes forholder det sig netop.

Den mekaniske Naturopfattelse stanser --- naar den gøres %
% --- p. 515 ---
\opage{515}
til afsluttende Verdensanskuelse --- ved en Mangfoldighed af Elementer (Atomer)
i indbyrdes Vexelvirkning. Den proklamerer en Pluralisme. Men hvorledes
forholde Elementerne sig til den Sammenhæng, i hvilken de staa? Kunne de have
selvstændig Bestaaen afset fra denne Sammenhæng, saa at denne er en for deres
Væsen ligegyldig Relation, eller bestemmes de ikke netop helt igennem ved den
Forbindelse, i hvilken de staa med Verden som Helhed? Vexelvirkning og
Sammenhæng kunne dog ikke svæve frit over eller mellem Elementerne, men
forudsætte en indre Enhed af disse. Ti antager jeg, at de to Elementer A og B
ere aldeles selvstændige, saa bliver Vexelvirkning mellem dem ubegribelig. En
Virkning kan ikke overføres som fuldt færdig Tilstand fra A til B. Hvad der
gaar for sig i A, kan kun faa Betydning for det, der gaar for sig i B, naar A
og B ikke i Virkeligheden ere selvstændige, absolut sondrede Væsener, men naar
deres Tilstande ere Tilstande i et og samme, mere omfattende Væsen. En
Mangfoldighed af uafhængige Væsener vilde gøre den mekaniske Vexelvirkning
uforstaaelig; Forstaaeligheden kommer kun med Antagelsen af et uendeligt,
altomfattende Væsen, hvis Momenter eller Aktionspunkter de enkelte Elementer
ere. Begrebet om en Krafts eller en Indflydelses »Overgang« fra et selvstændigt
Element til et andet er uholdbart. Ikke en \textit{overskridende} Aarsag (causa
transiens), kun en \textit{iboende} Aarsag (causa immanens) er forstaaelig. Et
og samme Væsens Tilstande kunne forholde sig til hverandre som Grund og Følge,
men to indbyrdes uafhængige Væseners Tilstande ikke. --- Blandt de talrige
Fremstillinger, Lotze har givet af denne Tankegang, maa fremhæves \textit{Logik
und Encyclopädie} § 99---100 og \textit{Drei Bücher der Metaphysik} § 50 --- 81.

Ved Analyse af Aarsagsforholdets og Vexelvirkningens Begreber ---
Grundbegreberne i den mekaniske Naturopfattelse --- er Lotze altsaa naat til
Ideen om en Ursubstans, et altomfattende Princip. Hans Tankegang bevæger sig
her ad samme Veje, som i sin Tid førte Spinoza til sit Substansbegreb, og som
Kant vandrede paa især i sine Ungdomsskrifter uden dog at forlade dem i sine
senere Arbejder. Ogsaa om Leibniz minder han paa flere Punkter, uden at der dog
her, ligesaa lidt %
% --- p. 516 ---
\opage{516}
som ved Spinoza, synes at have fundet nogen Paavirkning gennem Studium Sted. Da
han hørte, at man stillede ham sammen med Spinoza og Leibniz, udtalte han (se
\textsc{Falckenberg}: \textit{Hermann Lotze}. Stuttgart 1901, p. 85), at han
ikke følte sig at staa i noget indre Forhold til dem, skønt han indrømmede, at
der var Anledning nok til at anstille en Sammenligning.

Det Begreb, som Romantikens Filosofi vilde lægge til Grund, og af hvilket den
vilde udlede Alt, bliver for Lotze det sidste Postulat, eller som han ogsaa
udtrykker sig, den sidste Kendsgerning for vor Tænkning. En nærmere, anskuelig
Udvikling af det universelle Princip, som forudsættes ved det simpleste
Vexelvirkningsforhold, er umulig. Vi staa her ved et \textit{Grænsebegreb,} som
vi hverken kunne undvære eller udvikle. (Drei Bücher der Metaphysik. § 73 og
246). Det er dog dette Begreb, som gør det muligt for Lotze at fastholde, hvad
der for ham var det Væsentlige i den idealistiske Tankegang, han af sin Lærer
Weisse var bleven indført i. Ved det fandt han en Forbindelse mellem
Tankeverdenens modsatte Strømninger. Hverken Materialismens absolute Atomer,
Leibniz' Monader eller Herbarts Realer kunde for ham betegne Tænkningens
Afslutning: Pluralismen maa nødvendigvis give Plads for \textit{Monismen.} Hvad
særligt Atomerne angaar, hævder Lotze, at Naturvidenskabens Interesse kun fører
til at antage Elementer, som for vor Erfaring faktisk ere uopløselige, men at
Antagelsen af en Mangfoldighed af udstrakte Elementer --- selv om disse tænkes
forsvindende smaa --- ikke kan være nogen Afslutning for Tanken. Enten maa man
opgive Atomets Enhed eller dets Udstrækning: indenfor et udstrakt Atom vil
enhver Virkning tage Tid og gaa fra Del til Del --- og disse Dele ville da være
mere fundamentale Enheder end det hele »Atom«. Ved det afsluttende Atombegreb
maa vi se bort fra al Udstrækning og tænke os Atomerne som Kraftcentra, der ---
ifølge den ovenfor givne Analyse af Mekanismens Begreb --- hvert for sig ere
Udgangspunkter for Ursubstansens Virken. (Drei Bücher der Metaphysik. §
190---191 og 245, kfr. den interessante Anmeldelse af Fechners »Atomenlehre« i
»Kleine Schriften« III. p. 215---238). %
% --- p. 517 ---
\opage{517}

Lotzes Hævden af den mekaniske Naturopfattelse og de Konsekvenser, han drager
af denne, udgøre den betydningsfuldeste Del af hans Filosofi, uagtet den i
Regelen ikke drages frem, naar Lotze prises. Det er den spiritualistiske og
teologiserende Retning hos ham, der hyppigst betones, baade af Beundrere og af
Modstandere. Og dog har han indlagt sig den største Fortjeneste som Tænker ved
sin Analyse af Grundbegrebet i den videnskabelige Naturopfattelse. En Mangel er
det hos ham, at han ikke anstiller en erkendelsesteoretisk Undersøgelse af det
Grænsebegreb, han kommer til, men lader det staa i den Form, i hvilken
Romantikens Filosofi og den tidligere Dogmatisme havde opfattet det. Lotze
havde overhovedet ingen Forstaaelse af Erkendelsesteoriens Betydning.

\begin{center}\textit{β. Metafysisk Idealisme.}\end{center}

Ved Lotzes Tankegang, saavidt vi hidtil have fremstillet den, er der endnu en
vis Ubestemthed. Der kunde endnu spørges, om vi da aldeles ingen Udvej have til
at danne os en nærmere Forestilling om Elementerne og Ursubstansen. Idet vi nu
fremdrage Lotzes Svar paa dette Spørgsmaal, maa det allerførst fremhæves, at
han er aldeles klar over, at enhver Besvarelse af det støtter sig til Analogier
og bestemmes ved andre Motiver end de rent teoretiske. Den Analogiens Lov og
den Trang til i Tilværelsen at genfinde vort eget aandelige Væsen, som have
vist sig at ligge til Grund hos alle metafysiske Idealister (især Leibniz,
Berkeley, Schelling, Schopenhauer, Beneke), de fremtræde med fuld Bevidsthed
hos Lotze.

Lotze er Atomist --- men han opfatter ikke Atomerne som materielle, fordi
Udstrækningen ligesaa vel som alle andre sanselige Kvaliteter maa forklares ved
Vexelvirkning af Atomer, der altsaa selv ikke kunne have denne Egenskab.
Ligesom Livet og alle Erfaringskvaliteter skyldes Udstrækningens sanselige
Faktum et Sammenspil af Kraftpunkter, der igen ere at opfatte som
Begyndelsespunkter for det uendelige Urvæsens indre Virken. Lotze anser det
ikke for umuligt, at helt nye Begyndelser kunne indtræde under Verdensløbet:
det strider %
% --- p. 518 ---
\opage{518}
ikke mod den universelle Lovmæssighed, ti Loven udtrykker stedse kun den
Rækkefølge, i hvilke Tilstandene opstaa, ikke en ydre Skæbne. Elementernes
Antal behøver ikke at være konstant, men kan tænkes at stige eller falde, alt
efter det højeste Verdensformaals Krav. (Drei Bücher d. M. § 33 og 235). Og de
enkelte Elementer behøve heller ikke indbyrdes at være aldeles ensartede. En
vis Overensstemmelse eller Kommensurabilitet maa der vel være tilstede for at
en Verdensorden kan omfatte dem alle; men en fuldstændig Lighed er ikke
nødvendig. En Naturlov kan meget vel sammenknytte kvalitativt forskellige
Elementer, som vi ikke formaa at finde nogen Generalnævner for. Det er tilsidst
ikke en \textit{logisk,} men en \textit{æstetisk} Nødvendighed, som gør os
Tilværelsen forstaaelig (Drei Bücher d. M. § 59): ikke en formel Konsekvens i
Urvæsenets Virken, men den Rigdom og Fylde, hvormed Urvæsenets Virken foregaar,
er bestemmende for de enkelte Elementers Beskaffenhed. Paa dette Punkt gaar
Lotze tydeligt nok over fra de rent teoretiske Motiver til Følelsesmotiver,
uden at dog hverken disse eller hine formaa at føre ham til bestemtere
Fastsættelse af, hvor store \textit{Kvalitetsforskelligheder} der ere
forenelige med den \textit{Kommensurabilitet,} som Vexelvirkningens Faktum
forudsætter\textsuperscript{118}).

Ville vi nu idethele danne os en Forestilling om Elementernes indre Natur, saa
maa vi opfatte dem i Analogi med vort eget aandelige Væsen. Den mekaniske
Naturopfattelse belærer os egentligt kun om det indbyrdes Forhold mellem
Elementerne, ikke om disses indre Natur. Den angaar de ydre Omstændigheder,
bevæger sig hele Tiden udenom Tingene, --- er en cognitio circa rem. Den er
maaske endog --- ligesom den populære Verdensopfattelse --- tilbøjelig til at
mene, at det er aldeles ligegyldigt for Tingene, om vi opfatte dem eller ikke.
I Modsætning hertil hævder Lotze (allerede i sit Ungdomsarbejde
\textit{Metaphysik.} 1841. p. 313), at Subjektiviteten ligesaa vel er en Del af
Realiteten, som de ydre Objekter: Nicht nur was \textit{zwischen} den Wesen,
sondern auch was \textit{in} ihnen sich zuträgt, ist ein wahres und wirkliches
Geschehen. I sine senere Skrifter (først i \textit{Medicinische Psychologie)}
gaar han et Skridt videre, %
% --- p. 519 ---
\opage{519}
idet han hævder, at vort eget subjektive Væsen er det eneste Tilfælde, i
hvilket vi kende en Tings Indre og have en cognitio rei, ikke blot en cognitio
circa rem: den eneste Vej til at danne os en Forestilling om andre Tings Indre,
er da at opfatte dem i Analogi med os selv --- altsaa som følende Væsener (om
ikke som forestillende, da Følelse er en mere oprindelig Bevidsthedsytring end
Forestilling). Benyttelsen af denne Analogi er det eneste Middel til at tænke
Tingene som reale, for sig bestaaende Væsener, der ere Mere end blotte Billeder
for os. Vi gøre derved blot Brug af den almindelige Metode, der fører det
Ubekendte tilbage til det Bekendte. (Drei Bücher d. M. § 96---98). Og den
»æstetiske Nødvendighed« kræver ogsaa dette, da den ikke kan forenes med en
Tilværelse, i hvilken en stor, maaske den største Del blot skulde være et
dunkelt Grundlag for et kun til enkelte Regioner udstrakt Sjæleliv (se allerede
\textit{Allg. Physiologie} § 129). Tilværelsens Elementer maa altsaa antages
besjælede i mangfoldige Grader.

Som det forholder sig med Tilværelsens Elementer, forholder det sig ogsaa med
den altomfattende Verdensgrund, der gør deres Vexelvirkning forstaaelig. Følge
vi Analogiens Lov og forsøge vi at føre det Ubekendte tilbage til det Bekendte,
saa maa vi ifølge Lotze opfatte Verdensgrunden som et aandeligt Væsen. Naar
indre Tilstande skulle være forbundne til Enhed og Sammenhæng, saa er vort eget
indre, aandelige Liv det eneste Exempel, vi kende, i hvilket Muligheden af en
saadan Bevarelse af Enheden under de forskellige Tilstande er virkeliggjort ---
idetmindste tilnærmelsesvis virkeliggjort. I Tilslutning til sin Lærer C. H.
Weisse opfatter Lotze Verdensgrunden som en absolut Personlighed, og han
forsvarer denne Overførelse af Personlighedens Begreb paa det absolute Væsen
derved, at kun et absolut Væsen kan være Personlighed, fordi kun det har indre
Selvstændighed og Oprindelighed, medens Personlighedens Begreb kun ufuldkomment
finder Anvendelse paa endelige Væsener, der ere afhængige af ydre Betingelser.
Ganske vist indrømmer Lotze, at til et personligt Liv hører der Modstand at
overvinde og Evne til at lide og modtage saavel som til at virke. Men naar man
da spørger, hvorledes et absolut %
% --- p. 520 ---
\opage{520}
Væsen, som ingen Skranker er underlagt, kan lide, svarer Lotze
\textit{(Grundzüge der Religionsphilosophie} § 34, kfr. § 52), at det, der
sætter Guddommens Følelse i Bevægelse, maa være »indre Frembringelser af dens
egen skaberiske Fantasi!« --- Det er dog vel et stort Spørgsmaal, om en saadan
selvskabt Modstand kan betyde noget Alvorligt --- især da den kan
tilintetgøres, i hvilket Øjeblik det skal være. Det personlige Liv, vi kende,
har i hvert Tilfælde netop ikke selvskabte og ikke saa let bortryddelige
Skranker at kæmpe med, og den Analogi, til hvilken Lotze støtter sig, synes
derfor at svigte netop paa det afgørende Punkt. Og dertil kommer, at Lotze ---
efter den sandsynligste Fortolkning af hans uklare og skiftende
Ytringer\textsuperscript{119}) --- skiller sig fra Weisse deri, at han ikke
antager Tidsformen for anvendelig paa det absolute Væsen. Et personligt Væsen,
der ikke udvikler sig i Tiden, --- et tidløst Liv og en tidløs Liden og Virken:
det er dog Begreber, der stille altfor store Fordringer til vor Evne til at
analogisere! --- Endnu skal kun bemærkes, at Lotze --- baade hvad Elementerne
og Verdensgrunden angaar --- forudsætter, at Modsætningen mellem Aand og
Materie er kontradiktorisk, saa at der ingen flere Muligheder skulde kunne
tænkes end de to: kun under denne Forudsætning er der tvingende Nødvendighed
for at opfatte Elementerne og Verdensgrunden som aandelige Væsener, naar man
ikke kan eller vil opfatte dem som materielle. Men der er ingen logisk, kun en
faktisk Nødvendighed for at vælge mellem Aand og Materie ved Bestemmelsen af de
afsluttende Begreber. Vi kende faktisk kun disse to Former for Tilværelse; men
der vilde ikke være noget \textit{Utænkeligt} i, at der existerede en eller
flere andre Former. Derfor kan den metafysiske Idealisme ikke
bevises\textsuperscript{120}).

Lotze er sig nu ogsaa fuldt bevidst, at Afslutningen af hans hele
Verdensanskuelse beror paa Følelsesmotiver, ikke paa streng Tænkning.
Nødvendigheden af at gaa tilbage fra Elementernes Mangfoldighed til
Verdensgrundens Enhed hævder han vel af rent teoretiske Grunde: »Gar kein
Weltlauf, weder ein harmonischer noch ein unharmonischer, ist mir ohne die
Voraussetzung jener Einheit begreiflich, die alles gegenseitige %
% --- p. 521 ---
\opage{521}
Wirken des Einzelnen erst möglich macht; die Störungen der Dinge durch einander
bezeugen die beständige Gegenwart dieses Einen eben so eindringlich, wie das
Zusammenstimmen der Kräfte zum Zweck« (Drei Bücher d. M. § 233). Men naar han
fører Alt tilbage til et eneste Princip, saa véd han meget vel, at vi ikke
kunne konstruere Verdenssammenhængen og Elementerne ud fra det: vi mangle
stedse den anden Præmis! (Drei Bücher d. M. § 93). Kun i praktisk Overbevisning
kan da den Tanke holdes fast, at alle Væsener og alle Begivenheder have deres
sidste Grund i det af Verdensgrunden hævdede højeste Formaal. Hvilket dette
Formaal er, vide vi ikke, og ligesaa lidt, hvorledes Kræfternes Modstrid i
Naturen og det Ondes Realitet i den sædelige Verden ere forenelige med
Verdensplanens Gyldighed. \textit{(Mikrokosmus.} Buch IX, Kap. 5. --- Drei
Bücher d. M. § 233). Men i en etisk Ide fandt Lotze baade ved Begyndelsen og
ved Afslutningen af sin Løbebane den sidste Afslutning for sin Tænkning:
Metafysik gaar for ham tilbage til Etik; i det, som \textit{skal} være, ligger
den sidste Forklaring for det, som \textit{er. ---}

Hvis man ved Panteisme forstaar den Lære, der hævder et indre Forhold mellem
Gud og Verden, hvad enten Guddommen opfattes som et personligt Væsen eller
ikke, saa er Lotzes religions-filosofiske Standpunkt at betegne som etisk
Panteisme. Selv følte han sig nær beslægtet med Fichte; Slægtskabet med Spinoza
vilde han derimod ikke ret vide af. I Brugen af teologiske Ord og Vendinger gik
han (især i Mikrokosmus og i de efter hans Død udgivne Grundtræk af
Religionsfilosofi og Moralfilosofi) videre end tilbørligt og berettiget. Ved
hans Sprogbrug (ti Mere er det ikke) kom det til at se ud, som om han stod den
populære religiøse Forestillingskreds nærmere end han i Virkeligheden gjorde.
Dog betoner han lejlighedsvis (især Mikrokosmus. Buch VIII, Kap. 4) Forskellen
mellem den frie Religiøsitet og de positive Religioner. %
% --- p. 522 ---
\opage{522}

\begin{center}\textit{γ. Spiritualistisk Psykologi.}\end{center}

De Elementer, hvis Vexelvirkning frembringer de Fænomener, Erfaringen
frembyder, behøve --- som vi have set --- ifølge Lotze ikke at være absolut
ensartede. Der maa kun være en saadan Ensartethed tilstede, som selve
Vexelvirkningen forudsætter. Ved denne meget ubestemte Regel gør Lotze sig det
muligt at fastholde den populære Antagelse om en Vexelvirkning af Sjæl og
Legeme. De Elementer, Mekanismens Begreb forudsætter, ere for ham dels
psykiske, dels fysiske. Han bestrider Antagelsen af en kontinuerlig Række af
fysiske Fænomener. Paa visse Punkter afbrydes den fysiske Bevægelse for at
»absorberes«, det vil sige, gaa over til psykiske Tilstande. Det er vel
tvivlsomt, om Lotze, da han (i Aaret 1851) første Gang (Allg. Physiol. § 424)
udtalte denne Paastand, har havt en levende Forestilling om Betydningen af den
nysopdagede Lov om Energiens Bestaaen, skønt han heller ikke senere i den fandt
nogen Grund til at fravige sit en Gang grebne Standpunkt. Der var for ham
tvingende Grunde, som nødvendiggjorde Antagelsen af en særegen Sjælesubstans.
Der stillede sig for ham kun to Muligheder: enten at udlede de psykiske
Fænomener af en Sjæl som et for dem ejendommeligt Princip eller at forklare dem
ved de fysiske Kræfters Samvirken. (Se \textit{Selbstanzeige der medicinischen
Psychologie.} Kleine Schriften. III. p. 4). Da nu det Sidste er umuligt,
eftersom fysiske Kræfters Vexelvirkning ikke kan forklare den Enhed, hvoraf
selv den simpleste af Sjælelivets Ytringer er præget, saa bliver der for Lotze
kun den første Udvej tilbage. En stor Del af Lotzes psykologiske Undersøgelser
gaa da ud paa at udfinde den Maade, paa hvilken de psykiske Elementer staa i
Vexelvirkning med de fysiske Elementer. Hans Psykologi minder her om
Descartes'. Der gives efter hans Opfattelse nogle Omraader, hvor de fysiske
Elementer paavirke Sjælen; derefter kunne de sjælelige Processer en Tidlang
forløbe udelukkende efter deres egne Love, og de kunne da igen udløse en
mekanisk Kraft til Frembringelse af nye fysiske Forandringer. (Kleine
Schriften. III. p. 7). Uagtet Nervefysiologiens og Psykologiens nuværende %
% --- p. 523 ---
\opage{523}
Tilstand synes at maatte gøre en Sondring mellem disse forskellige Stadier
haabløs, skrider Lotze dog trøstig (især i \textit{Medicin. Psychol. og
Mikrokosm.)} til en nærmere Undersøgelse. Den Betydning, som de materielle
Processer have for Sjælelivet, kan --- ifølge den Opfattelse, Lotze har gjort
gældende af Vexelvirkningens Natur --- ikke bestaa i, at der fra dem tilføres
Sjælen færdige Fornemmelser eller Forestillinger; de kunne kun give Signaler,
som Sjælen maa oversætte paa sit eget Sprog, ligesom omvendt Sjælens indre
Tilstande ikke som saadanne overføres til de materielle Organer, men kun blive
Anledning til, at disse træde i Funktion. De materielle Organer arbejde i de
højere Aandsvirksomheders Tjeneste ved at overlevere og forberede det
Materiale, paa hvilket Sjælen skal øve sine Kræfter, og deri bestaar hele deres
Betydning. Selve de højere Aandsvirksomheder (Erindring, Tænkning, æstetisk og
moralsk Følelse) foregaa, naar Materialet er givet, i Sjælens Indre, og det er
aldeles unødvendigt at antage særegne materielle Processer, der skulde svare
til dem. Lotze vilde i hvert Tilfælde snarere fortolke saadanne materielle
Processer som Virkninger af de sjælelige Tilstande --- fremkaldte ved, at disse
Svingninger forplantede sig fra Sjæl til Hjerne --- end som Aarsager til disse.
Saaledes især, naar vi se en Følelsestilstand holde sig længe efter, at den
oprindelige Anledning er falden bort.

Paa de specielle psykologiske Teorier hos Lotze kunne vi her ikke gaa ind.
Fremhæves maa især hans geniale Lære om Lokaltegn, hvorved Rumsopfattelsens
Udvikling bliver forstaaeligere\textsuperscript{121}), og hans Hævden af
Følelsens Betydning som fundamentalt Element i Sjælelivet overfor den Hegelske
og Herbartske Psykologi. Paa en Mængde enkelte Punkter kaster Lotze ved sin
fine, aandfulde Opfattelse og sin ofte lykkelige Udtryksmaade nyt Lys. Mod god
psykologisk Metode strider det derimod, naar Lotze i Stedet for at søge de
Love, der forbinde de sjælelige Fænomener indbyrdes, opfatter Forholdet ved de
forskellige sjælelige Fænomeners Opstaaen saaledes, at Forestillingerne virke
tilbage paa »selve Sjælen« og ud af den udløse Følelserne, der igen maa
paavirke »selve Sjælen« for af den at udløse Villiesytringerne. Ligeledes naar
han lader en sær%
% --- p. 524 ---
\opage{524}
egen Tænkeevne optræde i Modsætning til Forestillingsassociationen. Her har
hans spiritualistiske Metafysik paatvunget Psykologien en altfor kunstig og
skolastisk Forklaringsmaade. Det hævner sig her, at Lotze ikke anerkender
Psykologien som selvstændig Erfaringsvidenskab, men betragter den som anvendt
Metafysik. \textit{(Logik und Encyclopädie.} § 93).

Naar Lotzes Psykologi paa væsentlige Punkter minder om Descartes', saa maa det
dog vel fastholdes, at denne Lighed kun er tilstede, saalænge Lotze fastholder
den populære Talebrug, han i sin Psykologi ikke mindre end i sin
Religionsfilosofi anvender i stor Udstrækning. Sjælen en Ting for sig og
Legemet en Ting (eller Gruppe af Ting) for sig --- det kunde efter mange
Udtalelser af Lotze synes at være hans Lære, ligesom det var Descartes' og
endnu er den populære Metafysiks. Men det er kun en foreløbig Opfattelse for
Lotze. Fra Descartes adskiller han sig først og fremmest ved sin Lære om
Udstrækningens Subjektivitet og ved sin Analyse af Materiens Begreb, der fører
ham til den Antagelse, at Materien kun er en fænomenal Form for psykiske
Elementers Vexelvirkning. Derved ophæves den dualistiske Opfattelse af
Forholdet mellem Sjæl og Legeme. Den kvalitative Differens mellem de
vexelvirkende Elementer reduceres derved betydeligt: i Stedet for en
Vexelvirken af psykiske og fysiske Elementer faas indbyrdes Vexelvirken af
psykiske Elementer. Og denne Reduktion foretager Lotze --- hvad der vel maa
mærkes --- ikke, fordi han finder nogen Vanskelighed ved den populære eller den
kartesianske Metafysik. Problemet om Forholdet mellem Sjæl og Legeme existerer
egentligt slet ikke for ham; Vexelvirkningen mellem dem frembyder for ham ikke
nogen større Gaade end det Vexelvirkningsforhold, der finder Sted mellem to
hinanden stødende Legemer\textsuperscript{122}). Det er almindelige filosofiske
Betragtninger, der her lede ham, Betragtninger, der vilde beholde deres
Gyldighed, selv om han, i Stedet for at lægge den kartesianske Dualisme til
Grund i sin Psykologi, havde bygget paa Spinozas Identitetshypotese. Den
metafysiske Idealisme er uafhængig af enhver mere speciel psykologisk Hypotese.
Da Lotze nu begyndte som Kartesianer, føre hine Betragtninger %
% --- p. 525 ---
\opage{525}
ham til en Opfattelse, der kan betegnes som en \textit{monistisk
Spiritualisme,} idet den betragter Sjæl og Legeme som to
\textit{forskellige}\textit{Substanser,} der dog ere af \textit{ensartet}
Beskaffenhed. Herbarts Teori og en Teori, der optraadte indenfor den Wolffiske
Skole (se Note 1 til dette Bind), danne Forgængere for denne ejendommelige Teori.

Paa et andet og for hele hans Filosofi mere væsentligt Punkt skiller Lotze sig
ogsaa fra Descartes (og her tillige fra Leibniz og Herbart). Som vi have set,
betragter han de enkelte Elementer i Tilværelsen som Momenter eller Ytringer
(Aktioner) af Ursubstansen, der er det ene i Sandhed Virkelige. Sjælen kan da
ligesaa lidt som det fysiske Atom være Substans i Ordets strengere Betydning.
De endelige Tings Selvstændighed er kun tilsyneladende. »Vor monistiske
Opfattelse« --- siger Lotze --- »har \textit{helt og hensynsløst} lagt Verdens
Ordning og enhver enkelt Tings Tilværelse og Virkeevne i Hænderne paa det ene
uendelige Væsen, paa hvilket alle Vexelvirkningers Mulighed ene beroede«. (Drei
Bücher d. M. § 245 kfr. § 72---73). Ordene Substans og Væsen tager han da, naar
han bruger dem om Sjælen, kun i Betydning af det, der kan virke og lide (ib. §
243), --- og han indrømmer tilsidst (ib. § 307), at det dog maaske var bedst at
undlade at bruge disse Ord, der let kunne lede vild. Allerede i sin Afhandling
\textit{Seele}\textit{und Seelenleben} (i Wagners »Physiologisches
Handwörterbuch«. 1846. »Kleine Schriften« II p. 198) udtalte Lotze, at enhver
Sjæl kun besidder den Realitet, der tilkommer den paa Grund af dens Betydning i
Verdens Helhed, og at dens Udødelighed ikke beror paa dens Natur, men paa dens
Plads indenfor den etiske Verdensorden. I »Mikrokosmus« og i sit Hovedværk
(Drei Bücher der Metaphysik) vender Lotze tilbage til den samme Tanke. Ligesom
det ikke er Interessen for Sjælens Udødelighed, der har ført ham til hans
spiritualistiske Psykologi, saaledes ser han klart, at dette Spørgsmaal slet
ikke kan besvares ad videnskabelig Vej. »Ingen anden Grundsætning staar her til
vor Raadighed end den almindelige Overbevisning: bestaa vil Alt, hvad der er
blevet til, naar dets Bestaaen hører med til Verdensplanen, og saalænge den
hører med til denne; forgaa %
% --- p. 526 ---
\opage{526}
vil Alt, hvis Virkelighed kun havde sin berettigede Plads i en forbigaaende
Fase af Verdensløbet. At denne Grundsætning ikke i menneskelige Hænder kan
anvendes i det Enkelte, behøver neppe at siges; vi kende visseligt ikke de
Fortjenester, der kunne give det ene Væsen Ret til evig Bestaaen, saalidt som
de Mangler, der kunne frakende et andet denne Ret«. (Drei Bücher d. M. § 245).

Naar man lægger Mærke til Overensstemmelsen mellem Lotzes tidligere (1846 og
1864) og senere (1879) Udtalelser om dette Punkt, saa vil man neppe føres til
med nogle nyere Kritikere af Lotze\textsuperscript{123}) at antage, at Lotze
har foretaget væsentlige Ændringer i sin Lære om Sjælens Substantialitet og
derved nærmet sig mere til Spinozismen end før. Men til Fordel for Klarheden
havde det vistnok været, om den udmærkede Tænker --- hvis Sind ganske sikkert
ikke tillod nogen uværdig Akkommodation --- allerede i sine tidligere
Fremstillinger havde fremhævet sin Afvigelse fra de gængse Antagelser ligesaa
bestemt som i de seneste Udtalelser.

\begin{center}2. GUSTAV THEODOR FECHNER.\end{center}
% PRINTED AS IS: „2.“ where the series a.-c. of section 3 calls for „b.“ (p. 526)

Fechner og Lotze staa som to Dioskurer i den tyske Filosofi i Aarhundredets
sidste Del. De ere beslægtede Aander ved deres ideale Sind, deres store
naturvidenskabelige Dannelse, deres poetiske Sans og deres Trang til Enhed i
Verdensopfattelsen. De forfølge beslægtede Formaal, men tildels ad forskellige
Veje. Derfor er en Sammenligning mellem dem interessant, især da der paa
afgørende Punkter maa vælges mellem dem. Valget vil --- da Grundsætningen,
idealistisk Verdensopfattelse paa realistisk Grundlag, er den samme --- bero
paa, hvem af dem der mest lader det realistiske Grundlag komme til dets fulde
Ret. Det er en Maalestok, Lotze har anerkendt i den Ytring, at Idealismen kun
kan hævdes ved Undersøgelser, der føres i Realismens Aand. Selv om Lotze
indtager den første Plads, hvad Gennemarbejdelsen af de almindelige filosofiske
Principer angaar, turde Fechner have Forrangen, hvad den konsekvente og %
% --- p. 527 ---
\opage{527}
resolute Fastholden af det naturvidenskabelige Grundlag angaar. Dette er
saameget des mærkeligere, som Fechners Fantasi fra først af faar langt større
Indflydelse paa hans Tankegang end det hos hans kritiske, betænksomme Fælle er
muligt. Fechner er --- ligesom Kepler, om hvem han i høj Grad minder --- et
interessant Exempel paa, hvorledes man gennem dristige og fantastiske
Spekulationer kan føres til positive og exakte Resultater, naar man blot
fastholder en bestemt Grundtanke og formaar at udskille den af det mystiske
Svøb. Ligesom Kepler fra mystiske Tankespekulationer successivt førtes til
Opdagelsen af sine berømte Love, der tilfredsstillede hans Trang til at finde
bestemte matematiske Forhold udtrykte i Universet, saaledes førtes Fechner
gennem dristige Analogier til Overbevisningen om et bestemt kvantitativt
Forhold mellem det Aandelige og det Materielle og kom ved den nærmere
Udarbejdelse af denne Tanke til at grundlægge Psykofysiken eller den
experimentale Psykologi.

Fechner er født i Lausitz 1801 og dyrkede i sin Ungdom medicinske og fysiske
Studier. 1835 blev han Professor i Fysik i Leipzig og gjorde sig bekendt ved
dygtige Arbejder i sit Fag. Religionsfilosofen C. H. Weisse hørte til hans
nærmeste Venner og fik stor Indflydelse paa hans --- ligesom paa Lotzes ---
religiøse Opfattelse. Det var dog tvivlsomt, om Fechner var kommen til at
spille saa stor en Rolle i Filosofiens Historie, som han har gjort, dersom han
ikke ved (i Vinteren 1839---40) at studere subjektive Lys- og Farvefænomener
havde ødelagt sine Øjne, saa at han efter lang Tids Lidelse maatte opgive sit
fysiske Professorat. I filosofisk Tænkning, i fri Fantaseren, i Fordybelse i
Sindets indre Verden søgte han Erstatning for den lyse og farverige Verden,
hans svækkede Syn ikke længere tillod ham at dvæle saa vedholdende i som før,
men som han dog ikke glemte. Tvertimod blev det hans Livs Opgave at
sammenarbejde de to Verdener og finde Loven for deres Sammenhæng. I Modsætning
til Romantikens Filosofi vilde han gaa nedefra opad, ikke oppefra ned. Og hvor
den erfaringsmæssige Forsken slipper op, lader han hellere Fantasien tumle sig
i de dristigste Analogier, end han lader den abstrakte Tæn%
% --- p. 528 ---
\opage{528}
ken spinde sine spinkle Traade. Sit Hang til Paradoxier gav han Luft i en Række
humoristiske Skrifter (under det pseudonyme Mærke: Dr. Mises). Hans øvrige
Skrifter dele sig naturligt i to Grupper. Den ene af disse --- indenfor hvilken
især \textit{Zendavesta} (1851), \textit{Ueber die Seelenfrage} (1861) og
\textit{Die drei Motive des Glaubens} (1863) ere at nævne --- udvikler hans
poetisk-spekulative Verdensanskuelse; til den anden høre hans mest epokegørende
Arbejder paa Naturfilosofiens og Psykologiens Omraader: \textit{Ueber die
physikalische und die philosophische Atomenlehre} (1855), \textit{Elemente der
Psychophysik} (1860) og \textit{Vorschule der Aesthetik} (1876).

Lige til sin høje Alderdom var Fechner aandeligt virksom og producerende. Endnu
i sit 86. Aar gav han et interessant Indlæg i Diskussionen om hans psykofysiske
Teori \textit{(Ueber die psychischen Massprincipien und das Weber'sche Gesetz.}
Wundts Philosophische Studien. IV). Kort efter døde han (1887)\textsuperscript{124}).

\begin{center}\textit{α. Poetisk spekulativ Verdensanskuelse.}\end{center}

Fechner bekæmper paa Verdensanskuelsens Omraade enhver Opfattelse, der antager
Sjælelivets rige og lyse Verden som baaret af eller Produkt af dunkle,
anskuelige Ting eller Væsener. Materialismens »Materie«, Spiritualismens
»Sjælesubstans« og Kants »Ding an sich« ere atter og atter Genstande for hans
Polemik. Ikke mindre bekæmper han dem, der skille Gud og Verden, Aanden og
Naturen ad. Han forkaster i lige Grad Troen paa en naturløs Gud og paa en
aandløs Materie, --- Ortodoxien ligesaa vel som Materialismen. Han bebrejder
den sædvanlige Verdensopfattelse, at den dualistisk skiller det Uendelige og
det Endelige ad: »Man stiller det Uendelige overfor, ovenover, hinsides,
udenfor, bagved det Endelige, anbringer maaske endog en uhyre Skillevæg mellem
dem, som om de slet ikke komme hinanden ved. Og dog ligge de slet ikke udenfor
hinanden. Tvertimod er det Endelige det Uendeliges Indhold. . . . Det Uendelige
er derfor heller ikke ufatteligt; det kan fattes ved utallige Hanke indenfor
det Endelige; kun omfattes kan det ikke«. (Ueber die Seelenfrage. p. 111). %
% --- p. 529 ---
\opage{529}

I nøje Sammenhæng med de anførte Fordomme staar efter Fechner den Mening, at
kun Mennesker og Dyr ere besjælede, ellers Intet, idetmindste paa vor Jord. Man
tror, at Erfaringen forbyder at udstrække Besjælingen videre. Men kun om min
egen Sjæl har jeg direkte Erfaring; til andre Sjæles Existens maa jeg slutte
mig ad Analogiens Vej. Og hvad kan da hindre mig i at strække Analogien videre,
fra Mennesker og Dyr til Planterne og til de himmelske Kloder, naar der viser
sig gode Grunde dertil? At Planterne ingen Nerver have, er intet Modbevis; ti
de laveste Dyrearter have heller ingen Nerver. Siger man, at Planten intet
sluttet Individ er, da er at svare, at Individualitet stedse er relativ, da jo
intet levende Væsen staar absolut sondret overfor Omverdenen. Der er en saa
kontinuerlig Overgang mellem Dyre- og Planteliv, at det er uberettiget at
fastslaa en Modsætning mellem dem som mellem Besjælet og Sjælløst. Plantens
Sjæleliv kunde jo staa ligesaa langt under Dyrets, som dettes under Menneskets.
--- Og hvorfor skulde Himmellegemerne ikke være besjælede? Mennesker og Dyr
høre jo med til Jorden, og Jordsjælen kunde forholde sig til de enkelte
Menneske- og Dyresjæle, som Jordlegemet til de enkelte Menneske- og
Dyrelegemer. Det er kun en kunstig Abstraktion, naar man stiller Menneske- og
Dyrelivet som besjælet i Modsætning til hele Klodens Liv. Det Levende fremgaar
ikke af det Livløse, men det Livløse er Produkt af organisk Stofskifte og af
Organismers Undergang; denne Opfattelse lader sig ifølge Fechner snarere
forsvare end den modsatte. I sit Skrift \textit{Ideen zur Schöpfungs- und
Entwickelungsgeschichte der Organismen} (1873) udvikler han især den Tanke, at
Variation og Differentiation er foregaaet med større Hurtighed i tidligere
Jordperioder end i den nuværende. --- De lavere Sjæle maa antages at forholde
sig til de højere, som Forestillinger og Motiver forholde sig til den enkelte
Sjæl. Tilsidst ere alle Sjæle Dele af den højeste, alt omfattende Sjæl, hvis
Liv og hvis Virkelighed kundgør sig i Aarsagsloven, Principet for alle enkelte
Naturlove og dermed for al Sammenhæng og Orden i Tilværelsen.

Ligesom for Lotze er altsaa den lovmæssige Sammenhæng %
% --- p. 530 ---
\opage{530}
i Verden det Grundlag, paa hvilket den religionsfilosofiske Anskuelse bygges.
Hvad der for Mange gør Gudsideen umulig eller overflødig, gør den for Fechner
og Lotze netop nødvendig. I hin universelle Lov finde de den højeste Enhed, det
Evige og Uforanderlige, der omfatter Alt, udtrykt. Lovens Begreb er det sidste
i al vor Erkendelse; paa det maa derfor efter Fechner vor højeste Ide bygges.
(Se især hans Afhandling \textit{Ueber das Causalgesetz} i »Berichte der sächs.
Societät der Wissenschaften« 1849. Kfr. Seelenfrage p. 205 f. Atomenlehre. 2.
Aufl. p. 125).

Verdensbegrebet gaar da for Fechner, ligesom for Lotze, tilbage i Gudsbegrebet.
Den videste Opfattelse af Gudsbegrebet anser Fechner for den mest begrundede;
den snevrere er i hans Øjne en Abstraktion. Og da Verdens Liv er Guds Liv, kan
dette ikke være afsluttet i sig selv, men udvikler og udfolder sig med og ved
Verdensudviklingen. Guds Fuldkommenhed bestaar ikke i færdig Afsluttethed, men
i ubegrænset Fremadskriden. Den guddommelige Verdensplan er ikke færdig en Gang
for alle, men udformes efter Verdensbegivenhederne, der tillige ere
guddommelige Oplevelser. --- I denne Lære (i hvilken Fechner følger sin Ven
Weisse, medens Lotze ikke mente at kunne tillægge Tiden Gyldighed for
Guddommen) ser Fechner en Gennemførelse af Kristendommens Udsagn om, at Gud er
Aand og skal dyrkes i Aand og Sandhed, og at i ham ere, leve og røres vi, ---
Udsagn, der for det Meste staa hen som blotte Ord! Han indrømmer, at hans Lære
ganske vist ikke er kristelig, naar man som væsentlig Bestanddel af
Kristendommen regner »Troen paa Æblebiddet i Paradis og dets mystiske Følger,
paa de ikke Udvalgtes evige Fordømmelse, paa Underne mod Naturens Love, paa
Guds Fjernhed fra Verden og paa alt det Uopbyggelige, hvoraf Teologerne opbygge
Kristendommen« (Seelenfrage p. 194). Fechners dristige Teologi gør ham det
lettere at forstaa, hvorledes Mennesket kan være Guds Medarbejder, og hvorledes
det Onde kan være fremmed for Guds Væsen og dog nødvendigt Led i
Verdensudviklingen. %
% --- p. 531 ---
\opage{531}

\begin{center}\textit{β. Psykofysik.}\end{center}

Tidligt var Fechner i Overensstemmelse med hele sin Grundanskuelse kommen til
den Overbevisning, at Forskellen mellem det Aandelige og det Materielle ikke
kunde være en Forskel mellem to Væsener, af hvilke det ene var immaterielt, det
andet aandløst. Den materielle Verden er Guddommens ydre, den aandelige Verden
er Guddommens indre Side. Forskellen er fænomenal, beror paa en Forskel i
Tilskuerens Standpunkt, ikke paa en Forskel i Substans. (Zendavesta II. p.
341). Eller, som han senere udtrykte sig (Elemente der Psychophysik I. p. 2),
Forskellen mellem det Aandelige og det Materielle er at ligne med Forskellem %
% PRINTED AS IS: „Forskellem“ (p. 531)
mellem den konkave og den konvexe Side af en og samme Cirkel: den, der staar
indenfor Cirkelen, ser kun den konkave Side, den, der staar udenfor Cirkelen,
ser kun den konvexe Side --- og den, der kan skifte Standpunkt, tror maaske, at
det er to forskellige Ting, han har for sig. Det er Spinozas
Identitetshypotese, Fechner her kommer til, ledet af Fysikerens Interesse for
at hævde de fysiske Processers Kontinuitet, og ledet af Religionsfilosofens
Interesse for at tænke Guddommen, det altomfattende aandelige Princip,
saaledes, at den paa ethvert Punkt udtrykker sig i materielle Fænomener, saa at
disse ikke løsnes fra den og falde udenfor den. Fechner faar, saa at sige, en
Teofysik, før han udvikler en Psykofysik.

En Bekræftelse af denne Grundopfattelse fandt han i Loven om Energiens
Bestaaen, og han er den Første, som benytter den nys opdagede Lov paa denne
Maade. (Elemente der Psychophysik. I. p. 21---45). Han indrømmer, at der ikke
er ført noget Bevis for, at denne Lov gælder for de materielle Processer, der
ere knyttede til de aandelige Virksomheder; men han mener, at alle Erfaringer
pege i denne Retning, og at vi altsaa, indtil et Modbevis foreligger, maa
antage Loven for gyldig. Erfaringen godtgør, at de materielle Processer, til
hvilke Bevidsthedslivet er knyttet, staa i Vexelvirkning med de andre
materielle Processer i og udenfor Legemet, saa at den samlede fysiske Energi,
vi raade over, snart kan forbruges overvejende ved hine, snart ved disse. Der
tæres paa den fysiske Energi, baade naar vi tænke, %
% --- p. 532 ---
\opage{532}
og naar vi hugge Brænde, --- og derfor kunne vi ikke gøre begge Dele paa én
Gang ligesaa godt, som hver Del for sig.

Fra Spinoza adskiller han sig i denne Lære kun ved at lægge stærk Vægt paa at
udfinde et matematisk Funktionsforhold\textsuperscript{125}) mellem
Tilværelsens to Sider. Han gik først ud fra, at der maatte være et ligefrem
proportionalt Forhold mellem dem. Senere (han har noteret Dato og Situation:
den 22. Okt. 1850 om Morgenen i Sengen) kom han paa den Tanke, at det Aandelige
ikke ligefrem steg og faldt med det Materielle, men at hints Ændringer svarede
til dettes forholdsmæssige Ændringer, saa at Ændringen (d\textit{γ}) af en
aandelig Tilstands Intensitet er at sætte lig med Forholdet mellem Ændringen
(d\textit{β}) i den tilsvarende materielle Tilstands Energi og den Energi
(\textit{β}), som i Forvejen var tilstede (altsaa d\textit{γ} = d\textit{β/β}).
Hermed var han fra de ubestemte Spekulationers Region naat ned til en
Antagelse, der kunde kontrolleres experimentalt. Han fandt den strax bekræftet
ved, at Lysfornemmelsens Styrke ikke voxer ligesaa hurtigt som den fysiske
Lysstyrke. Egne Forsøg og Efterforsken i Literaturen gav ham et Materiale, paa
Grundlag af hvilket han opstillede den almindelige Lov, at en Fornemmelses
Tilvæxt ikke ligefrem svarer til det fysiske Indtryks Tilvæxt, men til
Forholdet mellem denne Tilvæxt og det hele tilstedeværende Indtryk, saa at
Fornemmelsen ved en og samme Indtryksforøgelse voxer saameget desto mindre, jo
stærkere Indtrykket i Forvejen er. Denne Lov kalder Fechner den Weberske Lov,
til Minde om sin Lærer, Fysiologen E. H. Weber, fra hvem nogle af de
betydningsfuldeste Forsøg, hvorpaa den byggedes, skrev sig. --- Om Lovens
Gyldighed, Begrænsning og Forstaaelse skal her ikke nærmere tales; derom maa
henvises til nyere psykologiske Værker. Her skal kun bemærkes, at ved det
Skridt, Fechner saaledes gjorde, blev exakt Experimenteren mulig paa det
psykologiske Omraade: han gjorde (for at bruge et Udtryk af Galilei) Noget
maaleligt, der hidtil ikke havde været det. Gennem sit Forhold til det ydre
Indtryk kan Fornemmelsen, der i og for sig, som rent subjektivt Element, ikke
kan maales, bestemmes kvantitativt. Dermed var den Vej an%
% --- p. 533 ---
\opage{533}
vist, \textit{den experimentale Psykologi} maatte gaa: den havde at søge
Tilknytningspunkter mellem psykiske Processer og saadanne Processer, der
direkte kunne tælles, vejes eller maales.

Hvad Fechner først havde søgt efter, var ganske vist ikke det, den nye
Videnskab kom til at beskæftige sig med. Han havde villet finde Forholdet
mellem de sjælelige Virksomheder og de tilsvarende Hjerneprocesser. Men dette
Forhold er os utilgængeligt. Hvad vi umiddelbart opfatte, er fra den subjektive
Side den sjælelige Tilstand, fra den objektive Side det fysiske Indtryk. For
Forholdet mellem disse gælder den af Fechner fundne Lov (indenfor visse
Grænser), og det er sandsynligst, at Hjerneproces og sjælelig Proces ere
direkte proportionale; det stemmer ogsaa bedst med den af Fechner hyldede
Identitetshypotese. Psykofysiken er forsaavidt bleven noget Andet, end Fechner
selv havde tænkt sig. Men dette forringer ikke Betydningen af hans geniale Værk
\textit{Elemente der Psychophysik:} mellem Aandsvidenskab og Naturvidenskab er
der --- i Konsekvens af Arbejdsdelingens store Lov --- skudt en ny Disciplin
ind. Og det var i Fechners frie, paa Muligheder og Ideer rige Aand, at denne
Knopskydning fuldbyrdedes. Den Forsker, som næst efter Fechner har virket mest
paa den nye Videnskabs Omraade, udtalte ved hans Grav: »Psykofysiken, som han
grundlagde, var kun den første Erobring paa et Omraade, hvis fortsatte
Tilegnelse ikke mere kunde frembyde betydelige Vanskeligheder, efterat denne
Begyndelse én Gang var gjort«. (\textsc{Wilhelm Wundt}: \textit{Zur Erinnerung
an Fechner.} Philos. Studien. IV. p. 477).

Til Forsvar for den psykofysiske Lov, han havde opstillet, har Fechner
forfattet en Række Skrifter, der ere Mønstre paa sandhedskærlig og omhyggelig
Drøftelse og ridderlig Polemik; tillige præges de af Forfatterens Friskhed og
Humor. ---

Beslægtet med Fechners Bestræbelser for Grundlæggelsen af en experimental
Psykologi er hans Bestræbelse for en paa Erfaring grundet Æstetik. Hans
\textit{Vorschule der Aesthetik} (1876) er et epokegørende Arbejde, hvis
specielle Indhold vi dog her ikke kunne gaa nærmere ind paa. %
% --- p. 534 ---
\opage{534}

\begin{center}\textit{γ. Naturfilosofi.}\end{center}

Sin Modsætning til Romantikens Filosofi har Fechner især udtrykt i sin
»Atomenlehre«. Spekulationen, siger han, stiger ned i Naturen, som Bjørnen i et
Bistade. Den nyder den sammenhængende Honning, men ænser ikke, at den er bragt
sammen ved en Mængde enkelte smaa Væseners Virksomhed. Naturvidenskaben
paaviser --- baade i Fysiken og i Kemien --- at de materielle Processer kun
kunne forstaas ved at betragtes som Vexelvirkningsprocesser af Smaadele, som
for os ere udele lige. Enhver nok saa lille Del af den materielle Verden
opfatter Naturvidenskaben som en lille Verden med Kloder og Systemer af Kloder.
Hvad der for os synes at være en kontinuerlig Masse, er ligesom saa mange
Taagemasser paa Himlen en Gruppe af forskellige Dele. Naturvidenskaben antager
ikke absolute Atomer, men kun, at Deleligheden gaar videre, end Øjet og
Mikroskopet opdager. Kun relative, d. v. s. for os eller ved de os bekendte
Naturkræfter uopløselige Atomer antager den, ledet af Trangen til at henføre
enhver Virkning til et bestemt Sted i Rummet som Udgangspunkt. Den filosofiske
Atomistik gaar den Vej til Ende, som den fysiske Atomistik kun havde Grund til
at vandre et Stykke frem paa. En filosofisk Afslutning paa Naturopfattelsen faa
vi, naar vi antage, at jo mindre man tænker sig Atomerne, des nøjagtigere ville
de Resultater blive, man kan udlede ved at tænke sig Materien bestaaende af
Atomer. Det var jo for at have Udgangspunkter for Virkningerne, at
Naturvidenskaben tænkte sig Atomer. Den Andel, hvormed ethvert saadant
Udgangspunkt bidrager til en Naturlovs Opfyldelse, kalde vi den Kraft, der
udgaar fra dette Punkt. Hvad vi vide om Atomerne, er nu netop kun, at disse
Bidrag udgaa fra dem: derfor betegner den filosofiske Atomistik dem som
Kraftcentra. Der er altsaa ingen Grund til at tillægge Atomerne (de absolute,
filosofiske Atomer) Udstrækning, og Materiens Begreb er ikke længere
materialistisk. Filosofien staar her for Fechner som en Hypotese, i hvilken vi
tænke som absolut, hvad der i Virkeligheden stedse kun er givet relativt.

Grundlaget for hele vor Erkendelse af Tilværelsen er Loven %
% --- p. 535 ---
\opage{535}
for Fænomenernes Sammenhæng; udfra Loven komme vi til at bestemme Kræfterne i
de samvirkende Elementer. Fechner vil ikke med Leibniz, Herbart og Lotze i
Atomerne se psykiske Elementer. Disse ere for ham blot de sidste Punkter, som
de aandelige Væsener naa til, naar de analysere deres Bevidstheds Indhold.

En nærmere Forklaring anser Fechner for overflødig. Atomet er den lavere Grænse
for vor Erkendelse, ligesom den universelle Verdenslov, der vidner om et
altomfattende Væsens Realitet, er den højere Grænse. Mellem disse to
Grænsebegreber ligger al vor Viden. (Seelenfrage. p. 215---216).

\begin{center}c. EDUARD VON HARTMANN.\end{center}

I Aaret 1869 udkom en Bog, der gjorde stor Opsigt og oplevede en for
omfangsrige filosofiske Værker sjælden Udbredelse. Det var Eduard von Hartmanns
\textit{Philosophie des Unbewussten.} Dens Motto: »Speculative Resultate nach
inductivwissenschaftlicher Methode« viser, at Forfatteren tilsigtede en
lignende Begrundelse og Udvikling af filosofiske Ideer som Lotze og Fechner.
Ved sin klare og brede Fremstilling, ved den udførlige Brug af
naturvidenskabelige Exempler, ved paa én Gang at møde med en realistisk
Forgrund og en romantiskmystisk Baggrund og endeligt ved at optage Pessimismen,
der under Indflydelse dels af Schopenhauers nu meget læste Skrifter, dels af
herskende Tidsstrømninger var en udbredt Tankegang, som Led i sit System, kunde
Hartmann finde Læsere i store Kredse. Den Opmærksomhed, han vakte, ses af, at
der i Løbet af tyve Aar udkom ti Udgaver af hans Hovedværk, og at der blot i
Aarene 1870---1875 udkom 58 Skrifter, som handlede om ham!

\textsc{Eduard Von Hartmann} er født i Berlin 1842, var Søn af en General og
betraadte selv den militære Løbebane. I sin Fritid dyrkede han Musik, Maleri og
Filosofi. En Knæskade nødte ham i Aaret 1865 til at tage sin Afsked. Og da han
blev klar over, at hans egentlige Begavelse ikke laa paa Kun%
% --- p. 536 ---
\opage{536}
stens Omraade, hengav han sig --- »bankerot paa Alt, kun ikke paa Tanken« ---
til Filosofien. Allerede tidligere havde han filosoferet og udarbejdet
filosofiske Forsøg. Nu skred han til Udarbejdelsen af »Philosophie des
Unbewussten«, der stadigt maa betragtes som hans Hovedværk, uagtet det er
blevet efterfulgt af omtrent tredive større og mindre Værker. Det betydeligste
af disse er sikkert \textit{Phänomenologie des sittlichen Bewusstseins} (1879),
som idethele er Hartmanns bedste Værk. Til Karakteristik af hans filosofiske
Standpunkt vil jeg her holde mig til to Hovedpunkter, med fuld Bevidsthed om,
at Hartmanns utrættelige Virksomhed for at udvikle sine Ideer i alle Retninger
og belyse dem fra alle Synspunkter derved ikke helt kommer til sin Ret.
Hartmann har klaget over, at jeg har »halveret« ham ved kun at omtale den Del
af hans Virksomhed, der ligger før 1880. Dette laa nu dog i denne Bogs Plan;
tillige har jeg unægteligt den Opfattelse, at de to nævnte Værker ere og blive
Hartmanns betydeligste Bidrag til Filosofien. Dem, der ønske at supplere den
korte Fremstilling af Hartmanns Filosofi, jeg her kan give, maa jeg henvise til
det udførlige og af overvættes Beundring udsprungne Værk af ARTHUR DREWS:
\textit{Hartmanns philosophisches System im Grundriss} (Heidelberg 1902).

\begin{center}\textit{α. Naturfilosofi og Psykologi.}\end{center}

Medens Lotze og Fechner principielt lagde den naturvidenskabelige
Forklaringsmaade til Grund og vilde begrunde deres Filosofi ved at drage
Konsekvenserne af dens Forudsætninger og Resultater, har Hartmanns Filosofi
langt mere Præget af en nyromantisk Reaktion mod Naturvidenskabens Realisme.
Han gaar ud paa at vise, at den naturvidenskabelige Forklaringsmaade ikke er
tilstrækkelig, men at der sammen med de Aarsager, den mekaniske Naturopfattelse
antager, maa virke et aandeligt Princip, som han --- for ikke at gøre det for
menneskelignende --- kalder \textit{det Ubevidste.} Han appellerer til dette
Princip overalt, hvor han mener, at de Aarsager, Erfaringsviden%
% --- p. 537 ---
\opage{537}
skaben kan paavise, ikke slaa til. Det er saaledes, han ad »induktiv« Vej mener
at komme til »spekulative« Resultater.

Hartmann er enig med Fechner i, at Materien maa opfattes som et System af
Atomkræfter. Hvert enkelt Atom kan ingen Masse have, ligesaa lidt som man kan
tale om Størrelsen af en Ener. Holde vi os da til Atomkraften som det sidste
Element, til hvilket Analysen af den materielle Virkelighed fører, saa kan der
--- mener Hartmann --- kun findes bestemt Sammenhæng og Konsekvens i
Atomkraftens Stræben, naar denne opfattes som en Villen med en ubevidst
Forestilling om dens Maal. Vi forstaa kun Kraften ved at tænke os den som en
Villen. Selve Materien er altsaa Forestilling og Villie --- og Forskellen
mellem Materie og Aand falder bort. I den organiske Væxt ytrer sig ligeledes en
ubevidst Villen og Forestillen: der virkeliggøres jo ved Organismens
Tilblivelse et Formaal, et Formaal, til hvilket de enkelte Stoffer og Processer
ere at betragte som Midler, hvis Tilvejebringelse og Sammenføjning kun bliver
forstaaelig ved Formaalet, uagtet dette ikke er Genstand for Bevidsthed. Der er
kun en Gradsforskel mellem organisk Dannelse og Instinkthandling. Instinktet
kan, paastaar Hartmann, ikke forklares ved den materielle Organisation eller
bero paa en en Gang for alle indplantet Nervemekanisme; ti saa kunde det ikke
variere i sine Ytringer, naar Organisationen ikke forandres. Ligesaa lidt beror
det paa bevidst Overlæg. Der bliver da kun tilbage at opfatte det som ubevidst
Villen og Forestillen. I den menneskelige Aand ytrer det Ubevidste sig i
Sanseopfattelsen, idet Anskuelsen af den ydre Genstand bliver til ved ubevidst
Sammenarbejden af Fornemmelserne. Forestillingsassociationen gaar kun for sig
ved, at der uden bevidst Søgen fremdrages Forestillinger, som staa i Sammenhæng
med de nærværende Forestillinger. Følelser og Motiver opstaa paa Grund af
ubevidste Processer og ere derfor ofte uforstaaelige for os selv. Selv hvor
Villien er forbunden med Bevidsthed, maa der en ubevidst Villie til, forat den
kan føre til en villet Bevægelse: Bevidstheden véd jo ikke, hvilket
Nervecentrum i Hjernen der skal paavirkes, forat Bevægelsen kan komme i Stand!
I Historien formaar det Ubevidste at faa de Enkelte til at arbejde i store %
% --- p. 538 ---
\opage{538}
Verdensformaals Tjeneste, medens de tro blot at arbejde for deres egne,
begrænsede Formaal. Hvad er Skæbne eller Forsyn Andet end det Ubevidstes
Herredømme i Menneskenes Handlinger, saalænge deres bevidste Forstand ikke er
moden nok til at gøre Historiens Formaal til sine egne? --- Det Ubevidste er
overalt det, der gør Samvirken mellem individuelle Væsener mulig. Der er
Individer af alle mulige Grader --- ligefra Atomet til Alnaturen, og det er det
ene og samme Ubevidste, som rører sig i dem alle. Flerheden har kun fænomenal,
ikke metafysisk Betydning. Hartmann erklærer sig i Principet enig med den
spekulative Teisme (som den fremtræder hos Schelling, Weisse, Lotze og
Fechner), idet hans »Ubevidste« snarest er at kalde »overbevidst«. Men han
foretrækker det negative Udtryk for at holde antropomorfistiske Forestillinger
borte. Han overser herved, at selv om det Navn, han har valgt, fjerner saadanne
Forestillinger, saa begunstiges disse i høj Grad ved den Maade, paa hvilken han
lader »det Ubevidste« virke.

Det er en uvidenskabelig Fremgangsmaade, Hartmann anvender, naar han overalt,
hvor han mener at finde et Hul i den videnskabelige Forklaring, putter sit
»Ubevidste« ind som et magisk Middel, der virker lige godt overalt. Om det
gælder, hvad Galilei sagde om Appellen til en almægtig Villie: det forklarer
Intet, fordi det lige godt forklarer Alt. Den egentlige Interesse for det
Ubevidstes Begreb er hos Hartmann neppe heller opstaaet af Stræben efter
Forstaaelse af Fænomenerne. Den skyldes snarere en Trang til Hvile fra
Reflexion, Analyse og Kritik og en Tro paa det uvilkaarligt og uforberedt
Opdukkendes Værdi. Saaledes hedder det et Sted hos Hartmann (Philosophie des
Unbewussten. 3. Aufl. p. 369): »Den bevidste Fornuft er kun negerende,
kritiserende, kontrollerende, korrigerende, maalende, sammenlignende,
kombinerende, subordinerende, inducerende og generaliserende. Aldrig er den
skaberisk produktiv, aldrig opfindende; i denne Henseende er Mennesket ganske
afhængigt af det Ubevidste, og naar det mister Evnen til at opfatte det
Ubevidstes Indskydelser, mister det sin Livskilde, uden hvilken det vil
henslæbe sit Liv paa ensformig %
% --- p. 539 ---
\opage{539}
Maade under tør Skematiseren. Derfor er det Ubevidste uundværligt, og ve den
Tidsalder, der voldsomt undertrykker det og i ensidig Overvurdering af det
Bevidst-Fornuftige kun vil lade dette gælde!« Bevidsthed kan ikke være
Formaalet for Udviklingen, men kun være et Middel, som er nødvendigt paa et
vist Trin af Udviklingen. (Phil. d. Unbew. 3. Aufl. p. 739). --- Det er denne
psykologisk-historiske Erfaring om det ubevidste Livs Betydning, der er det
berettigede Grundlag for Hartmanns Tænkning. Men han gør saa det Ubevidste til
et mytologisk Væsen, en Dæmon, der overalt griber ind i den naturlige
Sammenhæng, dirigerer Atomerne, forskyder Hjernemolekulerne, lemper Forholdene.
Hartmanns Filosofi faar en mytologisk og dualistisk Karakter, som ganske vist
er mest fremtrædende i de første Udgaver af »Philosophie des Unbevussten«, %
% PRINTED AS IS: „Unbevussten“ (p. 539)
men er for nøje forbunden med hans Princip til, at den trods alle Korrektiver
helt kan forsvinde. Karakteristisk er især hans Stilling til Darwinismen, idet
han søger at vise, at Kvalitetsvalget ikke forklarer nye Formers Opstaaen, men
kun er et mekanisk Middel, »det Ubevidste« benytter sig af, naar det vil bane
Vej for saadanne Former! Darwin holder sig --- mener Hartmann --- til
Betingelserne, men overser den egentligt producerende Kraft!

Det var da intet Under, at Hartmann ofte blev angreben fra naturvidenskabelig
Side. For nu at vise, at han selv meget godt magtede de naturvidenskabelige
Synsmaader skrev han en anonym Kritik af sig selv: \textit{Das Unbewusste vom
Standpunkt der Physiologie und der Descendenzlehre} (1872). Kritiken var saa
god og rammende, at hans Modstandere roste det anonyme Skrift og brugte det
imod ham. I anden Udgave navngav Hartmann sig og føjede Noter til, i hvilke han
søgte at gendrive de Indvendinger, han selv havde rejst. Det er dog tvivlsomt,
hvorvidt det er lykkedes ham; han har fremmanet Aander, som ikke ville
forsvinde igen. Han indrømmer, at han havde undervurderet de mekaniske Aarsager
og overset mange Mellemled. \emph{Men} hans Fremgangsmaade er vedblivende den,
at han opsøger de Huller, den videnskabelige Forsken stedse lader tilbage, og
udfylder dem ved Hjælp af sin mystiske deus ex machina. Da der stedse er
Mulighed for, at Hullerne kunne udfyldes viden%
% --- p. 540 ---
\opage{540}
skabeligt, er det en mislig Metode at anvende. Lotze anvendte denne Metode i
sin spiritualistiske Psykologi; Hartmann udstrækker den til hele sin Filosofi
uden at agte paa Lotzes klare og energiske Fremhæven af Mekanismens Betydning.
Hans »spekulative Resultater« ere ikke vundne ved »induktiv Metode«. Han
indrømmer da ogsaa (Das Unbewusste etc. 2. Aufl. p. 359 f.) Induktionens
Utilstrækkelighed og tillægger den kun den Betydning at tjene til Bekræftelse
af, hvad der ad anden Vej er udspekuleret.

\begin{center}\textit{β. Pessimisme og Etik.}\end{center}

Hartmann erklærer det for en stor Misforstaaelse, at han selv i sin personlige
Stemning skulde være Pessimist; han lærer Pessimismen som \textit{teoretisk
Anskuelse,} efter at Skuffelser og Lidelser i hans Ungdom aabnede hans Øjne for
denne Opfattelse. » Med Kapitlet om Pessimismen [i »Philosophie des
Unbewussten«] har jeg«, siger han, »for stedse skrevet mig løs fra
Verdenssmerten som saadan og lutret den til en objektiv, rolig \textit{Viden}
om Tilværelsens Elendighed, men derved ogsaa tilbageerobret mig den
uforstyrrede Oprømthed, som den kan findes hos Filosofen, der svæver i den rene
Tankes Æter og udfra den betragter Verden og sin egen Smerte som et fremmed
Undersøgelsesobjekt«. \textit{(Mein Entwickelungsgang.} Gesammelte Studien und
Aufsätze. p. 38). Tillige fremhæver han, at Pessimismen end ikke teoretisk
udgør hele hans Filosofi, og ikke giver denne dens Hovedpræg. \textit{(Die
Stellung des Pessimismus in meinem philosophischen System.} Zur Geschichte und
Begründung des Pessimismus. 2. Aufl. p. 18---28). Han betegner som sin Opgave
at forene Pessimismen, som Resultat af Vurderingen af Verden efter
Lykkesynspunktet, med Udviklingsfilosofien. Han vil altsaa forsone
Schopenhauers Filosofi med Hegels, --- ligesom han allerede i selve sit Begreb
om det Ubevidste har søgt at kombinere Schopenhauers absolute Villie med Hegels
absolute Ide. Paa ejendommelig Maade finder Hartmann her et Tilknytningspunkt
hos Kant. Naar Kant erklærede ethvert Forsøg paa at løse det Ondes Problem for
haabløst, da var det fra hans Stand%
% --- p. 541 ---
\opage{541}
punkt konsekvent, da han gik ud fra det kristelige Gudsbegreb, med hvilket det
Ondes Realitet staar i uovervindelig Modsigelse. Men Problemet maa stilles
saaledes: hvorledes maa Verdens absolute Grund tænkes, naar den uden Modsigelse
med sig selv skal være kommen til at sætte en saadan Verden? Og naar
Spørgsmaalet stilles saaledes, har Kant banet Vej for dets Behandling, idet han
baade har givet Antydninger i pessimistisk og i evolutionistisk Retning.
Schopenhauer forfulgte den ene Antydning, Hegel den anden: nu er Tiden kommen
til at sammenfatte hvad der ligger i begge. \textit{(Kant als Vater des
modernen Pessimismus.} Zur Gesch. u. Begr. des Pess. 2. Aufl. p. 136 f.).
(Smlgn. Note 17 i nærværende Bind).

Da der nu --- saaledes ræsonnerer Hartmann --- for en uhildet Betragtning viser
sig mere Ulykke end Lykke i Verden, kan Verdens Tilblivelse ikke udledes af
Fornuften. Den kan ikke skyldes et rationelt, men maa skyldes et irrationelt
Princip. Forklaringen maa da søges i, at Villieselementet indenfor det
Ubevidste har sondret sig fra Forestillingselementet. Hartmann optager Böhmes
og Schellings mystiske Lære om et dunkelt Element i Guddommen, hvis Løsrivelse
fører til Verdensudviklingen. Paa grundløs Maade har den i sig selv blinde
Villen revet sig ud af Forbindelsen med Forestillingen eller Ideen. Men dette
seende Element i det Ubevidste virker stedse med i Verdensudviklingen og søger
at føre hint blinde og stridige Element til Harmoni og Forsoning. Deraf det
dobbeltsidige Indtryk, vi faa af Verden. Baade Pessimismen og Evolutionismen
have Ret. Der er to Principer, som virke i Verden. Hartmanns Religionsfilosofi
kommer her til at minde om Voltaires, Rousseaus og Stuart Mills, --- kun at
Udvikiingen %
% PRINTED AS IS: „Udvikiingen“ (p. 541)
efter hans Tro er en fremskridende Forløsning af den lidende Gud, idet det
gælder om gennem Ophævelsen af Villien at gøre det ugjort, som skete, da Verden
blev en Virkelighed. Af stor Betydning er det i denne Henseende, at Erkendelsen
af Livets Elendighed vaagner, og derfor mener Hartmann, at der med Schopenhauer
er begyndt en ny Verdensperiode. (Phänomenologie des sittlichen Bewusstseins.
p. 782). Kun langsomt opgaar denne Erkendelse i Menneskeheden. Det sker gennem
en successiv Op%
% --- p. 542 ---
\opage{542}
hævelse af Illusioner. Først venter man Lykken i det nærværende Liv. Naar man
opdager, at de Goder, det kan bringe, kun ere tilsyneladende, sætter man sin
Lid til et andet Liv. Det er Illusionens andet Stadium. Troen paa det Hinsidige
kan ikke holde sig; man faar paany Tillid til det jordiske Liv, men belært af
Erfaring venter man ikke en umiddelbar Nydelse af Lykken, men haaber en
Lyksalighedstilstand for Menneskeheden engang i Fremtiden og anspænder alle
Kræfter for at bane Vejen for den. Det er Illusionens tredie og sidste Stadium,
i hvilket Menneskeheden sandsynligvis vil blive staaende i lange Tider. Men i
selve det Faktum, at en lille Minoritet har gennemskuet Illusionen i alle dens
tre Former, ser Hartmann tilstrækkelig Berettigelse for den pessimistiske
Teori; efter Afstemning kan dens Rigtighed ikke afgøres.

Hartmanns Etik hænger nøje sammen med hans pessimistiske Teori. Der er for ham
en uforsonlig Modstrid mellem Kultur og Lykke. Kulturens Fremskridt ere
forbundne med Tilbagegang i Lykke. Jo mere kompliceret Livets Apparat bliver,
desmere Anledning til Ulyst gives der. Modtageligheden for Smerte bliver
større, og den stigende Reflexion gennemskuer lettere Illusionerne. Kulturen
lader Fornødenhederne voxe hurtigere end Midlerne til deres Tilfredsstillelse.
Man maa da vælge mellem Kulturen og Lykken, mellem Udviklingslæren og
Lyksalighedslæren. Lykken forudsætter Fred og Hvile --- og vil derfor føre
Stagnation og Opløsning med sig. Udviklingen fører stedse videre --- indtil
alle Muligheder ere udtømte\textsuperscript{126}). Ved Udviklingen forøges
Utilfredsheden og ophæves Illusionerne, indtil Livet tilsidst stanser. Da er
Forløsningen naat, og den blinde Villies Løsrivelse er ophævet. Den højeste
Form for Etik er efter Hartmann den, der ikke blot bygger paa Menneskekærlighed
og Udviklingstrang, men gør Medlidenheden med Gud, Trangen til at udløse det
bundne og lidende Verdensprincip, til Moralens Grundlag. Man gør da det
Ubevidstes Formaal til sine egne Formaal.

Hartmanns Filosofi er omvendt Teologi. Han opererer trygt med Tanken om et
absolut Verdensformaal, uden at indlade sig for Alvor paa de Vanskeligheder,
denne Tanke indeholder. I %
% --- p. 543 ---
\opage{543}
hans etiske Værk (Phän. d. sittl. Bew.) findes dog fortræffelige Enkeltheder,
især i Karakteristiken og Kritiken af de forskellige Standpunkter. Og i det
Enkelte er hans Tankegang uafhængig af Pessimismens Metafysik. Han har
overvundet denne i endnu højere Grad end han siger. I Praxis er han
Evolutionist og Optimist. Ti foreløbigt er Bekræftelsen af Villien til Livet
det ene Rigtige, da der kun kan udrettes Noget for Verdensprocessen ved fuld
Hengivelse til Livet og dets Smerter, ikke ved fejg Forsagelse og
Selvopgivelse. Pessimismens Tid kommer egentligt først om lange Tider --- saa
lange Tider, at man undrer sig dobbelt over Hartmanns sikre Viden om, hvad der
vil ske, naar de ere forbi. Ved denne lange Udsættelse mister Pessimismen den
personlige, praktiske Betydning, den havde hos Schopenhauer. Hartmann forholder
sig til Schopenhauer, som de moderne Kristne, for hvem Dommedag er skudt ud i
det Ubestemte, til de første Kristne, hvis hele Liv bestemtes ved Forventningen
om dens snare Komme. Fjerne Muligheder kunne ikke være bestemmende for vor
Livsførelse og kunne ikke rokke ved Livets Værdi. Kun det Liv, vi kende og
leve, kunne vi vurdere, og indenfor selve dette Liv maa vi finde det, vi skulle
kunne anerkende som værdifuldt.

\begin{center}4. KRITICISME OG POSITIVISME.\end{center}

\begin{center}α. FRIEDRICH ALBERT LANGE.\end{center}
% PRINTED AS IS: „α.“ where „a.“ is called for: the next part of section 4 is „b. EUGEN DÜHRING“ (p. 543)

Det var naturligt, at Opløsningen af Romantikens Filosofi, Naturvidenskabens
Opsving og Materialismens Fordring paa at være Videnskabens sidste Ord maatte
bringe den Retning, der i Aarhundredets første Halvdel kun havde været en
Understrøm, frem til større Indflydelse. Allerede den større Opmærksomhed og
Tilslutning, Schopenhauers og Herbarts Ideer nu bleve Genstand for, ere Tegn i
denne Retning. Man trænger til en kritisk Revision af selve Grundlaget for vor
Erkendelse for at kunne stille sig paa rette Maade til det nye Erkendelses%
% --- p. 544 ---
\opage{544}
stof. Hos Lotze og Fechner, tildels ogsaa hos Hartmann, findes vel
Undersøgelser i denne Retning, men det egentlige Erkendelsesproblem tages ikke
op. De tage Realismen som given og søge saa, mere eller mindre konsekvent, at
finde et Grundlag for deres idealistiske Konstruktioner enten i Realismens
Forudsætninger eller i dens Huller. Saa stor Interesse navnligt Lotzes Forsøg i
denne Retning har, saa var det dog mere Tilværelsesproblemet end
Erkendelsesproblemet, han behandlede. En af de Første, der henledte
Opmærksomheden paa Nødvendigheden af at genoptage Kants Arbejde paa
Erkendelsesproblemet, var den berømte Filosofihistoriker \textsc{Eduard Zeller}
i et lille Skrift \textit{Ueber die Bedeutung und Aufgabe der
Erkenntnisstheorie} (1862), der vakte saa meget des større Opsigt, som
Forfatteren var udgaaet fra Hegels Skole. Faa Aar senere henviste
\textsc{Friedrich Albert Lange} i sit glimrende historisk-filosofiske Værk
\textit{Geschichte des Materialismus} (1866) til Erkendelsesteorien som en ikke
blot af den romantiske, men ogsaa af den materialistiske Filosofi overset
Instans, og han søgte paa dens Basis at føre Kampen for en idealistisk
Verdensanskuelse.

Lange er en af de mest tiltalende Skikkelser i den nyere tyske Filosofis
Historie. Ideel Begejstring, skarp Kritik, stor Kundskabsfylde og kraftigt
realistisk Blik have hos ham indgaaet en sjelden Forening. Derhos raader han
over et stort Talent som Forfatter. Desværre blev hans Løbebane kun kort. Han
er født 1828 ved Solingen i Rhinpreussen, men levede sin første Ungdom i
Zürich, hvorhen hans Fader, den bekendte Teolog J. P. Lange, blev kaldet som
teologisk Professor i Aaret 1841. Han fik saaledes Grundlaget for sin aandelige
Udvikling i den frie Schweizerluft, og det sporedes gennem hele hans Liv og
Stræben. Sine Universitetsstudier dyrkede han i Bonn. Som Hovedformaal stod
pædagogisk Virksomhed for ham, og han studerede Filologi som Middel dertil,
uden dog at forsømme naturvidenskabelige og filosofiske Studier. Hans
filosofiske Retning er tidligt blevet grundlagt. I et Brev fra Aaret 1851 (hos
\emph{Ellissen}: \textit{Friedrich Albert Lange.} Eine Lebensbeschreibung.
Leipzig 1891. p. 69---73) udvikler han, at alle Forestillinger, som Mennesket
kan danne sig om Tilværelsen (om Gud og Verden, %
% --- p. 545 ---
\opage{545}
om Godt og Ondt o. s. v.) udvikles efter den menneskelige Aands Love: disse
Love blive da vor Erkendelses sidste Grundlag, udover hvilket vi ikke kunne
komme. Han var her ført ind paa det psykologisk-erkendelsesteoretiske
Standpunkt, som Ludwig Feuerbach gennem Kritik af de religiøse og spekulative
Ideer var naat til, uden at Feuerbach dog synes at have paavirket ham. Hvad der
har paavirket ham af Filosofi i hans Ungdom, var især Hegels »Fænomenologi«,
Herbarts Psykologi og Schleiermachers Skrifter; om denne sidste taler han i et
Brev af 1849 med Begejstring.

Efter nogen Tids Docent- og Skolevirksomhed virkede Lange som Redaktør og
politisk og social Agitator. Han deltog i Oppositionen mod Bismarcks
Indenrigspolitik og i de tyske Arbejderforeningers Bestræbelser for at vinde en
fast Organisation. Skriftet \textit{Die Arbeiterfrage} (1865) skriver sig fra
denne Tid. Det er et af de første Skrifter (næst efter Mills politiske
Økonomi), i hvilke det sociale Spørgsmaal opfattes fra et mere almindeligt
Synspunkt og paa upartisk Maade. Tillige har det Interesse ved, at de Darwinske
Ideer bringes i Forbindelse med Arbejderspørgsmaalet, idet dette ses udspringe
af Kampen for Tilværelsen paa det sociale Omraade. Lange opfatter Sagen
saaledes, at det gælder om at føre en Kamp mod Kampen for Tilværelsen, en Kamp,
som netop udgør Menneskets højere aandelige Bestemmelse. (Da paa denne Tid kun
»Arternes Oprindelse«, ikke »Menneskets Afstamning« var udkommen, kunde Lange
ikke benytte de Synspunkter, som Darwin har givet til at opfatte ogsaa de
etiske Bestræbelser som en Kamp for Tilværelsen). Han saa Arbejderspørgsmaalets
politiske Betydning i, at det øver et stadigt Tryk paa de konservative
Institutioner, hvis Dæmninger ville blive gennembrudte, hvis der ikke i Tide
graves Kanaler. De frisindede Mænd i de højere Klasser have nu Valget imellem
at befæste disse Dæmninger eller at hjælpe med til at grave Kanalerne. Langes
eget Valg var let truffet. Men i Tyskland havde man mest travlt med
Dæmningerne, og Lange udvandrede derfor til Schweiz. Midt i sin journalistiske
og agitatoriske Virksomhed havde han dog faaet Tid til at lægge den sidste
Haand paa sit Hovedværk \textit{Geschichte des Material%
% --- p. 546 ---
\opage{546}
ismus} (som ifølge Ellissen p. 145 udkom i Efteraaret 1865, skønt der paa
Titelbladet staar 1866). Det er ikke et rent historisk Værk. Gennem en aandfuld
Skildring af Materialismens Hovedformer og Naturvidenskabernes Udviklingsgang
søger Lange at belyse Problemet om det Aandeliges og det Ideales Forhold til
det Materielle, og han sætter tillige det teoretiske Spørgsmaal i Forbindelse
med det praktisk-sociale. Skønt naturligvis mange Partier af Værket nu ere
forældede, vil dog Ingen læse det uden Berigelse og Vækkelse. --- Kort efter
udkom et nyt social-filosofisk Skrift af Lange: \textit{Mill's Ansichten über
die soziale Frage} (1866).

I Schweiz arbejdede Lange som Boghandler, Lærer og Redaktør i Winterthur og
udfoldede en forbavsende alsidig Virksomhed. Han tog ivrig Del i Agitationen
for en Revision af Kanton Zürichs Forfatning. Senere blev han Professor i
Zürich og kunde nu igen leve for sine Studier. Man blev dog snart opmærksom paa
ham i hans gamle Fædreland, som han ogsaa længtes tilbage til, og i Aaret 1872
blev han Professor i Marburg. Men da havde han allerede i flere Aar lidt af
Kræft i Underlivet, og hans sidste Aar vare smertefulde. Hans ideale Sind og
store Begejstring for Forskning og Frihed holdt ham oppe; sin aandelige Næring
fandt han især i Schillers filosofiske Digte. Han døde 1875. --- Et ufuldendt,
men interessant Arbejde over Logik \textit{(Logische Studien.} 1877) blev
udgivet kort efter hans Død. Senere (1897) har Ellissen udgivet
\textit{Einleitung und Kommentar zu Schillers Philosophischen Gedichten,} et
Arbejde, Lange havde forberedt under sit Ophold i Zürich, men lagt til Side for
andre Arbejder. Schiller var, efter hvad Lange udtalte i Fortalen til sit
Hovedværk, den Kantianer, han følte sig mest beslægtet med. Schillers Digte
indeholdt Langes Religion.

Som Titelen paa Langes Bog angiver, er hans Problem det, hvorledes man skal
stille sig overfor Materialismen. Hans Standpunkt er her det samme som
Fechners. Han vil overvinde Materialismen netop ved at gennemføre den. Hvad han
priser hos Materialismen, og hvad der for ham udgør dens store historiske
Betydning, det er dens konsekvente Fordring om Paa%
% --- p. 547 ---
\opage{547}
visning af materielle Aarsager til alle Naturfænomener. I Idealismen ser Lange
derimod paa det teoretiske Omraade en stor Fare, da den let fører til at slaa
af paa den strenge mekaniske Forklaring og til at indskyde subjektivt
konstruerede Led i de objektive Aarsagers Række. Ogsaa Nervesystemets og
Hjernens Funktioner maa forklares efter de almindelige fysiske og kemiske Love.
Man maa --- støttende sig til Loven om Stoffets og Kraftens Bestaaen ---
forfølge Paavirkningerne udefra gennem Sanseorganer, indadgaaende Nerver,
Hjerne, Bevægelsesnerver, Muskler, og søge en fuldstændig og ubrudt Sammenhæng
i denne hele Proces. Det er uberettiget paa noget Punkt at antage Sammenhængen
for brudt. Selve Bevidstheden er ikke noget Led i denne Sammenhæng. Den er den
subjektive Tilstand hos det Individ, hos hvem Processen foregaar, en anden Side
af Processen, men ikke selv en Del af den. Her har Materialismen sin Grænse:
netop fordi Bevidsthedstilstandene ikke ere Led eller Dele i de materielle
Processer, ikke kunne forklares efter Loven om Stoffets og Kraftens Bestaaen,
netop derfor falde de udenfor Materialismen, og denne kommer i Modsigelse med
sig selv, naar den mener at have forklaret Alt. Den kan i det Højeste forklare
den objektive, materielle Proces, som er den Maade, paa hvilken det bevidste
Individs subjektive Tilstand fremtræder for den ydre Betragtning. Lange anser
det for tænkeligt, at man engang vil kunne bestemme den Del af de materielle
Processer, der i Tiden falder sammen med en vis Bevidsthedstilstand hos
Individet; men derimod mener han ikke, at man nogensinde vil kunne give en
nøjere Bestemmelse af Forholdet mellem den subjektive Bevidsthedstilstand og
den objektive Nerveproces. Det er en berettiget Formodning, at der bag de to
korresponderende Verdener, Materiens Verden og Bevidsthedens Verden, ligger et
ubekendt Tredie som deres fælles Aarsag.

Ejendommeligt er det at se hvorledes det 17. Aarhundredes store Hypoteser
fornyes efter det naturvidenskabelige Opsving ved Midten af det 19.
Aarhundrede. Som Vogt, Moleschott og Büchner forny Hobbes' Standpunkt, fornyer
Lotze Descartes', Fechner og Lange Spinozas. Det er et Vidnesbyrd %
% --- p. 548 ---
\opage{548}
om, at disse Hypoteser angive de Muligheder, mellem hvilke vor Tænkning her maa
ræffe sit Valg. Ved nærmere Undersøgelse vil man finde, at det ikke er en
simpel Genoptagelse, men en Fornyelse af de ældre Hypoteser, der er sket. Ikke
blot det naturvidenskabelige Grundlag, men ogsaa den filosofiske Udformning er
en anden end hos det 17. Aarhundredes store Dogmatikere. Der er større kritisk
Besindighed tilstede, klarere Bevidsthed om, at det her drejer sig om at finde
de bedste og mest konsekvente Synsmaader i en Diskussion, der neppe vil kunne
afsluttes, ikke om at opstille definitive Systemer. Man mærker, at den kritiske
Filosofi, hvis Gennembrud ligger imellem de to Epoker, ikke har arbejdet forgæves.

Langes Standpunkt, som fra en Side viser tilbage til Spinoza, viser nu netop
fra en anden Side tilbage til Kant. Han bliver ikke staaende ved den dobbelte
Form for Tilværelsen --- som Bevægelse og som Bevidsthed, men gør opmærksom
paa, at hele Materien, Verden udenfor os saavel som vor egen Organisme, Hjerne
og Sanseorganer medindbefattede, kun er til for vor Bevidsthed, er et
Bevidsthedsprodukt. Det er i Kraft af vor Organisation, at vi opfatte Verden,
som vi opfatte den. Det har Kant paavist, og den senere Sansefysiologi har
bekræftet det. Atter her overvindes Materialismen ved at gennemføres
konsekvent. Ti selv om det indrømmes den, at hele vor Verdensopfattelse er et
Produkt af vor materielle Organisation, saa er selve denne materielle
Organisation dog kun et Bevidsthedsobjekt, existerer umiddelbart kun i og for
Bevidstheden. Materie og Organisation ere vore Forestillinger, selv om det er
Forestillinger, vi nødvendigvis danne, og danne efter bestemte Love. Striden
mellem Materie og Aand er saaledes endt til den sidstes Fordel. Den naive Tro
paa Materien maa forsvinde ved nærmere Undersøgelse, og den kan forsvinde, uden
at Naturforskningens Erobringstog derved i mindste Maade hæmmes. --- I den
Maade, Lange her stiller sig overfor Materialismen, har han en Forgænger i
Schopenhauer. Omtrent samtidigt med Lange udtalte den berømte Wienerpatholog
\textit{Rokitansky} en lignende Opfattelse\textsuperscript{127}).

Der er i Langes Fremstilling i »Geschichte des Material%
% --- p. 549 ---
\opage{549}
ismus« en vis Vaklen med Hensyn til, hvor vidt han vil strække
Subjektivismen\textsuperscript{128}). Lejlighedsvis udtaler han, at selve
Modsætningen mellem Tingen i sig og Fænomenet er et Produkt af vor
Organisation, og at vi derfor ikke vide, om denne Modsætning har nogensomhelst
Betydning udenfor vor Erfaring. Andre Steder (især i 2\textit{.} Udg.) taler
han om en til vor Erkendelse svarende Tingenes Orden. I et Brev til
Naturforskeren Dohrn (hos Ellissen p. 258---262) udvikler Lange, at der med
Hensyn til hele dette Spørgsmaal maa skelnes mellem fire forskellige Stadier:
1) den naive Tro paa det sanselige Skin; 2) den Erkendelse, at Sanserne ikke
give os Tingene, som de i sig selv ere, men Luftsvingninger som Lyd,
Ætersvingninger som Farve o. s. v. --- det er Naturvidenskabens Standpunkt; 3)
den Erkendelse, at vor Forstand lige saavel som vore Sanser høre med til vor
Organisation, og at derfor selv den videnskabeligt korrigerede Sanseopfattelse
ikke giver os Tingen i sig --- det er Kants Standpunkt; 4) endeligt den
Erkendelse, at selve Modsætningen mellem Ting i sig og Fænomen er et Produkt af
vor Organisation, hvorfor der kan tvivles om dens Gyldighed. Denne fundamentale
Tvivl kan efter Lange ikke hæves ad Erkendelsens Vej. Den betegner Videnskabens
Grænse, ud over hvilken kun Digtningen, ikke Tænkningen kan føre. Lange
bebrejder Kant, at han blev staaende paa det tredie Stadium; Kants romantiske
Efterfølgere bebrejder han, at de »tumlede ud over det fjerde Stadium ind i
Digtningen og blandede Kritik og Digtning i deres Systemer istedetfor at bringe
dem i et klart afgrænset Forhold til hinanden«.

Ved denne Opfattelse har Lange dog overset, at naar vor Erkendelse paa selve
sin Grænse anvender de samme Principer som indenfor Grænsen, maa den
nødvendigvis fastholde Modsætningen mellem Subjekt og Objekt, Erkendelse og
Ting: i modsat Fald vilde den under en eller anden Form komme til at antage en
Skabelse af Intet, en Kausalitet ud fra en absolut Enhed. Det er en Lov for al
Erkendelse, at hvad vi forstaa, kun forstaas ved Hjælp af \textit{flere}
Præmisser eller Betingelser\textsuperscript{129}). Anerkender man ikke denne
Lov, hildes man i Vanskeligheder, %
% --- p. 550 ---
\opage{550}
som ligefra Nikolaus Cusanus af atter og atter dukke op i Tænkningens Historie.

Lange benytter ligesom Schopenhauer Idealismen (i Betydning af Subjektivisme)
til at lette Materialismen af Sadlen. Men han hylder ogsaa Idealismen i Ordets
praktiske Betydning. Det er ikke nok for ham at hævde Virkelighedens subjektive
Karakter; derved vilde jo Virkeligheden ikke blive anderledes, end den er, slet
eller god. Men han hævder, at ligesaa dybt begrundet i den menneskelige Natur
som Trangen til Opfattelse og Forklaring af Virkeligheden efter
naturvidenskabelige Love, er Trangen til at danne ideale Billeder og betragte
dem som Udtryk for den højeste Virkelighed. Den producerende Kraft, der ytrer
sig i Sansning og Tænkning, slaar sig ikke til Ro med det Givne, men higer
udover det for at danne en Idealverden. Mennesket kan ikke undvære en
Supplering af Virkeligheden. Alle Religioner og alle spekulative Systemer ere
Produkter af denne Trang. De adskille sig fra det videnskabelige
Virkelighedsbillede ved, at det ikke er den for Alle fælles Organisation,
Slægtens Natur, der virker ved deres Opstaaen, men individuelt-personlige
Tendenser og Fornødenheder, og tillige ved, at en fuldstændig
Erfaringsbekræftelse ikke er mulig. Det er af stor Betydning at fastholde denne
Forskel, for at man ikke skal slaa fast som bogstavelig Sandhed, hvad der kun
er et Symbol. Enhver Betydning for Erkendelsen af Virkeligheden maa fraskrives
de religiøse og spekulative Ideer. Men paa den anden Side maa der ogsaa skelnes
mellem saadanne Ideer og blotte Hjernespind. Skønt Lange enkelte Steder lader
Digtningen begynde, hvor Erkendelsen ender, hævder han dog andre Steder (1.
Udg. p. 539, 541. 2. Udg. II. p. 540), at Spekulationen staar imellem
Erfaringsvidenskaben og Poesien: hist raader Stoffet, her den frie Form;
Spekulationen søger at forene begge. Tydeligere er den Forskel, Lange finder i,
at Ideen skal bedømmes efter sin Værdi for vort aandelige Livs Udvikling, ikke
efter sin Begrundelse og Oprindelse (1. Udg. p. 346 f. 2. Udg. II. p. 545). Vi
betragte jo dog Kølner Dom som Andet og Mere end en Stenhob! Ja, Lange taler
endog om, at den praktiske Idealisme kunde hænge sammen med den %
% --- p. 551 ---
\opage{551}
ubekendte Sandhed, skønt paa en anden Maade end Materialismen gør det: Ideen er
Billede og Symbol for et uerkendeligt Absolut! (1. Udg. p. 346, 539, 541, 545.
--- 2. Udg. II. p. 494 o. fl. St.).

Atter her viser det sig, hvor vanskeligt det er at anbringe en absolut Grænse.
Lange bliver her uvilkaarligt Platoniker. Hvoraf ved han, at de Ideer, der for
os udtrykke det højeste Værdifulde, ere Afbilleder eller Symboler? Og paa den
anden Side --- dersom hans skarpe Forskel mellem Ideer og videnskabelig
Erkendelse holdes fast --- hvem vil saa arbejde paa Ideernes Udviking, hvis de
slet ingen real Betydning have? De spekulative Ideer kunne kun hævdes i deres
Betydning ved at ses som Hypoteser, i hvilke den menneskelige Tanke drager de
sidste Konsekvenser af, hvad der til den givne Tid staar fast for den. Det maa
da overlades nærmere Prøvelse, hvilken teoretisk og praktisk Betydning de kunne
have; men praktisk Betydning kunne de neppe have, naar de slet ingen teoretisk
Betydning have. Hvad endeligt selve den praktiske Betydning angaar, giver Lange
intet bestemt Kriterium, ingen Maalestok. Hvor han antyder sin etiske
Opfattelse, siger han, at den naturlige Sympati, der for Auguste Comte var
Etikens Grundlag, ikke er nok til en idealistisk Etik, og han betragter Ideen
om en Totalitet, til hvilken vi høre, som det aprioriske Grundlag for Etiken,
ved hvilket den moralske Lovs ubønhørlige Strenghed ene bliver forstaaelig. Men
strax bliver han igen betænkelig: fører Ideen ikke ofte vild, og kan det ikke
især overfor Religionssystemerne spørges, om det ikke er bedre simpelthen at
overlade sig til den naturlige Sympatis forædlende Virkning end at høre paa
Profetstemmer, der kun altfor ofte have ført til den frygteligste Fanatisme?
--- En nærmere Undersøgelse af Forholdet mellem Formalisme og Realisme i Etiken
har Lange ikke givet. I sine etiske Ideer var han Discipel af Kant, eller
snarere af Schiller, skønt han saa Begrænsningen af deres Opfattelse.

Sin Lære om de religiøse Ideers symbolske Karakter og praktiske Værdi anvender
Lange ikke blot paa den filosofiske Spekulation, men ogsaa paa de positive
Religioner. Naar Mate%
% --- p. 552 ---
\opage{552}
rialismen angriber Kristendommen som Kulsviertro, saa overser den --- hævder
Lange --- Muligheden af en friere, ideal Opfattelse af de kirkelige
Traditioners Indhold. Man skal se paa Ideernes Værdi og ikke paa deres
bogstavelige Formulering. I en vis Forstand ere Religionens Ideer
uforgængelige. Hvem vil gendrive en Messe af Palæstrina eller godtgøre, at
Rafaels Madonna er en Vildfarelse? I det inderligt bevægede Aandsliv hos de
virkeligt Fromme rører sig Noget, som den idealt sindede Fritænker maa forstaa
og maa sympatisere med! Lange ønsker en Omtydning af de kirkelige Lærdomme, i
Kants, Fichtes og Schleiermachers Aand, for at Kristendommens etisk-ideale
Indhold kan uddrages af Bogstavhylsteret; Hegels Religionsfilosofi var ham
derimod for konservativ. Direkte Polemik mod den positive Religion vil han
ikke, for at Folkets Enhed og Aandslivets Kraft ikke skal lide. Et af to, vil
han hellere for en Stund ofre Oplysningen end Kraften, naar blot Lærefriheden
sikres og Præstevælde holdes borte. Han har et aabent Blik for de store
Modsætninger i vor Tid og venter fra Filosofiens Side som Bidrag til at mildne
Kampen især mere indtrængende Forstaaelse af den menneskelige Udviklings Natur
og Love. Han slutter sit geniale Værk med at udtale det Haab, at alvorligt
Arbejde i denne Retning vil bære sin Frugt. ---

Den store Interesse for Kants Filosofi, som Langes Værk dels vidnede om dels
vakte, førte til et ivrigt Studium af den store Filosof, som Romantikens
Filosofi troede at have overvundet. Paa flere Punkter viste der sig Uenighed
om, hvorledes Kant var at opfatte. En rig Kantliteratur opblomstrede, der i
adskillige Henseender førte til en nøjagtigere Opfattelse af Kant, hans Opgave,
hans Metode og hans Stilling i Filosofiens Historie. Selv om man ofte fortabte
sig i Undersøgelser, der havde mere filologisk end filosofisk Betydning, saa
blev det dog af stor Betydning, at Kant forstodes rigtigere end før. Dette
skyldes især Arbejder af \textsc{cohen, fr. Paulsen, laas} og
\textsc{vaihinger}, hvortil senere kom de Bidrag til Belysning af Kants
Udvikling og Tankegang, som \textsc{b. Erdmann, r. Reicke} o. Fl. fremdrog af
Kants uudgivne Optegnelser og Breve, I den af Berlinerakademiet begyndte store
Udgave af Kants Vær%
% --- p. 553 ---
\opage{553}
ker vil Interessen for den kritiske Filosofi sætte sig et smukt Minde.

I en aandfuld Afhandling \textit{Was uns Kant sein kann} (Vierteljahrsschr. für
wissensch. Philos. 1881) har \textsc{Friedrich Paulsen} søgt at bestemme den
Betydning, som Kant kan have for vor Tid. Denne Betydning beror ikke paa en
ligefrem Genoptagelse af hans Lære, men paa en Behandling af Problemerne i hans
Aand. Den saakaldte \textit{Nykantianisme} betegner derfor ikke nogen sluttet
Skole, men en Bestræbelse for erkendelsesteoretisk Prøvelse af de Begreber, vi
operere med. \textsc{Wilhelm Wundt} har i sin Afhandling Was \textit{soll uns
Kant nicht sein?} (Philosophische Studien VII) taget Ordet mod nogle
Nykantianeres vel dogmatiske Tilslutning til den gamle Mester og paavist,
hvorfor Kants Lære ikke kan opstilles som en for vor Tids Videnskab gældende
Erkendelsesteori.

\begin{center}b. EUGEN DÜHRING.\end{center}

Aaret før Langes »Geschichte des Materialismus« udkom et af de mest udmærkede
erkendelsesteoretiske Arbejder i den sidste Halvdel af Aarhundredet, nemlig
Eugen Dührings \textit{Natürliche Dialektik.} Det tager forsaavidt sit
Udgangspunkt i den kritiske Filosofi, som det paa mere radikal Maade vil
genoptage Kauts %
% PRINTED AS IS: „Kauts“ (p. 553)
Problem, som Romantikens »Syndflod« havde skyllet bort. Blandt senere tyske
Arbejder nævner det kun Trendelenburgs »Logische Untersuchungen« med stor
Anerkendelse som Forgænger; af udenlandske Tænkere nævnes Comte og Stuart Mill,
som dog ikke særligt have behandlet Bogens Emne. Bogen er allerede nu en
literær Sjældenhed, da Forfatteren ikke har udgivet den paany, maaske fordi
hans Standpunkt senere har ændret sig i mere positivistisk Retning. Og dog er
det sikkert den bedste af alle Dührings Bøger, baade hvad Form og Indhold
angaar. Den er et interessant Indlæg til Belysning af Forholdet mellem kritisk
Filosofi og Positivisme, to Retninger, der brydes ikke blot hos Dühring, men i
hele den nyeste Tids Filosofi. Ved to Grundtanker nærmer Dühring sig allerede her %
% --- p. 554 ---
\opage{554}
Positivismen. Endnu skarpere end Kant og Schopenhauer hævder han, at
Begrundelsens Lov (den tilstrækkelige Grunds Princip) kun er en Lov for vor
Tænken, ikke for Virkeligheden, der er mere omfattende end vor Tænken. Og han
opstiller paa den anden Side den manglende Grunds Princip, hvorved Bevisbyrden
paalægges den, der opstiller en ny Paastand, idet det hidtil Givne bliver at
fastholde, saalænge der mangler Grund til at forlade det. Ved den første
Grundtanke afviser han den rationalistike %
% PRINTED AS IS: „rationalistike“ (p. 554)
Dogmatisme, ved den anden Idealismen og Dualismen; gennem begge arbejder han
sig frem til at bygge paa den givne Virkeligheds Grund uden Indblanding af
Forestillinger, som ikke ere hentede fra den. Han fortsætter Ludwig Feuerbachs
Værk og er Positivismens vigtigste Repræsentant i Tyskland. Selv kalder han sin
Filosofi for Virkelighedsfilosofi. De store Grundtræk af Virkeligheden, som
Erfaringen lærer os at kende, skulle danne Basis baade for Livsanskuelse og for
Livsførelse. Filosofien er for ham ikke blot Teori, men Udtryk for personligt
Sindelag.

Efter Dührings Opfattelse har det 19. Aarhundrede en afgjort reaktionær
Karakter --- det tekniske Omraade gør ene en Undtagelse. Det 17. Aarhundrede
stod for ham som den i videnskabelig Henseende klassiske Tid. Alvorlig
Videnskab begynder først med det, og vort Aarhundrede har Intet frembragt, som
kan fordunkle det. Dühring tænker her især paa Galilei, Huyghens og Newton i
Naturvidenskaben og paa Bruno, Hobbes og Spinoza i Filosofien. Det er
Grundtankerne hos disse store Forskere, vi endnu tære paa. Det er ejendommeligt
for Dühring, at de store, kraftige og frugtbare Begyndelsestanker for ham have
langt større Betydning end den specielle Gennemarbejdelse og
Erfaringsbekræftelse. Det hænger sammen med hans historiske Blik og med hans
Sans for det Geniale; men det fører ham til at miskende den Betydning, det har,
at de store Tanker finde deres Bekræftelse i specielle Erfaringer, hvilken
Betydning dog en Positivist allermindst burde miskende. Denne Tendens hos
Dühring, som sikkert blev næret ved, at at hans Blindhed hindrede ham i Selvsyn
og Selvorientering, tog til hos ham i Løbet af hans Udvikliug %
% PRINTED AS IS: „Udvikliug“ (p. 554)
og blev af skæbne%
% --- p. 555 ---
\opage{555}
svanger Betydning ikke blot for hans Filosofi, men ogsaa for hans ydre Livsforhold.

Medens han priser det 17. Aarhundrede som de store Tankers Tid, priser han det
18. Aarhundrede for den reformatoriske Anvendelse af disse Tanker. Naar hans
egen Tid forekommer ham ringe, saa vender han sig med des større Haab mod
Fremtiden. Han føler sig som Bebuder af en Retning i Tænkning og Liv, der først
i en fjern Fremtid vil kunne vinde Udbredelse. Hans store Tankeenergi og
Tankefylde forenes med Reformatorens Selvfølelse --- desværre ogsaa med en
Bitterhed og en Mistænksomhed, som hans Skæbne har næret, men som gøre hans
senere Skrifter i mange Partier uhyggelige ved grove Udfald mod dem, han kalder
sine Fjender. Hans Skrifter have et stærkt personligt Præg, baade i Ordets gode
og i dets mindre gode Betydning. Selv har han skildret sin Personlighed og sit
Liv i den mærkelige Bog\textit{ Sache, Leben und Feinde. Als Hauptwerk und
Schlüssel zu seinen sämmtlichen Schriften} (Karlsruhe und Leipzig 1882). Han
har Ret i, at hans Sag og hans Liv hænge nøje sammen, og til hans Ære skal det
siges, at Sagen hos ham staar foran Livet; mindre Ret har han maaske i, at hans
Fjender hænge saa nøje sammen med hans Sag og hans Liv, som han antyder i Titelen.

\textsc{Eugem Dühring} %
% PRINTED AS IS: „Eugem“ (p. 555)
er født i Berlin 1833. Hans Slægt stammede --- hvad han satte stor Pris paa ---
fra Sverig. Af sin Fader blev han opdragen i Rousseau'sk Aand, under
Indflydelse af en frisindet religiøs Tænkemaade. Opdragelsen udviklede
Selvstændighed og Fasthed i Tankegangen. Tidligt bleve Matematik og Astronomi
Hovedinteresser. Efter Faderens Ded kom han paa et Gymnasium, der tillige var
Pensionat, men hvor han var ligesaa lidt fornøjet med den aandelige som med den
legemlige Kost. Siden studerede han Jura og arbejdede en Tid som praktisk
Jurist. En Øjensygdom, der medførte Blindhed, gjorde denne Virksomhed umulig og
gav Anledning til, at han besluttede at benytte sit store Kundskabs- og
Tankeforraad til at betræde den videnskabelige Forfatterbane. Hans Hustru,
senere hans Søn, var hans Sekretær. En Række af filosofiske, nationaløkonomiske
og videnskabshistoriske Værker vidne om, %
% --- p. 556 ---
\opage{556}
at den energiske Tænker ikke forgæves satte sin Lid til de indre Hjælpekilder,
da de ydre lukkedes for ham. Af hans Skrifter ere, foruden det allerede
omtalte, at nævne \textit{Der Werth des Lebens} (1865), \textit{Cursusder %
% PRINTED AS IS: „Cursusder“, without a space (p. 556)
Philosophie} (1875) (omarbejdet til \textit{Wirklichkeitsphilosophie} 1895),
\textit{Logik und Wissenschaftstheorie} (1878), \textit{Ersatz der Religion
durch Vollkommeneres} (1883). Han virkede som Privatdocent ved
Berlineruniversitetet, hvor hans Forelæsninger hørte til de mest besøgte. Hans
Harme over at være bleven forbigaaet ved Besættelsen af Professorater og hans
Kritik af Universitetsforholdene bragte ham i en skarp Modsætning til de
ledende Kredse ved Universitetet, og da han mente, at Helmholtz havde tilranet
sig Æren for Opdagelsen af Loven om Energiens Bestaaen, der tilkom Robert
Mayer, i hvem Dühring saa det 19. Aarhundredes Galilei, --- fik hans
Indignation særligt Luft i gentagne voldsomme Angreb paa den berømte Fysiker.
Ministeriet straffede ham efter det filosofiske Fakultets Indstilling med at
berøve ham Retten til at holde Forelæsninger. At dette var en ubillig Haardhed
mod den blinde Filosof, synes efter Alt, hvad der foreligger, at være klart,
selv om hans Polemik var utilbørlig baade i Indhold og i Tone. Den Beskyldning,
Dühring rettede mod Helmholtz, bliver kun forklarlig ved, at han paa Grund af
sin Blindhed ikke fulgte omhyggeligt med Literaturen; ti Helmholtz havde i
Virkeligheden, saasnart han opdagede, at han havde en Forgænger i Mayer, ikke
forsømt nogen Lejlighed til at give denne den Ære, der tilkom ham. Men tillige
virkede her hos Dühring den Overvurdering af de første Ideer i Modsætning til
den exakte Gennemførelse, som vi ovenfor have omtalt; hvad der lod ham stille
Robert Mayer saa højt over alle hans Medbejlere, var den samme
Grundbetragtning, der lod ham sætte det 17. Aarhundrede saa højt over det 19.
Og endeligt medvirkede hans Mistænksomhed, der sandsynligvis næredes ved den
Isolation, Blindheden medførte, og som lod ham finde »Fjender« paa alle Kanter.
Næst Professorerne vare Socialdemokraterne og Jøderne hans »Fjender«.
»Modstandere«, hvem det var om Sandheden og ikke om egen Fordel at gøre, mener
han derimod ikke at have truffet paa sin Vej! (Leben, Sache und Feinde. Kap.
14). Han er %
% --- p. 557 ---
\opage{557}
saa sikker paa at have alle gode Magter paa sin Side, at han siger i Slutningen
af sin Selvbiografi: »Ueberhaupt sind meine Feinde wesentlich keine andern
Elemente gewesen, als diejenigen, welche auch Feinde der besseren Menscheit %
% PRINTED AS IS: „Menscheit“ (p. 557)
sind!«

Efter i nogle Aar at have virket ved frie Foredrag har Dühring nu trukket sig
tilbage til en lille By i Berlins Nærhed, beskæftiget med naturvidenskabelige
og literaturhistoriske Arbejder. En stor og ædel Kraft er her trængt til Side,
skønt ikke brudt, ved Ulykke og ved egen og Andres Skyld.

\begin{center}\textit{α. Erkendelsesteori.}\end{center}

I sin Undersøgelse af Erkendelsen arbejder Dühring (især i »Natürliche
Dialektik«) ganske i den kritiske Filosofis Aand. Han søger at paavise vor
Erkendelses Ejendommelighed i Form og Fremgangsmaade for at kunne afgøre,
hvorvidt vi tør betragte dens Former og Resultater som Udtryk for Tilværelsen.
Han forsøger altsaa en kritisk Prøvelse af Forholdet mellem Tænken og Realitet.

Et Hovedpunkt er det her for ham, at Tænkningen stedse søger efter
\textit{kontinuerlig Sammenhæng} og stedse søger at \textit{fortsætte} sin
Gang. I Matematiken opstaar derved Uendelighedsbegrebet, idet der ingen Grænser
kunne sættes for en Størrelses Aftagen eller Tiltagen. I Logiken komme vi i
Kraft af den tilstrækkelige Grunds Princip ind i en uendelig Række, idet det
ene Hvorfor stedse afløser det andet. Overfor denne Tendens til stadig
Sammenhæng og Fortsættelse stiller Dühring, hvad han kalder \textit{det
bestemte Antals Lov.} Ethvert \textit{virkeligt} Givet, ethvert
\textit{faktisk} Resultat er begrænset og udviser et bestemt Antal Dele i Tid
og i Rum. Uendelighed og Ubegrænsethed betyder kun Mulighed af Fortsættelse, en
Mulighed, som ikke altid er realt berettiget. Og hvor Fortsættelsen virkeligt
kommer, sker den gennem successive Tilføjelser af enkelte, bestemte Elementer.
--- Der fremtræder altsaa her efter Dühring en karakteristisk Forskel mellem
Tanken og Virkeligheden.

Tanken maa være smidig og utrættelig forat kunne følge Virkeligheden i alle
dens Svingninger, forat kunne stige og %
% --- p. 558 ---
\opage{558}
synke med Fænomenerne og opdage alle Afhængighedsforhold mellem disse. Men den
Lov, den maa følge for at kunne fyldestgøre sin Opgave, maa man --- indskærper
Dühring --- ikke overføre paa Naturen, selve den virkelige Tilværelse. Det ses
klart ved at sammenligne den reale Videnskab med den formale. I Astronomien
staa de bestemte givne Størrelser i Modsætning til den rene Matematiks
ubegrænsede Frembringelsesmulighed. I Kemien ere Atomtallene Exempler paa det
bestemte Antals Lov. Ja, alle virkeligt givne Tal og Størrelser ere endelige:
der er et bestemt Antal Kloder i Himmelrummet, et bestemt Antal Atomer i
Materien; i hvert Øjeblik er et aldeles bestemt Antal Tidsdele forløbet; Jorden
har bevæget sig om Solen et vist bestemt Antal Gange; Aarsagsrækken bestaar af
et endeligt Antal Led.

Med stor Resoluthed drager Dühring af det bestemte Antals Lov den Slutning, at
Naturprocessen, Indbegrebet af alle Forandringer i Naturen, maa have havt en
Begyndelse. En uendelig Regres kan ikke tænkes. Forud for den af Forandringer
opfyldte Tid ligger da en evig Væren, hvor Differenserne, de indre Forskelle,
der medføre Succession og Forandring, endnu ikke fremtraadte \textit{som
successive.} Samtidig Mangfoldighed kan meget vel have været til Stede og maa
have været til Stede, ti den indeholder Spirerne (eller, som Dühring siger,
Rødderne) til Alt, hvad vi betegne som Proces og opfatte som bestemt Existens
(Wirklichkeitsphilosophie. p. 15). Tiden og Aarsagsrækken bleve først til ved
Overgangen fra denne ensartede og uforanderlige Tilstand til Udfoldelsen af
Differenser og Forandringer. --- Denne regressive Afslutning af Aarsagsrækken i
en uforanderlig Tilstand er for Dühring tilfredsstillende ikke blot fra
Tankens, men ogsaa fra Følelsens Synspunkt.

Men Dühring har fremhævet Modsætningen mellem Tænkningens ideale Kontinuitet og
den stadigt i Sondring og Udstykning givne Erfaringsvirkelighed saa stærkt, at
han ender i en uløselig Gaade. Hvorledes er nemlig en Overgang fra absolut
Ligevægt og Ensartethed til Succession og Forandring mulig? Han er her i sin
positivistiske Tænkning kommen til at staa overfor det samme Problem, som Böhme
og Schelling havde fremsat %
% --- p. 559 ---
\opage{559}
i mytologiserende og spekulativ Form. Han vil mildne Paradoxet ved at henvise
til, at hver Gang, der indenfor vor Erfaring opstaar et nyt Fænomen [en ny
Kvalitet], have vi egentligt en absolut Begyndelse. (Cursus p. 79. Logik p.
191). Men Forskellen er, at det Nye og Begyndende i vor Erfaring vækker
Forundringen og giver Tanken et Problem, der kan sysselsætte den, medens hin
skarpe Modsætning mellem absolut Identitet og Forandring en Gang for alle gør
Tankens Arbejde haabløst. Ikke blot en Forskel, men en Strid mellem Tænken og
Virkelighed kommer Düring %
% PRINTED AS IS: „Düring“ (p. 559)
her til at statuere. --- Naar man har Valget mellem at antage en saadan Strid
og at antage, at Tilværelsen er uafsluttelig ligesom Tænkningen, synes denne
sidste Mulighed at være at foretrække. Der finder et stadigt Vexelspil Sted
mellem Tanken og Virkeligheden: der viser sig, efterhaanden som Forskningen
skrider frem, ny Kontinuitet bag den tilsyneladende Diskontinuitet, og
forsaavidt bekræfter Virkeligheden Tankens Ideal; men tillige opdages ny
Diskontinuitet, hvor der syntes at være Kontinuitet i det Givne, og Problemet
om, hvorvidt Virkeligheden kan forstaas, stilles da paany. Dogmatikeren har hos
Dühring faaet Overhaand over den kritiske Filosof paa Grund af hans store Iver
for at lade Virkeligheden komme til sin Ret.

Paa den anden Side lægger Dühring meget stor Vægt paa den indre Sammenhæng
mellem Tænkningen og Virkeligheden. Tilværelsen selv fortsætter sig ind i vor
Tænkning, hvor denne udfolder sig efter sit sande Væsen. Til Forholdet mellem
Forudsætninger og Slutning (Grund og Følge) svarer i de virkelige Processer
Forholdet mellem Aarsag og Virkning. Til Logikens Identitetsprincip svarer den
enkelte Tings Identitet med sig selv, uagtet den forekommer i forskellig
Sammenhæng. Til vor Tænknings Kombinationer og Deduktioner svarer Naturens
Enhed og Vexelspillet af Kræfter ved Fænomeners Opstaaen. Vi opfatte herved
ikke Naturen i Analogi med menneskelig Bevidsthed, men som virkende paa en
saadan Maade, at den menneskelige Erkendelse faar et til dens egen Beskaffenhed
svarende Materiale at undersøge. Og selv om vor Analyse aldrig vil kunne
udtømme den store Sammenhæng i Naturen, ophører %
% --- p. 560 ---
\opage{560}
ikke Nødvendigheden af at hævde Slægtskabet mellem det, der virker i Tingene,
og det, der virker i Forstanden.

Hvad der fremtræder i vore videnskabelige Principer, er efter Dühring Frugten
af Kræfter, der længe have arbejdet i en bestemt Retning, før de komme til klar
Bevidsthed. Bevidsthed og Tanke bero paa Noget, der ligesaa lidt er Bevidsthed
og Tanke, som den bevægende mekaniske Kraft selv er Bevægelsesfænomen. Men det
er en berettiget Hypotese at forestille os den ikke i Bevidstheden fremtrædende
Grund til vor Tankevirksomhed som stemmende med Grunden til, at
Natursammenhængen viser sig for os at være forstandsmæssigt ordnet. (Nat. Dial.
p. 155, 225. Logik p. 173. Leben, Sache etc. p. 303). Hvad der faar os til at
tænke, er altsaa det Samme, der faar Naturen til at virke.

Denne sidste Tankegang faar nu efterhaanden en saadan Overvægt over Dühring, at
han i sine senere Skrifter fører en voldsom Polemik mod den kritiske Filosofi,
der hævder, at vi ikke erkende Tingene i sig selv, fordi vor Erkendelse ikke
blot er bestemt ved Tingenes Natur, men ogsaa ved vor egen Natur. Han ser i
denne Opfattelse en Forsumpning af Videnskaben, en hyklerisk Lempen af
Tænkningens Resultater efter hvad Trangen til at antage en hinsidig Verden
fordrer. Atter her er det det bestemte Antals Lov, der ligger til Grund: Tallet
er Exempel paa, at vi raade over Begreber, der udtrykke det Værende paa den
Maade, at der imellem Virkeligheden og Begrebet slet ikke finder noget ved
Jegets Beskaffenhed betinget Indskud Sted. (Logik p. 165).

Ligesaa skarpt som Dühring ellers fremhæver Modsætningen mellem Tænken og
Væren, ligesaa bestemt fremhæver han her Muligheden af, at de fuldstændigt
kunne dække hinanden. Hist har han nærmest den formale Viden for Øje, her den
reale. Men de kunne ikke holdes helt ude fra hinanden, hvad der viser sig deri,
at netop de bestemte Tals Lov stiller Forskningen nye Opgaver --- nemlig
Udfindelse af, hvorfor Forandringer netop ske, naar disse bestemte Tal ere naaede.

Det er Dührings Fortjeneste af den nyere Erkendelsesteori, at han sætter
Problemet om Forholdet mellem Kvalitet og Kon%
% --- p. 561 ---
\opage{561}
tinuitet, mellem de bestemte Vendepunkter og de sammenhængende Processer i saa
klart og tydeligt et Lys. Hans hele Standpunkt er bestemt ved en Brydning af
Kriticisme og Positivisme, og indeholder en Opfordring til en nærmere
Undersøgelse af Forholdet mellem disse to Retninger.

\begin{center}\textit{β. Verdensopfattelse.}\end{center}

Dühring lægger stor Vægt paa Udviklingen af en \textit{Verdensskematik,} der
skal give et udformet Billede af den virkelige Tilværelses almindelige Træk.
Deri bestaar Virkelighedsfilosofiens Opgave. Den skal føre til et Erfaringens
System, medens den tidligere Metafysik var en Slags Alkymi, der ad mystiske
Veje vilde udgrunde Tilværelsens Væsen. --- Paa dette Punkt minder han om Comte
og Spencer. Han er i sin »Skematik« mere abstrakt end Spencer og mindre
original end Comte.

Mit System, siger han, kender ikke to Virkeligheder, men kun en eneste, Naturen
og dens Dele. Han nærmer sig til Materialismen ved at fremhæve, at Alt i Verden
har en materiel Side, og at man ved at følge de materielle Fænomener i deres
anskuelige Sammenhæng maa kunne finde Alt, hvad der har Realitet. Men han vil
ikke blive staaende ved de blotte Fænomener. Derved fjerner han sig fra »den
sædvanlige, ret plumpe Materialisme«. Ofte bruger han vel Ordet
\textit{Materie} i Betydning af \textit{Væren, Verdensmedium, Bæreren af alt
Virkeligt.} Men han fremhæver udtrykkeligt, at Materiebegrebet maa udvides og
uddybes saaledes, at det bliver klart, at Anlæggene til Bevidsthedsfænomenerne
ligesaa vel findes i Naturen som Mulighederne til alle andre Fænomener.
\textit{Det filosofiske Materiebegreb} indeholder Mere end Fysikernes og
Kemikernes, da det skal være Begrebet om det, der ligger til Grund for alle
legemlige og aandelige Tilstande, og som omfatter den hele og fulde
Virkelighed. (Cursus p. 62, 75. Leben, Sache etc. p. 302 f.). Materialismen er
for Dühring kun en »Piedestal«, et realt Grundlag, ydre, mekanisk Ramme for
Verdensopfattelse og Livsanskuelse, ikke selve den udtømmende, fyldige
Forklaring af Tilværelsen. %
% --- p. 562 ---
\opage{562}

En Lære, man ikke venter at finde hos Dühring, er Antagelsen af
\textit{Formaal} i Naturen. For at forstaa ham ret paa dette Punkt maa man
lægge Mærke til, at han skelner mellem Formaal (Zweck) og Hensigt (Absicht).
(Logik p. 203). Naar han tillægger Formaalsbegrebet Gyldighed for vor
Naturopfattelse, er det, fordi de i Naturen virkende Kræfter og Dispositioner
stedse udvise visse Tendenser, arbejde i visse Retninger, og først og fremmest,
fordi de \textit{virke sammen} og derved gøre, at Naturen fremtræder som en
systematisk Enhed. Hine Tendenser føre ikke altid --- ligesaa lidt i den ydre
Natur som i det bevidste Menneskeliv --- til afsluttede Resultater. Men hvad
enten Maalet naas eller ikke, fremtræder \textit{den kombinerede Virken af
mangfoldige Elementer} som et karakteristisk Hovedtræk ved alle Processer. Ved
Opfattelsen af Naturen maa man ikke blot lægge Mærke til de almindelige Love,
men ogsaa til Retninger og Tendenser af de Processer, som gaa for sig, og til
de Typer og Former, i hvilke Naturvirksomheden fremtræder. Naturen danner en
Trinrække, idet de lavere Tilværelsesformer ere Grundlag for de højere.
Maalestokken, ved hvilken vi skelne mellem lavere og højere Former, kan kun
søges i de forskellige Tendensers Resultater, i den Grad, i hvilken de føre til
deres naturlige Maal. Og det sidste Maal maa være Opstaaen af Væsener, som ikke
blot existere og virke, men ogsaa blive sig deres Existens og deres Virken
bevidste: kun som Grundlag for bevidste Væseners Historie har det mekaniske
Naturheles Historie en Mening; en gennemgaaende bevidstløs Verden vilde ---
siger Dühring (Cursus p. 104) --- være en taabelig Halvhed, et Teater uden
Skuespillere og Tilskuere. Og at paa den anden Side Bevidsthedsfænomenerne
væsentligt ere at betragte som Formaal, ses af, at ubevidst og bevidst Virken
kunne udrette Et og det Samme; det vilde jo være blot Luxus, at Mennesker og
Dyr blive sig deres Tilskyndelssr %
% PRINTED AS IS: „Tilskyndelssr“ (p. 562)
og Motiver bevidste, naar det kun drejede sig om ydre Resultaters Opnaaelse, og
naar vi her ikke stode overfor Noget, der i sig selv var Formaal. Særligt have
vore Følelser (f. Ex. Elskovsfølelsen, den sociale Følelse) ikke blot deres
Værdi ved at være Midler eller Motiver, der føre til Handlinger, som ere
nødvendige for Livets %
% --- p. 563 ---
\opage{563}
Bestaaen; de have umiddelbar Værdi som Udtryk for et højere Livstrin end det
rent elementære. (Cursus p. 158 f. Der Werth des Lebens. 2. Aufl. p. 241).

Hvis man vilde indvende mod denne sidste Tanke, at Bevidsthedslivet dog ikke
altid er forbundet med Lyst, men at Smerten indtager saa stor en Plads, at
Bevidsthedslivet ikke synes at være et virkeligt »Formaal«, saa vilde Dühring
henvise til, hvad han kalder \textit{Differensens Lov,} en Lov, som under
forskellige Former og Benævnelser spiller en stor Rolle i den nyere Psykologi.
(Den hensigtsmæssigste Benævnelse turde være: Forholdsloven). For Dühring er
den ikke blot en psykologisk Lov, men en Verdenslov. Kræfternes Antagonisme er
et Hovedtræk i Verdensskematiken. Al Kraftytring, al Bevægelse og al Udvikling
er betinget ved Afstand, Forskel, Modsætning. Ogsaa al Bevidsthed forudsætter
Forskel. Uden Differens ingen Fornemmelse. Indenfor Bevidstheden gælder det
Samme. Livets Spil vilde tabe sin Tiltrækning, hvis der ikke var Modstand og
Hindring at overvinde, og hvis Trang og Tilfredsstillelse ikke afløste
hinanden. Det er Tilværelsens Rytmer og Differenser, der gør den værdifuld for
os; vor Livsfølelse sættes i Bevægelse ved Overgangene mellem de indbyrdes
forskellige Tilstande. Uden det Drøje, Bitre og Smertelige vilde den dybe
Tilfredsstillelse ved Livet ikke føles. --- Den samme Betonen af Følelsens
Bestemthed ved Modsætninger, der gjorde Schopenhauer til Pessimist, gør Dühring
til Optimist. Schopenhauer er, hvad »Differensloven« angaar, konsekventere end
Dühring, forsaavidt han antager, at Uroen og Bevægelsen ligger som en
oprindelig Tendens i Verdensgrundens Natur, medens det for Dührings Filosofi
bliver en stor Gaade, hvorledes successive Differencer og rytmisk skiftende
Tilstande kunne have afløst den hvilende Væren, som den antager i Kraft af det
bestemte Antals Lov.

De forskellige Kræfters \textit{Kombination} er det, der gør Udvikling mulig.
Hvad enten denne Kombination fører til en Ligevægtstilstand eller til en ny
Bevægelse, saa bevares de enkelte Kræfters Tendenser i Udviklingens Resultat.
En stadig Strid mellem Kræfterne vilde være en Absurditet --- men den %
% --- p. 564 ---
\opage{564}
findes ikke i Naturen. Undgaaen af det Absurde er et naturlogisk Princip. Kun i
Undtagelsestilfælde fører dette Princip til, at uholdbare Former tilintetgøres
og opløses. Som Regel gaar Naturen den positive Vej. Derfor er, efter Dührings
Opfattelse, Ideen om Kampen for Tilværelsen falsk: den fremhæver jo ensidigt
det negative Synspunkt og lægger Vægten paa Antagonismen i Stedet for paa
Kombinationen. Dühring giver Lamarcks Teori Fortrinnet frem for Darwins.

\begin{center}\textit{γ. Etik.}\end{center}

Det Etiske har efter Dühring sit Grundlag i de menneskelige Drifter. Kun sit
Grundlag --- ti Naturen kan fejle, og derfor er en nærmere Udvikling og
Korrigeren nødvendig. Jo højere et Væsen staar, des større Mulighed er der for,
at det tager fejl. Desuden gælder det etisk set om at kombinere Elementer, der
fra Naturens Haand endnu ikke ere forbundne, men maaske endog i indbyrdes Strid.

Spirerne til det Gode finder Dühring, ligesom Comte, i de sympatiske
Instinkter. Naturen har selv sørget for, at fremmed Lidelse bestemmer vor egen
Følelse paa smertelig Maade, og der gives maaske ingen Følelse, der under
Kulturens Gang har faaet en saa iøjnefaldende Udvikling som denne. Den fører
ikke til at kue og ophæve vor egen Individualitet. Tvertimod bestaar den etiske
Udvikling baade i en Individualisering og i en Socialisering; begge høre
nødvendig sammen, da den fulde Udvikling af den enkelte Individualitet først
bliver mulig i et højere udviklet Samfund. Under ufuldkomnere Samfundsforhold
hindres de Enkeltes frie, ejendommelige Udvikling. Men da vore nuværende Stater
væsentligt ere byggede paa Tvang, hæmme de stadigt denne Udvikling. Et frit
Samfund, til hvilket Spirerne ere givne i de ved fri Sammenslutning dannede
Organisationer, vil først baade lade Individualiteten og Fællesskabet komme til
sin Ret. I et saadant Samfund ville alle Produktions- og Konsumtionsforhold
være inddragne under social Ledelse, saa at den Enkeltes hele Interesse kan
knyttes til Arbejdet, ikke til Arbejdets Udbytte. Først derved bliver %
% --- p. 565 ---
\opage{565}
en gennemgaaende Forædling af Livet muligt. Fra Socialismen adskiller Dührings
Opfattelse sig dels ved, at han tænker sig frie Arbejdskommuner, ingen absolut
Centralisation af Produktionen\textsuperscript{130}), dels ved, at han ikke
anser det for nødvendigt, at Udviklingen afbrydes ved et historisk Mirakel:
Fremtidens Bygning vil efter hans Opfattelse opstaa ikke (som Karl Marx mener)
derved, at det Onde tager til, men derved, at det Gode i Stilhed voxer. Ogsaa
paa dette Punkt er Dühring afgjort Optimist, saa megen Nederdrægtighed han end
mener at kunne konstatere i Verden. Han er en bestemt Modstander af
Pessimismen. Kun Indignationspessimismen (Entrüstungspessimismus), saaledes som
den repræsenteres af Byron og fremkaldes ved bestemte sociale Tilstande, har
hans Billigelse, hvorimod han foragter den romantiske »Hinsidighedspessimisme«
og den i de højere Stænder hyppige »Raaddenskabspessimisme«.

Udover Følelsen for de gode Kræfter, der arbejder sig frem indenfor den
menneskelige Verden, opstaar der en \textit{universel Affekt,} naar Blikket
hæver sig til den store Enhed, af hvilken den menneskelige Verden kun er en
enkelt Forgrening, --- naar Tænkeren føler sig som hørende til denne Enhed og
bevæget af dens Kræfter, --- og naar Tanken om hans egen Skæbne gaar op i
Tanken om den store Tingenes Orden, indenfor hvilken saa mange gode Spirer have
kunnet udfolde sig. ---

Ikke blot ved de Problemer, den stiller, men ogsaa ved den ædle, antike Stil, i
hvilken den er gennemført, og ved den Sammenhæng, den sætter mellem Tænkning og
Personlighed, staar Dührings Filosofi som et af de mest karakteristiske
Tankeforsøg i den Menneskealder, som fulgte efter Romantikens Fald og
Naturvidenskabens nye Opsving ved det 19. Aarhundredes Midte. Aaret 1880 danner
her et karakteristisk Vendepunkt, idet en Kreds af betydeligere Forskere
afslutte deres Virksomhed omtrent ved denne Tid, nemlig, foruden Dühring,
Lotze, Fechner, Darwin og Spencer; kort forud var Lange bukket under. I et
særligt Skrift skulle nogle Hovedformer af Filosofiens Historie i det følgende
Fjerdedelsaarhundrede fremstilles. 

% ======== NOTER TIL ANDET BIND (vol. II, pp. 566-600): the notes to the Tenth Book, 113-130 ========
% (p. 596 begins with the end of the notes to the Ninth Book, 110-112, not transcribed here)
\clearpage
\begin{center}\large Noter til Tiende Bog\end{center}
% (heading supplied: the print's heading, "NOTER TIL ANDET BIND.", stands on p. 566, before the notes to the First Book)

\begin{small}
% --- p. 596 ---
\opage{596}
\textsuperscript{113}) S. 495. Om »de store Opdagelsers Periode« i Slutningen
af det 18. Aarhundrede se \emph{Rasmus} \emph{Pedersen}: \textit{Planternes
Næringsstoffer.} Historisk Indledning. København 1883. p. 44 ff.

\textsuperscript{114}) S. 496. Se om Vitalismen i dens forskellige Former
\emph{Louis} \emph{Peisse}: \textit{La Médicine et les Médicins.} Paris 1857.
I. p. 226---297. --- Det maa vel fastholdes, at der her ved Vitalismen menes en
Retning, der finder \textit{Forklaringen} af Livsfænomener ved Antagelsen af en
»Livskraft«, --- ikke %
% --- p. 597 ---
\opage{597}
en Retning, der hævder, at der er Meget i Livsfænomenerne, der hidtil ikke er
forklaret ved fysiske og kemiske Aarsager. Se min Afhandling \textit{Om
Vitalisme} (i »Mindre Arbejder«).

\textsuperscript{115}) S. 497. Kfr. \emph{Mayers} Brev til Griesinger 20. Juli
1844: »Wenn ich sage, Wärme lässt sich in Bewegung verwandeln, und umgekehrt,
so will dies nichts meinen, als dass zwischen Wärme und Bewegung finden hin und
her dieselben \textit{qvantitativen Beziehungen statt.} (\textit{Kleine
Schriften und Briefe} von \emph{Robert} \emph{Mayer}. Herausg. von
\emph{Weyrauch}. Stuttgart 1893. p. 225). --- I sin Afhandling \textit{Robert
Mayers Entdeckung und Beweis des Energieprincipes} (i Festskriftet til Sigwart.
Tübingen 1900) belyser A. \emph{Riehl} Mayers Lære fra et erkendelsesteoretisk
Synspunkt.

\textsuperscript{116}) S. 499. If the resistance to electrolysis wich %
% PRINTED AS IS: „wich“ (p. 597)
is over and above that due to chemical change were not accounted for elsewhere,
\textit{it would prove the annihilation of a part of the power of the circuit,
without any corresponding effect.} We shall see that this is not the case, but
that in the evolution of heat, where the excess of resistance takes place, an
exact equivalent is restored. (\emph{Joule}: \textit{On the Heat evolved during
the Electrolysis of Water.} 1843. Scientific Papers. I. London 1884. p. 115).
\textit{We might reason, a priori, that absolute destruction of living force
cannot possibly take place,} because it is manifestly absurd to suppose that
the powers with which God has endowed matter can be destroyed . . . but we are
not left with this argument alone. decisive as it must be to every unprejudiced
mind. (\emph{Joule}: \textit{Matter, Living Force and Heat.} 1847. Scientif.
Pap. I. p. 269). --- Joules Standpunkt er altsaa ikke principielt forskelligt
fra Mayers og Coldings. For denne Sidstes Vedkommende smlgn. en Udtalelse i
hans \textit{Undersøgelse om de almindelige Naturkræfter og deres gensidige
Afhængighed.} (Videnskabernes Selskabs Skrifter. Femte Række. Naturv. mathem.
Afd. II. p. 129): »Den Tanke, at en Virksomhed skulde kunne forsvinde i det
Legemlige uden igen at fremtræde som virkende Aarsag, forekommer mig
fornuftstridig«. --- Det er i denne Sammenhæng af Interesse at erindre, at
ogsaa \emph{Lavoisier} i sin Tid gik ud fra en Grundsætning, han betragtede som
selvfølgelig, nemlig den, at et sammensat Legemes Vægt er lig Vægten af de
Stoffer, af hvilke det bestaar. Saa selvfølgelig denne Forudsætning var, saa
maate den dog bevises for dem, der antog Varmen for et Stof; ti hvis denne
Antagelse var rigtig, maatte den Udskillen af Varme, som ledsager en kemisk
Forbindelse, medføre en Vægtforringelse. Smlgn. \emph{Ernst} \emph{von}
\emph{Meyer}: \textit{Geschichte der Chemie.} Leipzig 1889. p. 137.

\textsuperscript{117}) S. 500. \emph{Mayer}: \textit{Kleinere Schriften und
Briefe.} p. 339 f. 460. \textit{Die Mechanik der Värme.} 3. Ausg. p. 356 f. ---
\emph{Colding}: \textit{Naturvidenskabelige Betragtninger over Slægtskabet
mellem det aandelige Livs Virksomheder og de almindelige Naturkræfter.}
(Oversigt over det kgl. danske Videnskabernes Selskabs Forhandlinger. 1856). p.
155. 166. --- \emph{Joule}: \textit{Scientific Papers.} I. p. 269. 273. ---
Hvad Mayer angaar, synes det efter hans Foredrag i Innsbruck (Die Mech. d.
Wärme p. 357), at han ikke har %
% --- p. 598 ---
\opage{598}
antaget noget Brud paa de materielle Fænomeners Række, men har antaget en
Parallelisme mellem de aandelige Processer og Hjerneprocesserne.

\textsuperscript{117\,b}) S. 506. Den Konsekvens, Czolbe drager af
Materialismen, drog allerede \emph{Henry} \emph{More}, idet han (\textit{De
anima.} Roterodami 1677. p. 64) opstiller følgende Grundsætning: »Dersom
Sansning kan tillægges Materien, saa er en materiel Dels Bevægelse eller
Tilbagevirken paa en anden materiel Del, eller dog Fortsættelsen af denne
Bevægelse, Et og det Samme som Sansning og Opfattelse. Denne Sætning . . .
bestyrkes ved min vigtigste Modstander Hobbes' Stemme.« Han henviser saa til
Hobbes' Lære om Sansningen som en Bevægelse. Hobbes negtede dog hin Konsekvens,
idet han udtrykkeligt siger, at den Sætning, at Sansning er Bevægelse, ikke kan
vendes om. (\textit{De corpore} XXV, 5).

\textsuperscript{118}) S. 518. Om Ubestemtheden hos Lotze paa dette Punkt se
\emph{Max} \emph{Wentscher}: \textit{Lotzes Gottesbegriff und dessen
metaphysische Begründung.} Halle 1893. p. 19---29.

\textsuperscript{119}) S. 520. \emph{Reinhold} \emph{Geijer} henledte i sin
Afhandling \textit{Hermann Lotzes tankar om tid och timlighet i kritisk
belysning} (Lund 1886) Opmærksomheden paa de skiftende Standpunkter, Lotze
synes at have indtaget overfor Spørgsmaalet om Tidens absolute Gyldighed. I min
Afhandling \textit{Lotze og den svenske Filosofi} (Letterstedtske Tidsskrift
1888) sluttede jeg mig til Geijers Fortolkning, skønt jeg maatte opfatte selve
Problemet anderledes end min svenske Kollega. Dog synes Sagen at stille sig
anderledes, end Geijer, og jeg med ham, havde Grund til at antage. Efter hvad
\emph{Richard} \emph{Falckenberg} har oplyst i sin Afhandling \textit{Die
Entwickelung der Lotzeschen Zeitlehre} (Zeitschrift für Philosophie und
philosophische Kritik. 105. Bd.), skrive nemlig de efter Lotzes Død udgivne
\textit{Grundzüge der Metaphysik} sig allerede fra 1865, og dermed bortfalder
en væsentlig Støtte for den Antagelse, at Lotze først mod Slutningen af sit Liv
skulde have negtet Tidens Realitet. Naar han i tidligere Skrifter uden videre
taler om Tidsforholdet som absolut gyldigt, da skal dette i Følge Falckenberg
være en berettiget Brug af populære Forestillinger, hvor Sammenhængen ikke
nødvendiggjorde en stringentere Talebrug. Lotzes definitive Opfattelse har
været den i \textit{Drei Bücher der Metaphysik} (1879) og i \textit{Grundzüge
der Religionsfilosofi} (fra 1878---79) fremstillede, efter hvilken Successionen
gælder for de endelige Væsener, medens Guddommen er hævet over enhver
Tidsforskel. Dén afviger fra hans ældste Fremstilling (\textit{i Metaphysik.}
1841) ved at hævde denne Distinktion, som unegteligt ikke gør dette hele
Spørgsmaal klarere (skønt den \textit{psykologiske} Distinktion mellem
Succession og abstrakt Tidsform vel lader sig forsvare). --- Ligesaa lidt som
min svenske Kollega formaar jeg at forbinde nogen Mening med Udtrykket »ein
zeitloses Geschehen und Wirken« (Grundzüge der Metaphysik § 58). Min tyske
Kollega finder derimod en betydningsfuld Tanke deri.

\textsuperscript{120}) S. 520. Smlgn min \textit{Psykologi} 4. p. 80.
\textit{Psykologiske Undersøgelser.} p. 91 f. \textit{Filosofiske Problemer}
(Universitetsprogram 1902). p. 62---65. --- %
% --- p. 599 ---
\opage{599}
Interessant er det, at Grundlæggeren af den metafysiske Idealisme, Leibniz, i
sin tidligere Tid, ligesom Spinoza, lagde stor Vægt paa, at Toheden
Aand---Materie er rent faktisk, ikke logisk nødvendig. \textit{Philos. Schr.}
ed. Gerhardt I. p. 237, 242, 268. Smlgn. \emph{Ludw}. \emph{Stein}:
\textit{Leibniz und Spinoza.} p. 93 f.

\textsuperscript{121}) S. 523. Se herom \emph{Reinhold} \emph{Geijer}:
\textit{Hermann Lotzes lära om rummet} (Nyt svensk tidskrift 1880) og
\textit{Darstellung und Kritik der Lotzeschen Lehre von den Localzeichen}
(Philosophische Monatshefte 1885). Kfr. min Afhandling \textit{Lotze og den
svenske Filosofi} (Letterstedtske Tidsskrift 1890) og min \textit{Psykologi.} V
C, 7 c.

\textsuperscript{122}) S. 524. \textit{Mikrokosmus} Buch III, Kap. 1: »Warum
sollte nicht ein Atom des Nervensystems ebenso auf die Seele oder sie auf jenes
stossen und drücken können, da doch jeder gemeine Stoss und Druck sich für die
nähere Betrachtung nicht als ein Mittel zur Wirkung, sondern nur als die
anschauliche Form eines viel zarteren Ereignissen zwischen den Elementen ausweist?«

\textsuperscript{123}) S. 526. \emph{Kresto} \emph{Krestoff}: \textit{Lotzes
metaphysischer Seelenbegriff.} Halle 1890. p. 46 f.; 73. --- \emph{Max}
\emph{Wentscher}: \textit{Lotzes Gottesbegriff und dessen metaphysische
Begründung.} Halle 1893. p. 11 ff. --- Den Modsætning, i hvilken Lotzes
Opfattelse fra første Færd staar til Herbarts, trods alle de Berøringspunkter,
den har med denne, vidner ogsaa om, at der ikke i Aarenes Løb er foregaaet
nogen principiel Ændring i Lotzes Tankegang. Smlgn. herom \emph{Max}
\emph{Nath}: \textit{Die Psychologie Hermann Lotzes in ihrem Verhältniss zu
Herbart.} Halle 1892. --- Med Lotzes etiske Monisme hænger hans Optimisme
sammen. Uagtet han selv oftere fremhæver det fysiske og moralske Ondes Realitet
som en Hindring for Gennemførelsen af sin Verdensanskuelse, har dog en ellers
velvillig Kritiker anket over, at han mangler Interesse for »alle Religioners
Grunddogma, Forsoningen«, fordi han ikke tilstrækkeligt tager Hensyn til det
Onde. G. \emph{Vorbrodt}: \textit{Principen der Ethik und Religionsphilosophie
Lotzes.} Dessau und Leipzig 1891. p. 97.

\textsuperscript{124}) S. 528. Interessante biografiske Notitser hos A.
\emph{Elsas}: \textit{Zum Andenken Gustav Theodor Fechners} (Grensboten 1888).
En smuk Karakteristik af sin store Forgænger gav \emph{Wundt} i sine Mindeord
ved hans Grav (trykte i Philosophische Studien«. IV) og senere, mere udførligt,
i sin Tale paa Hundredaarsdagen for hans Fødsel (\textit{Gustav Theodor
Fechner.} Rede. Leipzig 1901). En indholdsrig Biografi har J. E. \emph{Kuntze}
givet (G. \textit{Th. Fechner.} Ein deutsches Gelehrtenleben. Leipzig 1892), og
\emph{Kurt} \emph{Lasswitz} har givet en kritisk indgaaende Fremstilling af
Fechners videnskabelige Standpunkt \textit{(G. Th. Fechner.} Stuttgart 1896,
det første Bind af »Frommanns Klassiker der Philosophie«).

\textsuperscript{125}) S. 532. I min ældre Fremstilling (\textit{Filosofien i
Tyskland.} 1872. p. 264 f.) har jeg af nogle Ytringer i \textit{Zendavesta} (II
p. 352 ff.) ladet mig forlede til at sige, at Fechner antager en Vexelvirkning
mellem Aand og %
% --- p. 600 ---
\opage{600}
Materie. Hvad Fechner der udvikler, er dog kun Vigtigheden af at skifte
Standpunkt og at udfinde Formlen for de to Siders indbyrdes Forhold. Ligesaa
lidt der som senere antager han en Overgang fra den ene Side til den anden. I
\textit{Elemente der Psychophysik} I p. 8 hedder det: »Wir nennen das
Psychische \textit{Function} des Physischen, \textit{davon abhängig} und
umgekehrt, insofern eine derartige constante oder gesetzliche Beziehung
zwischen beiden besteht, dass von dem Dasein und den Veränderungen des Einen
auf die des Anderen geschlossen werden kann.« For denne Art Afhængighed bruger
Pechner %
% PRINTED AS IS: „Pechner“ (p. 600)
oftere (f. Ex. Elemente I p. 18; II p. 393) Udtrykkene
\textit{Simultanbedingungen, simultane oder Wechselabhängigkeit.} Efter hans
Opfattelse kunne de til Bevidsthedsvirksomhederne svarende Hjerneprocesser ikke
(som Lotze og Spiritualisterne mene) betragtes som udløsende Indtryk (Reize)
overfor Sjælen (Elemente I p. 18).

\textsuperscript{126}) S. 542. Om Forholdet mellem Kultur og Lykke smlgn. min
\textit{Etik.} Kap. 7, og mine \textit{Etiske Undersøgelser.} Kap. 2.

\textsuperscript{127}) S. 548. Vel fremkom \emph{Karl} \emph{Rokitanskys}
aandfulde Afhandling \textit{Der selbständige Werth des Wissens}
(Sitzungsberichte der Wienerakademie 1867) først over et Aar efter Gesch. d.
Mat., men allerede 1862 havde han i en Tale henvist til Kants Idealisme som den
naturvidenskabelige Tænknings sande Konsekvens. Smlgn. Th. \emph{Meynert}:
\textit{Karl Rokitansky.} Ein Nachruf. (Sammlung von populär-wissenschaftlichen
Vorträgen. Wien und Leipzig 1892. p. 76 f.). --- Ogsaa \emph{Helmholtz} (i sin
\textit{Physiologische Optik.} 1867. p. 443 ff.) hævder, at det i sidste
Instans er uberettiget at tillægge vore Forestillinger anden Gyldighed end en
praktisk og symbolsk. Til et lignende Resultat kommer fremdeles Fysiologen A.
\emph{Fick} i sit Foredrag \textit{Die Welt als Vorstellung.} Würzburg 1870.

\textsuperscript{128}) S. 549. I sin interessante Afhandling \textit{Der
Idealismus Friedrich Albert Langes} (Vierteljahrschrift für wissensch. Philos.
I p. 183---185) har \emph{Max} \emph{Heinze} paavist dette.

\textsuperscript{129}) S. 549. Smlgn. min \textit{Psykologi.} V D, og
\textit{Filosofiske Problemer} (Universitetsprogram 1902). p. 48---52.

\textsuperscript{130}) S. 565. Om Dührings Stilling til det sociale Spørgsmaal
og om de forskellige Standpunkter, han her har indtaget, henviser jeg til B.
\emph{Friedländer}: \textit{Die vier Hauptrichtungen der modernen socialen
Bewegung.} Berlin 1901.
\end{small}

\end{document}
